\documentclass[12pt, a4paper]{article}

% Codigo y leng 
\usepackage[utf8]{inputenc}
\usepackage[T1]{fontenc}
\usepackage[spanish]{babel}

% Tipografia
\usepackage{lmodern}
\usepackage{microtype}

% Matematica
\usepackage{amsmath}
\usepackage{amssymb}
\usepackage{bm}

% Gráficos y Figuras 
\usepackage{graphicx}
\usepackage{float}
\usepackage{caption}
\usepackage{subcaption}
\usepackage{tikz}
\usetikzlibrary{shapes, arrows.meta, positioning, calc}

% Tablas
\usepackage{booktabs}
\usepackage{array}
\usepackage{tabularx}

% Código fuente
\usepackage{listings}
\usepackage{xcolor}

% Algoritmos
\usepackage{algorithm}
\usepackage{algpseudocode}

% Refs
\usepackage[hidelinks]{hyperref}
\usepackage[backend=biber, style=ieee, sorting=none]{biblatex}
\addbibresource{refs/references.bib}

% Layout
\usepackage{geometry}
\geometry{margin=2.5cm}
\usepackage{setspace}
\onehalfspacing
\usepackage{parskip}

% Listings
\definecolor{codebg}{RGB}{245,245,245}
\definecolor{codegreen}{RGB}{0,128,0}
\definecolor{codeblue}{RGB}{0,0,180}
\definecolor{codegray}{RGB}{128,128,128}
\lstdefinestyle{python}{
    backgroundcolor=\color{codebg},
    commentstyle=\color{codegreen},
    keywordstyle=\color{codeblue}\bfseries,
    numberstyle=\tiny\color{codegray},
    stringstyle=\color{orange},
    basicstyle=\ttfamily\footnotesize,
    breakatwhitespace=false,
    breaklines=true,
    captionpos=b,
    keepspaces=true,
    numbers=left,
    numbersep=5pt,
    showspaces=false,
    showstringspaces=false,
    showtabs=false,
    tabsize=4,
    frame=single,
    language=Python
}
\lstset{style=python}

\title{
    \vspace{-1cm}
    {\large Inteligencia Artificial y Sistemas Expertos} \\[0.5em]
    \textbf{ Transformers } \\
    Fundamentos Teóricos, Arquitectura y Prototipo Mínimo \\[0.5em]
    {\normalsize Del perceptrón al mecanismo de atención}
}
\author{Ulises Ramirez}
\date{\today}
\begin{document}

\maketitle
\thispagestyle{empty}

\begin{abstract}
Este trabajo analiza la arquitectura Transformer como evolución natural
de las redes neuronales clásicas, contrastándola con los modelos estudiados
en la cátedra: perceptrón, perceptrón multicapa, redes de Hopfield y mapas
autoorganizados. Se presenta el mecanismo de auto-atención como solución al
problema de dependencias contextuales en secuencias, y se desarrolla un
prototipo mínimo funcional de un modelo de lenguaje basado en Transformer,
implementado desde cero para evidenciar la correspondencia entre teoría e
implementación.
\end{abstract}

\tableofcontents
\newpage

\section{Introducción}
\label{sec:introduccion}

En el campo de la Inteligencia Artificial, las redes neuronales artificiales
han sido históricamente una de las principales herramientas para el modelado
de funciones complejas y el aprendizaje a partir de datos. En cursos
introductorios de la disciplina, se estudian arquitecturas clásicas como el
perceptrón, el perceptrón multicapa (MLP), las redes de Hopfield y los mapas
autoorganizados (SOM), las cuales permiten comprender los fundamentos del
aprendizaje supervisado, no supervisado y asociativo.

Sin embargo, estas arquitecturas presentan limitaciones estructurales cuando
se enfrentan a problemas que requieren el modelado de dependencias contextuales
complejas, como es el caso del procesamiento del lenguaje natural (NLP, por sus
siglas en inglés). El lenguaje humano es inherentemente secuencial, ambiguo y
dependiente del contexto, características que dificultan su tratamiento mediante
modelos diseñados para entradas de tamaño fijo o relaciones puramente locales.

En este contexto, la arquitectura Transformer, introducida por
\textcite{vaswani2017attention}, representa un cambio de paradigma en el diseño
de redes neuronales para el procesamiento de secuencias, al basarse en el
mecanismo de atención en lugar de recurrencia o convolución. Su impacto fue tal
que dio origen a la familia de modelos de lenguaje de gran escala que hoy
sustentan aplicaciones como traducción automática, generación de texto y
asistentes conversacionales~\cite{rothman2022transformers}.

Este trabajo tiene como objetivo analizar los principios teóricos que
fundamentan la arquitectura Transformer, contrastándolos con los modelos
clásicos estudiados en la cátedra, y desarrollar un prototipo mínimo funcional
que permita evidenciar dichas diferencias de manera concreta. El prototipo
consiste en un modelo de lenguaje autoregresivo a nivel carácter, construido
desde cero en PyTorch, sin recurrir a implementaciones de alto nivel de la
librería.

\subsection*{Alcance y nivel de tratamiento}

Las arquitecturas y conceptos presentados en este trabajo se describen con
un nivel de detalle acorde al carácter introductorio de la cátedra. En
particular, los modelos clásicos (perceptrón, MLP, Hopfield, SOM) se
retoman como punto de contraste conceptual, sin pretender agotar su
formalización matemática completa. Del mismo modo, la descripción del
Transformer y su implementación apuntan a la comprensión de los principios
fundamentales, y no a reproducir el estado del arte en escala o rendimiento.
El objetivo es siempre la claridad conceptual por sobre la exhaustividad
técnica.

\subsection*{Organización del documento}

El documento se organiza de la siguiente manera:
la Sección~\ref{sec:redes_clasicas} presenta las arquitecturas clásicas y sus
limitaciones estructurales;
la Sección~\ref{sec:problema_nlp} caracteriza el problema del procesamiento del
lenguaje natural;
la Sección~\ref{sec:mecanismo_atencion} introduce el mecanismo de atención como
respuesta a dichas limitaciones;
la Sección~\ref{sec:arquitectura_transformer} describe en detalle la arquitectura
Transformer;
la Sección~\ref{sec:transformer_minimo} define formalmente el modelo implementado;
la Sección~\ref{sec:conclusiones} presenta las conclusiones del trabajo; y el
Anexo~\ref{sec:implementacion} incluye la implementación comentada.

\section{Modelos clásicos de inteligencia artificial y sus limitaciones}
\label{sec:redes_clasicas}

Esta sección presenta las arquitecturas de redes neuronales abordadas en la
cátedra, analizando sus principios de funcionamiento y las limitaciones que
las hacen insuficientes para el procesamiento del lenguaje natural. Al final
se presenta una tabla comparativa que sintetiza el contraste con la arquitectura
Transformer.

%------------------------------------------------------------
\subsection{El perceptrón}
\label{subsec:perceptron}

El perceptrón es la unidad computacional más elemental de las redes
neuronales artificiales~\cite{goodfellow2016deeplearning}. Recibe un vector
de entradas $\bm{x} \in \mathbb{R}^n$, calcula una combinación lineal
ponderada con pesos $\bm{w}$ y un sesgo $b$, y aplica una función de
activación $\varphi$ para producir la salida $y$:

\begin{equation}
    y = \varphi\!\left(\sum_{i=1}^{n} w_i x_i + b\right)
\end{equation}

% ============================================================
% fig_perceptron.tex: Diagrama del perceptrón
% Uso: % ============================================================
% fig_perceptron.tex: Diagrama del perceptrón
% Uso: \input{figures/fig_perceptron}
% ============================================================
\begin{figure}[H]
\centering
\begin{tikzpicture}[
    node distance=1.4cm,
    input/.style={circle, draw=gray!60, fill=gray!15, minimum size=0.8cm, font=\small},
    neuron/.style={circle, draw=blue!60, fill=blue!10, minimum size=1.0cm, font=\small\bfseries},
    output/.style={circle, draw=green!60!black, fill=green!10, minimum size=0.8cm, font=\small},
    arrow/.style={-{Stealth}, thick}
]

% Entradas
\node[input] (x1) at (0, 2)   {$x_1$};
\node[input] (x2) at (0, 0.7) {$x_2$};
\node        (xd) at (0,-0.1) {$\vdots$};
\node[input] (xn) at (0,-1)   {$x_n$};

% Neurona
\node[neuron] (n) at (3, 0.7) {$\Sigma$};

% Salida
\node[output] (y) at (5.5, 0.7) {$y$};

% Bias
\node[input] (b) at (1.5, -1.8) {$b$};

% Función de activación (caja)
\node[draw=orange!70, fill=orange!10, rounded corners, minimum width=1.1cm,
      minimum height=0.7cm, font=\small] (phi) at (4.4, 0.7) {$\varphi$};

% Flechas entradas → neurona
\draw[arrow] (x1) -- node[above, font=\tiny] {$w_1$} (n);
\draw[arrow] (x2) -- node[above, font=\tiny] {$w_2$} (n);
\draw[arrow] (xn) -- node[below, font=\tiny] {$w_n$} (n);
\draw[arrow] (b)  -- (n);

% Neurona → activación → salida
\draw[arrow] (n)   -- (phi);
\draw[arrow] (phi) -- (y);

% Etiquetas
\node[font=\footnotesize, gray] at (0, -2.2) {Entradas};
\node[font=\footnotesize, gray] at (3, -2.2) {Suma ponderada};
\node[font=\footnotesize, gray] at (5.5, -2.2) {Salida};

\end{tikzpicture}
\caption{Estructura del perceptrón. Las entradas $x_i$ se combinan linealmente
con pesos $w_i$ y un sesgo $b$; la suma se pasa por una función de activación
$\varphi$ para producir la salida $y$.}
\label{fig:perceptron}
\end{figure}

% ============================================================
\begin{figure}[H]
\centering
\begin{tikzpicture}[
    node distance=1.4cm,
    input/.style={circle, draw=gray!60, fill=gray!15, minimum size=0.8cm, font=\small},
    neuron/.style={circle, draw=blue!60, fill=blue!10, minimum size=1.0cm, font=\small\bfseries},
    output/.style={circle, draw=green!60!black, fill=green!10, minimum size=0.8cm, font=\small},
    arrow/.style={-{Stealth}, thick}
]

% Entradas
\node[input] (x1) at (0, 2)   {$x_1$};
\node[input] (x2) at (0, 0.7) {$x_2$};
\node        (xd) at (0,-0.1) {$\vdots$};
\node[input] (xn) at (0,-1)   {$x_n$};

% Neurona
\node[neuron] (n) at (3, 0.7) {$\Sigma$};

% Salida
\node[output] (y) at (5.5, 0.7) {$y$};

% Bias
\node[input] (b) at (1.5, -1.8) {$b$};

% Función de activación (caja)
\node[draw=orange!70, fill=orange!10, rounded corners, minimum width=1.1cm,
      minimum height=0.7cm, font=\small] (phi) at (4.4, 0.7) {$\varphi$};

% Flechas entradas → neurona
\draw[arrow] (x1) -- node[above, font=\tiny] {$w_1$} (n);
\draw[arrow] (x2) -- node[above, font=\tiny] {$w_2$} (n);
\draw[arrow] (xn) -- node[below, font=\tiny] {$w_n$} (n);
\draw[arrow] (b)  -- (n);

% Neurona → activación → salida
\draw[arrow] (n)   -- (phi);
\draw[arrow] (phi) -- (y);

% Etiquetas
\node[font=\footnotesize, gray] at (0, -2.2) {Entradas};
\node[font=\footnotesize, gray] at (3, -2.2) {Suma ponderada};
\node[font=\footnotesize, gray] at (5.5, -2.2) {Salida};

\end{tikzpicture}
\caption{Estructura del perceptrón. Las entradas $x_i$ se combinan linealmente
con pesos $w_i$ y un sesgo $b$; la suma se pasa por una función de activación
$\varphi$ para producir la salida $y$.}
\label{fig:perceptron}
\end{figure}


El conocimiento del perceptrón reside íntegramente en los pesos $\bm{w}$,
que se ajustan durante el entrenamiento. Sin embargo, su capacidad
representacional es fundamentalmente limitada: solo puede producir
fronteras de decisión lineales en el espacio de entrada. Cualquier problema
con estructura no lineal —como la función XOR— no puede ser resuelto por
un único perceptrón. Esta limitación motivó el desarrollo del perceptrón
multicapa.

%------------------------------------------------------------
\subsection{El perceptrón multicapa, backpropagation y descenso del gradiente}
\label{subsec:mlp}

El perceptrón multicapa (MLP) supera las limitaciones del perceptrón simple
apilando capas de neuronas con funciones de activación no lineales. El teorema
de aproximación universal establece que un MLP con al menos una capa oculta
suficientemente amplia puede aproximar cualquier función continua en un
compacto~\cite{goodfellow2016deeplearning}.

\begin{figure}[H]
\centering
\begin{tikzpicture}[
    neuron/.style={circle, draw=blue!50, fill=blue!8, minimum size=0.75cm, font=\scriptsize},
    out/.style={circle, draw=green!60!black, fill=green!8, minimum size=0.75cm, font=\scriptsize},
    arrow/.style={-{Stealth}, thick},
    back/.style={-{Stealth}, thick, dashed, red!70},
    layer/.style={font=\small\bfseries, gray}
]

% Capa de entrada (3 nodos)
\foreach \i/\y in {1/2, 2/0.7, 3/-0.7}{
    \node[neuron] (i\i) at (0, \y) {$x_{\i}$};
}

% Capa oculta 1 (4 nodos)
\foreach \i/\y in {1/2.7, 2/1.2, 3/-0.2, 4/-1.5}{
    \node[neuron] (h1\i) at (2.5, \y) {};
}

% Capa oculta 2 (4 nodos)
\foreach \i/\y in {1/2.7, 2/1.2, 3/-0.2, 4/-1.5}{
    \node[neuron] (h2\i) at (5, \y) {};
}

% Capa de salida (2 nodos)
\foreach \i/\y in {1/1.0, 2/-0.3}{
    \node[out=45, in=135] (o\i) at (7.5, \y) {$y_{\i}$};
}

% Conexiones entrada → oculta1
\foreach \i in {1,2,3}{
    \foreach \j in {1,2,3,4}{
        \draw[-{Stealth}, gray!50, thin] (i\i) -- (h1\j);
    }
}
% Conexiones oculta1 → oculta2
\foreach \i in {1,2,3,4}{
    \foreach \j in {1,2,3,4}{
        \draw[-{Stealth}, gray!50, thin] (h1\i) -- (h2\j);
    }
}
% Conexiones oculta2 → salida
\foreach \i in {1,2,3,4}{
    \foreach \j in {1,2}{
        \draw[-{Stealth}, gray!50, thin] (h2\i) -- (o\j);
    }
}

% Flecha de forward pass (arriba)
\draw[arrow, blue!70] (-0.6, 3.5) -- node[above, font=\small] {Forward pass $\longrightarrow$} (8.1, 3.5);

% Flecha de backward pass (abajo)
\draw[back] (8.1, -2.3) -- node[below, font=\small] {$\longleftarrow$ Backpropagation $\nabla\mathcal{L}$} (-0.6, -2.3);

% Etiquetas de capas
\node[layer] at (0,   -3.3) {Entrada};
\node[layer] at (2.5, -3.3) {Oculta 1};
\node[layer] at (5,   -3.3) {Oculta 2};
\node[layer] at (7.5, -3.3) {Salida};

\end{tikzpicture}
\caption{Perceptrón multicapa (MLP). El \textit{forward pass} propaga la
información de izquierda a derecha. La retropropagación (\textit{backpropagation})
calcula los gradientes de la pérdida $\mathcal{L}$ en sentido inverso,
actualizando los pesos mediante descenso del gradiente.}
\label{fig:mlp}
\end{figure}


\subsubsection*{Backpropagation}

El entrenamiento del MLP se realiza mediante el algoritmo de retropropagación
del error (\textit{backpropagation}), que calcula eficientemente el gradiente
de la función de pérdida $\mathcal{L}$ respecto a cada parámetro de la red
aplicando la regla de la cadena en sentido inverso a través de las capas:

\begin{equation}
    \frac{\partial \mathcal{L}}{\partial w_{ij}^{(l)}} =
    \frac{\partial \mathcal{L}}{\partial a_j^{(l)}} \cdot
    \frac{\partial a_j^{(l)}}{\partial w_{ij}^{(l)}}
\end{equation}

\noindent donde $a_j^{(l)}$ es la activación de la neurona $j$ en la capa $l$
y $w_{ij}^{(l)}$ es el peso que la conecta con la neurona $i$ de la capa
anterior. El proceso recorre la red de salida hacia entrada (\textit{backward}),
acumulando gradientes en cada capa.

\subsubsection*{Descenso del gradiente}

Los gradientes obtenidos por backpropagation se usan para actualizar los
parámetros mediante descenso del gradiente. La regla de actualización es:

\begin{equation}
    w \leftarrow w - \eta \, \frac{\partial \mathcal{L}}{\partial w}
\end{equation}

\noindent donde $\eta > 0$ es la \textit{tasa de aprendizaje}. Intuitivamente,
cada actualización mueve los parámetros en la dirección que reduce la pérdida.
Si $\eta$ es demasiado grande el proceso puede divergir u oscilar; si es
demasiado pequeño, la convergencia es lenta o puede quedar atrapada en
mínimos locales subóptimos.

% ============================================================
% fig_gradient_descent.tex: Descenso del gradiente
% Uso: % ============================================================
% fig_gradient_descent.tex: Descenso del gradiente
% Uso: \input{figures/fig_gradient_descent}
% ============================================================
\begin{figure}[H]
\centering
\begin{tikzpicture}[scale=1.1]

% Ejes
\draw[-{Stealth}, thick] (-0.3, 0) -- (6.5, 0) node[right] {$w$};
\draw[-{Stealth}, thick] (0, -0.3) -- (0, 4.2) node[above] {$\mathcal{L}(w)$};

% Curva de pérdida (parábola suavizada)
\draw[blue!70, very thick, domain=0.3:6.1, samples=80]
    plot (\x, {0.35*(\x-3.2)^2 + 0.5});

% Mínimo
\fill[green!60!black] (3.2, 0.5) circle (3pt) node[below right, font=\small] {mínimo};

% Puntos del descenso
\coordinate (p1) at (1.0, 2.34);
\coordinate (p2) at (1.7, 1.44);
\coordinate (p3) at (2.3, 0.87);
\coordinate (p4) at (2.75, 0.59);
\coordinate (p5) at (3.05, 0.51);

\fill[red!70]    (p1) circle (3pt);
\fill[red!60]    (p2) circle (3pt);
\fill[orange!80] (p3) circle (3pt);
\fill[orange!60] (p4) circle (3pt);
\fill[yellow!80!orange] (p5) circle (3pt);

% Flechas de actualización
\draw[-{Stealth}, red!60, thick]    (p1) -- (p2);
\draw[-{Stealth}, red!40, thick]    (p2) -- (p3);
\draw[-{Stealth}, orange!80, thick] (p3) -- (p4);
\draw[-{Stealth}, orange!50, thick] (p4) -- (p5);

% Etiqueta tasa de aprendizaje
\draw[dashed, gray] (1.0, 0) -- (p1);
\draw[dashed, gray] (1.7, 0) -- (p2);
\node[font=\tiny, gray] at (1.35, -0.2) {$\eta \cdot \nabla\mathcal{L}$};
\draw[{Stealth}-{Stealth}, gray, thin] (1.05, -0.15) -- (1.65, -0.15);

% Fórmula
\node[draw=blue!30, fill=blue!5, rounded corners, font=\small,
      inner sep=6pt] at (5.0, 3.5)
      {$w \leftarrow w - \eta \dfrac{\partial \mathcal{L}}{\partial w}$};

\node[font=\footnotesize, gray] at (3.2, -0.5) {$w^*$};

\end{tikzpicture}
\caption{Descenso del gradiente. Partiendo de un valor inicial de $w$, cada
iteración actualiza el parámetro en dirección opuesta al gradiente de la
pérdida $\mathcal{L}$, con magnitud controlada por la tasa de aprendizaje
$\eta$. El proceso converge progresivamente al mínimo $w^*$.}
\label{fig:gradient_descent}
\end{figure}

% ============================================================
\begin{figure}[H]
\centering
\begin{tikzpicture}[scale=1.1]

% Ejes
\draw[-{Stealth}, thick] (-0.3, 0) -- (6.5, 0) node[right] {$w$};
\draw[-{Stealth}, thick] (0, -0.3) -- (0, 4.2) node[above] {$\mathcal{L}(w)$};

% Curva de pérdida (parábola suavizada)
\draw[blue!70, very thick, domain=0.3:6.1, samples=80]
    plot (\x, {0.35*(\x-3.2)^2 + 0.5});

% Mínimo
\fill[green!60!black] (3.2, 0.5) circle (3pt) node[below right, font=\small] {mínimo};

% Puntos del descenso
\coordinate (p1) at (1.0, 2.34);
\coordinate (p2) at (1.7, 1.44);
\coordinate (p3) at (2.3, 0.87);
\coordinate (p4) at (2.75, 0.59);
\coordinate (p5) at (3.05, 0.51);

\fill[red!70]    (p1) circle (3pt);
\fill[red!60]    (p2) circle (3pt);
\fill[orange!80] (p3) circle (3pt);
\fill[orange!60] (p4) circle (3pt);
\fill[yellow!80!orange] (p5) circle (3pt);

% Flechas de actualización
\draw[-{Stealth}, red!60, thick]    (p1) -- (p2);
\draw[-{Stealth}, red!40, thick]    (p2) -- (p3);
\draw[-{Stealth}, orange!80, thick] (p3) -- (p4);
\draw[-{Stealth}, orange!50, thick] (p4) -- (p5);

% Etiqueta tasa de aprendizaje
\draw[dashed, gray] (1.0, 0) -- (p1);
\draw[dashed, gray] (1.7, 0) -- (p2);
\node[font=\tiny, gray] at (1.35, -0.2) {$\eta \cdot \nabla\mathcal{L}$};
\draw[{Stealth}-{Stealth}, gray, thin] (1.05, -0.15) -- (1.65, -0.15);

% Fórmula
\node[draw=blue!30, fill=blue!5, rounded corners, font=\small,
      inner sep=6pt] at (5.0, 3.5)
      {$w \leftarrow w - \eta \dfrac{\partial \mathcal{L}}{\partial w}$};

\node[font=\footnotesize, gray] at (3.2, -0.5) {$w^*$};

\end{tikzpicture}
\caption{Descenso del gradiente. Partiendo de un valor inicial de $w$, cada
iteración actualiza el parámetro en dirección opuesta al gradiente de la
pérdida $\mathcal{L}$, con magnitud controlada por la tasa de aprendizaje
$\eta$. El proceso converge progresivamente al mínimo $w^*$.}
\label{fig:gradient_descent}
\end{figure}


\subsubsection*{Limitaciones para el modelado de secuencias}

A pesar de su capacidad expresiva, los MLPs presentan limitaciones
estructurales que los hacen inadecuados para procesar lenguaje natural:

\begin{itemize}
    \item \textbf{Entradas de tamaño fijo.} El MLP requiere que la dimensión
    de la entrada esté fija en el momento del diseño. Para modelar texto,
    esto implica definir una ventana de contexto de longitud constante,
    truncando o rellenando todas las secuencias a esa longitud. Esto hace
    imposible adaptarse a oraciones de longitud variable sin pérdida de
    información~\cite{bengio2003neural}.

    \item \textbf{Ausencia de codificación del orden.} La posición de cada
    elemento en el vector de entrada no tiene representación interna
    diferenciada: si se permutan los elementos, la red no puede distinguir
    el nuevo orden del original. Para el lenguaje esto es crítico: \textit{``el
    perro muerde al hombre''} y \textit{``el hombre muerde al perro''} son
    semánticamente opuestas, pero un MLP que reciba las mismas palabras en
    distinto orden producirá la misma activación interna si no hay codificación
    posicional explícita~\cite{goodfellow2016deeplearning}.

    \item \textbf{Sin mecanismo relacional dinámico.} En un MLP los pesos son
    estáticos: la red no puede decidir en tiempo de inferencia a qué parte de
    la entrada prestarle más atención según el contexto. Dos palabras pueden
    ser muy relevantes entre sí en una oración y completamente irrelevantes
    en otra, pero el modelo las procesa siempre con los mismos pesos.

    \item \textbf{Dependencias limitadas al tamaño de la ventana.} Toda
    relación que exija considerar elementos separados por más posiciones que
    el ancho de la ventana es inaccesible para el modelo. En textos largos,
    esto descarta dependencias fundamentales: el sujeto de una oración puede
    estar decenas de palabras antes del verbo que lo acompaña~\cite{bengio2003neural}.
\end{itemize}

%------------------------------------------------------------
\subsection{Redes de Hopfield}
\label{subsec:hopfield}

Las redes de Hopfield~\cite{hopfield1982neural} son redes recurrentes
completamente conectadas diseñadas como modelos de \textit{memoria asociativa}:
dado un patrón de entrada incompleto o ruidoso, la red recupera el patrón
almacenado más cercano. Cada unidad $s_i \in \{-1, +1\}$ se conecta
simétricamente con todas las demás ($w_{ij} = w_{ji}$), y la dinámica del
sistema minimiza una función de energía global:

\begin{equation}
    E = -\frac{1}{2} \sum_{i \neq j} w_{ij} \, s_i \, s_j
\end{equation}

La regla de actualización de cada neurona es:

\begin{equation}
    s_i \leftarrow \text{sgn}\!\left(\sum_{j} w_{ij} \, s_j\right)
\end{equation}

La red converge cuando ninguna actualización reduce la energía, es decir,
cuando alcanza un mínimo local. Estos mínimos corresponden a los patrones
almacenados durante el entrenamiento. La Figura~\ref{fig:hopfield_energia}
muestra este paisaje de energía conceptualmente.

% ============================================================
% fig_hopfield_energia.tex — Paisaje de energía de Hopfield
% ============================================================
\begin{figure}[H]
\centering
\begin{tikzpicture}[scale=1.1]

% Ejes
\draw[-{Stealth}, thick] (-0.3, 0) -- (7.2, 0) node[right, font=\small] {espacio de estados};
\draw[-{Stealth}, thick] (0, -0.3) -- (0, 4.0) node[above, font=\small] {$E$};

% Curva de energía con varios mínimos locales
% Fix: usar grados directamente (sin sufijo r) con factores escalados
\draw[blue!70, very thick, domain=0.2:7.0, samples=200]
    plot (\x, {
        0.6*sin(51.43*\x) +
        0.4*sin(102.86*\x) +
        0.3*cos(77.14*\x) +
        1.8
    });

% Mínimos (patrones almacenados)
\fill[green!60!black] (1.05, 0.82) circle (3.5pt);
\fill[green!60!black] (3.55, 0.62) circle (3.5pt);
\fill[green!60!black] (5.90, 0.95) circle (3.5pt);

\node[green!50!black, font=\footnotesize, above right] at (1.05, 0.82) {patrón $\xi^1$};
\node[green!50!black, font=\footnotesize, above right] at (3.55, 0.62) {patrón $\xi^2$};
\node[green!50!black, font=\footnotesize, above right] at (5.90, 0.95) {patrón $\xi^3$};

% Estado inicial ruidoso → convergencia hacia mínimo
\fill[red!70]    (2.3, 2.85) circle (3pt) node[above, font=\footnotesize, red!70] {entrada ruidosa};
\draw[-{Stealth}, red!60, thick, dashed]
    (2.3, 2.85) .. controls (2.1, 2.0) and (1.7, 1.4) .. (1.05, 0.85);
\node[red!60, font=\tiny] at (1.5, 1.9) {convergencia};

% Etiqueta eje Y
\node[font=\footnotesize, gray, rotate=90] at (-0.55, 2.0) {Energía $E$};

\end{tikzpicture}
\caption{Paisaje de energía de una red de Hopfield. Cada mínimo local corresponde
a un patrón almacenado $\xi^k$. Ante una entrada ruidosa o incompleta, la dinámica
de la red desciende por el gradiente de energía hasta converger al mínimo más
cercano, recuperando el patrón original.}
\label{fig:hopfield_energia}
\end{figure}


\subsubsection*{Ejemplo numérico: recuperación de un patrón}

Para ilustrar el mecanismo, considérese una red de 4 neuronas con un patrón
almacenado $\xi^1 = [+1,-1,+1,-1]$. Los pesos se calculan mediante la regla
de Hebb:
$w_{ij} = \frac{1}{n} \sum_\mu \xi_i^\mu \xi_j^\mu$.
Si se presenta la entrada ruidosa $\bm{s} = [+1,+1,+1,-1]$ (bit $s_2$
invertido), la regla de actualización corrige el error en una sola iteración,
recuperando el patrón original.

% ============================================================
% fig_hopfield_ejemplo.tex: Ejemplo numérico paso a paso (4 neuronas)
% ============================================================
\begin{figure}[H]
\centering
\begin{tikzpicture}[
    cell/.style={draw, minimum width=0.55cm, minimum height=0.55cm, font=\footnotesize},
    on/.style={cell, fill=blue!60, text=white},
    off/.style={cell, fill=gray!20},
    lbl/.style={font=\footnotesize\bfseries},
    stepnum/.style={font=\small\bfseries, text=blue!70},
    arrow/.style={-{Stealth}, thick, gray}
]

% ── Fila 0: encabezados ─────────────────────────────────────
\node[lbl] at (-1.8, 3.8) {Paso};
\node[lbl] at ( 0.4, 3.8) {Estado $\bm{s}$};
\node[lbl] at ( 3.6, 3.8) {Descripción};

% ── Paso 0: patrón almacenado ───────────────────────────────
\node[stepnum]  at (-1.8, 3.0) {0};
\foreach \v/\i in {+1/0, -1/1, +1/2, -1/3}{
    \pgfmathparse{(\v>0)?1:0}
    \ifnum\pgfmathresult=1
        \node[on]  at (\i*0.65, 3.0) {\v};
    \else
        \node[off] at (\i*0.65, 3.0) {\v};
    \fi
}
\node[font=\footnotesize, align=left] at (5.9, 3.0)
    {Patrón almacenado $\xi^1 = [+1,-1,+1,-1]$};

% ── Separador ───────────────────────────────────────────────
\draw[dashed, gray!50] (-2.5, 2.55) -- (5.0, 2.55);

% ── Paso 1: entrada ruidosa ─────────────────────────────────
\node[stepnum]  at (-1.8, 2.1) {1};
\foreach \v/\i in {+1/0, +1/1, +1/2, -1/3}{
    \pgfmathparse{(\v>0)?1:0}
    \ifnum\pgfmathresult=1
        \node[on]  at (\i*0.65, 2.1) {\v};
    \else
        \node[off] at (\i*0.65, 2.1) {\v};
    \fi
}
\draw[red!70, very thick] (0.65-0.27, 2.1-0.27) rectangle (0.65+0.27, 2.1+0.27);
\node[font=\footnotesize, align=left] at (5.6, 2.1)
    {Entrada ruidosa: $s_2$ invertido {\color{red!70}$\leftarrow$ error}};

% ── Flecha de actualización ─────────────────────────────────
\draw[arrow] (0.95, 1.78) -- (0.95, 1.35);
\node[font=\tiny, gray] at (2.5, 1.57) {actualización: $s_i \leftarrow \text{sgn}\!\left(\sum_j w_{ij} s_j\right)$};

% ── Paso 2: después de una iteración ────────────────────────
\node[stepnum]  at (-1.8, 1.05) {2};
\foreach \v/\i in {+1/0, -1/1, +1/2, -1/3}{
    \pgfmathparse{(\v>0)?1:0}
    \ifnum\pgfmathresult=1
        \node[on]  at (\i*0.65, 1.05) {\v};
    \else
        \node[off] at (\i*0.65, 1.05) {\v};
    \fi
}
\node[font=\footnotesize, align=left] at (6.0, 1.05)
    {Tras actualizar $s_2$: patrón recuperado $\xi^1$ {\color{green!60!black}$\checkmark$}};

% ── Estado estable ───────────────────────────────────────────
\draw[dashed, gray!50] (-2.5, 0.6) -- (6.8, 0.6);
\node[font=\footnotesize, gray, align=left] at (2.0, 0.25)
    {Estado estable: $E$ alcanzó mínimo local $\Rightarrow$ la red converge.};

% ── Leyenda ─────────────────────────────────────────────────
\node[on,  label=right:{\footnotesize $s_i = +1$ (activa)}]  at (-1.2, -0.35) {};
\node[off, label=right:{\footnotesize $s_i = -1$ (inactiva)}] at ( 2.8, -0.35) {};

\end{tikzpicture}
\caption{Ejemplo de recuperación de un patrón en una red de Hopfield con
4 neuronas. El patrón almacenado $\xi^1 = [+1,-1,+1,-1]$ se recupera
correctamente a partir de una versión ruidosa en la que el bit $s_2$ estaba
invertido. La regla de actualización $s_i \leftarrow \text{sgn}\!\left(\sum_j w_{ij} s_j\right)$
corrige el error en una sola iteración.}
\label{fig:hopfield_ejemplo}
\end{figure}


\subsubsection*{Conexión con los Transformers}

La red de Hopfield introduce un concepto que reaparecerá de forma central en
los Transformers: la \textit{influencia global entre unidades}. El estado de
cada neurona está determinado por el estado de todas las demás, lo que
constituye una forma primitiva de atención distribuida. Sin embargo, la
capacidad de almacenamiento clásica es limitada (aproximadamente $0.14n$
patrones) y el modelo no está diseñado para representar relaciones
secuenciales. Trabajos recientes han demostrado que el mecanismo de atención
puede interpretarse como una generalización de Hopfield con capacidad de
almacenamiento exponencial~\cite{ramsauer2020hopfield}.

%------------------------------------------------------------
\subsection{Mapas autoorganizados (SOM)}
\label{subsec:som}

Los mapas autoorganizados (SOM), propuestos por Kohonen~\cite{kohonen1990som},
son redes no supervisadas que aprenden a proyectar datos de alta dimensión en
una grilla discreta de baja dimensión (usualmente 2D), preservando la topología
del espacio original. Cada neurona de la grilla tiene asociado un vector de
pesos $\bm{w}_i \in \mathbb{R}^d$ del mismo tamaño que los datos de entrada.

\subsubsection*{Proceso de entrenamiento}

En cada iteración se presenta una muestra $\bm{x}$ y se identifican:

\begin{enumerate}
    \item \textbf{Neurona ganadora (BMU):} aquella cuyo vector de pesos es más
    cercano a $\bm{x}$ en distancia euclidiana:
    $c = \arg\min_i \|\bm{x} - \bm{w}_i\|$.

    \item \textbf{Actualización de la BMU y sus vecinas:} los pesos de la BMU
    y las neuronas en su entorno se acercan a $\bm{x}$:
    \begin{equation}
        \bm{w}_i(t{+}1) = \bm{w}_i(t) + \alpha(t)\,h_{ci}(t)\,
        [\bm{x}(t) - \bm{w}_i(t)]
    \end{equation}
    donde $\alpha(t)$ es la tasa de aprendizaje y $h_{ci}(t)$ es la función
    de vecindad (usualmente gaussiana), que decrece con la distancia en la
    grilla entre la BMU $c$ y la neurona $i$.
\end{enumerate}

Tras muchas iteraciones, neuronas contiguas en la grilla responden a entradas
similares: la topología del espacio de entrada queda reflejada en la grilla 2D.

% ============================================================
% fig_som_proceso.tex: Proceso de entrenamiento SOM (pasos)
% Uso: % ============================================================
% fig_som_proceso.tex: Proceso de entrenamiento SOM (pasos)
% Uso: \input{figures/fig_som_proceso}
% ============================================================
\begin{figure}[H]
\centering
\begin{tikzpicture}[
    neuron/.style={circle, draw=blue!50, fill=blue!8, minimum size=0.55cm},
    bmu/.style={circle, draw=red!70, fill=red!20, minimum size=0.55cm, very thick},
    neighbor/.style={circle, draw=orange!70, fill=orange!10, minimum size=0.55cm, thick},
    input/.style={circle, draw=green!60!black, fill=green!15, minimum size=0.6cm, font=\tiny},
    arrow/.style={-{Stealth}, thick},
    label/.style={font=\small\bfseries, gray}
]

% ─── Grilla 3×3 en t=0 (aleatoria) ──────────────────────────
\node[label] at (1.0, 4.1) {$t = 0$ (inicial)};
\foreach \row in {0,1,2}{
    \foreach \col in {0,1,2}{
        \node[neuron] at (\col*0.9, 3.2-\row*0.9) {};
    }
}
\node[font=\tiny, gray] at (1.0, 0.8) {pesos aleatorios};

% ─── Flecha paso 1 ──────────────────────────────────────────
\draw[arrow, blue!60] (2.4, 2.2) -- (3.2, 2.2);
\node[font=\tiny, align=center, blue!70] at (2.8, 2.55)
    {1. presentar\\entrada $\bm{x}$};

% ─── Grilla con BMU marcada ──────────────────────────────────
\node[label] at (5.0, 4.1) {BMU encontrada};
% Entrada
\node[input] at (5.0, 4.55) {$\bm{x}$};
\foreach \row in {0,1,2}{
    \foreach \col in {0,1,2}{
        \pgfmathparse{int(\row==1 && \col==1)}
        \ifnum\pgfmathresult=1
            \node[bmu]      at (3.6+\col*0.9, 3.2-\row*0.9) {};
        \else
            \pgfmathparse{int(abs(\row-1)<=1 && abs(\col-1)<=1)}
            \ifnum\pgfmathresult=1
                \node[neighbor] at (3.6+\col*0.9, 3.2-\row*0.9) {};
            \else
                \node[neuron]   at (3.6+\col*0.9, 3.2-\row*0.9) {};
            \fi
        \fi
    }
}
\draw[red!60, dashed] (5.0, 4.4) -- (4.5, 2.3);
\node[font=\tiny, gray] at (5.0, 0.8) {BMU=rojo, vecinos=naranja};

% ─── Flecha paso 2 ──────────────────────────────────────────
\draw[arrow, blue!60] (6.4, 2.2) -- (7.2, 2.2);
\node[font=\tiny, align=center, blue!70] at (6.8, 2.55)
    {2. actualizar\\BMU+vecinos};

% ─── Grilla actualizada ─────────────────────────────────────
\node[label] at (9.0, 4.1) {pesos actualizados};
\foreach \row in {0,1,2}{
    \foreach \col in {0,1,2}{
        \pgfmathparse{int(\row==1 && \col==1)}
        \ifnum\pgfmathresult=1
            \node[bmu]      at (7.6+\col*0.9, 3.2-\row*0.9) {};
        \else
            \pgfmathparse{int(abs(\row-1)<=1 && abs(\col-1)<=1)}
            \ifnum\pgfmathresult=1
                \node[neighbor] at (7.6+\col*0.9, 3.2-\row*0.9) {};
            \else
                \node[neuron]   at (7.6+\col*0.9, 3.2-\row*0.9) {};
            \fi
        \fi
    }
}
\node[font=\tiny, gray, align=center] at (9.0, 0.8) {BMU y vecinos se\\acercan a $\bm{x}$};

% ─── Leyenda regla de actualización ─────────────────────────
\node[draw=blue!20, fill=blue!5, rounded corners, font=\footnotesize,
      inner sep=5pt, align=left] at (5.5, -0.5)
      {$\bm{w}_i(t{+}1) = \bm{w}_i(t) + \alpha(t)\, h_{ci}(t)\,[\bm{x}(t) - \bm{w}_i(t)]$};

\end{tikzpicture}
\caption{Proceso de entrenamiento de un SOM en una iteración. (1) Se presenta
una entrada $\bm{x}$ y se identifica la neurona ganadora o \textit{Best Matching
Unit} (BMU, en rojo), la de menor distancia a $\bm{x}$. (2) La BMU y sus
vecinas (naranja) actualizan sus pesos acercándose a $\bm{x}$ según la regla
mostrada, donde $\alpha(t)$ es la tasa de aprendizaje y $h_{ci}(t)$ es la
función de vecindad.}
\label{fig:som_proceso}
\end{figure}

% ============================================================
\begin{figure}[H]
\centering
\begin{tikzpicture}[
    neuron/.style={circle, draw=blue!50, fill=blue!8, minimum size=0.55cm},
    bmu/.style={circle, draw=red!70, fill=red!20, minimum size=0.55cm, very thick},
    neighbor/.style={circle, draw=orange!70, fill=orange!10, minimum size=0.55cm, thick},
    input/.style={circle, draw=green!60!black, fill=green!15, minimum size=0.6cm, font=\tiny},
    arrow/.style={-{Stealth}, thick},
    label/.style={font=\small\bfseries, gray}
]

% ─── Grilla 3×3 en t=0 (aleatoria) ──────────────────────────
\node[label] at (1.0, 4.1) {$t = 0$ (inicial)};
\foreach \row in {0,1,2}{
    \foreach \col in {0,1,2}{
        \node[neuron] at (\col*0.9, 3.2-\row*0.9) {};
    }
}
\node[font=\tiny, gray] at (1.0, 0.8) {pesos aleatorios};

% ─── Flecha paso 1 ──────────────────────────────────────────
\draw[arrow, blue!60] (2.4, 2.2) -- (3.2, 2.2);
\node[font=\tiny, align=center, blue!70] at (2.8, 2.55)
    {1. presentar\\entrada $\bm{x}$};

% ─── Grilla con BMU marcada ──────────────────────────────────
\node[label] at (5.0, 4.1) {BMU encontrada};
% Entrada
\node[input] at (5.0, 4.55) {$\bm{x}$};
\foreach \row in {0,1,2}{
    \foreach \col in {0,1,2}{
        \pgfmathparse{int(\row==1 && \col==1)}
        \ifnum\pgfmathresult=1
            \node[bmu]      at (3.6+\col*0.9, 3.2-\row*0.9) {};
        \else
            \pgfmathparse{int(abs(\row-1)<=1 && abs(\col-1)<=1)}
            \ifnum\pgfmathresult=1
                \node[neighbor] at (3.6+\col*0.9, 3.2-\row*0.9) {};
            \else
                \node[neuron]   at (3.6+\col*0.9, 3.2-\row*0.9) {};
            \fi
        \fi
    }
}
\draw[red!60, dashed] (5.0, 4.4) -- (4.5, 2.3);
\node[font=\tiny, gray] at (5.0, 0.8) {BMU=rojo, vecinos=naranja};

% ─── Flecha paso 2 ──────────────────────────────────────────
\draw[arrow, blue!60] (6.4, 2.2) -- (7.2, 2.2);
\node[font=\tiny, align=center, blue!70] at (6.8, 2.55)
    {2. actualizar\\BMU+vecinos};

% ─── Grilla actualizada ─────────────────────────────────────
\node[label] at (9.0, 4.1) {pesos actualizados};
\foreach \row in {0,1,2}{
    \foreach \col in {0,1,2}{
        \pgfmathparse{int(\row==1 && \col==1)}
        \ifnum\pgfmathresult=1
            \node[bmu]      at (7.6+\col*0.9, 3.2-\row*0.9) {};
        \else
            \pgfmathparse{int(abs(\row-1)<=1 && abs(\col-1)<=1)}
            \ifnum\pgfmathresult=1
                \node[neighbor] at (7.6+\col*0.9, 3.2-\row*0.9) {};
            \else
                \node[neuron]   at (7.6+\col*0.9, 3.2-\row*0.9) {};
            \fi
        \fi
    }
}
\node[font=\tiny, gray, align=center] at (9.0, 0.8) {BMU y vecinos se\\acercan a $\bm{x}$};

% ─── Leyenda regla de actualización ─────────────────────────
\node[draw=blue!20, fill=blue!5, rounded corners, font=\footnotesize,
      inner sep=5pt, align=left] at (5.5, -0.5)
      {$\bm{w}_i(t{+}1) = \bm{w}_i(t) + \alpha(t)\, h_{ci}(t)\,[\bm{x}(t) - \bm{w}_i(t)]$};

\end{tikzpicture}
\caption{Proceso de entrenamiento de un SOM en una iteración. (1) Se presenta
una entrada $\bm{x}$ y se identifica la neurona ganadora o \textit{Best Matching
Unit} (BMU, en rojo), la de menor distancia a $\bm{x}$. (2) La BMU y sus
vecinas (naranja) actualizan sus pesos acercándose a $\bm{x}$ según la regla
mostrada, donde $\alpha(t)$ es la tasa de aprendizaje y $h_{ci}(t)$ es la
función de vecindad.}
\label{fig:som_proceso}
\end{figure}


% ============================================================
% fig_som_resultado.tex — Resultado: mapa semántico organizado
% Uso: % ============================================================
% fig_som_resultado.tex — Resultado: mapa semántico organizado
% Uso: \input{figures/fig_som_resultado}
% ============================================================
\begin{figure}[H]
\centering
\begin{tikzpicture}[scale=0.9]

% ─── Antes: datos sin organizar ──────────────────────────────
\node[font=\small\bfseries, gray] at (2.0, 5.2) {Datos de entrada};
\node[font=\tiny, gray]           at (2.0, 4.85) {(alta dimensión, sin orden)};

\foreach \x/\y/\col/\lbl in {
    0.4/4.2/red!60/animal,
    1.1/3.5/red!60/animal,
    0.6/3.1/red!60/animal,
    1.8/4.3/blue!60/vehículo,
    2.5/3.8/blue!60/vehículo,
    2.1/3.0/blue!60/vehículo,
    3.2/4.1/green!60!black/fruta,
    3.7/3.4/green!60!black/fruta,
    3.2/2.8/green!60!black/fruta
}{
    \fill[\col] (\x,\y) circle (3.5pt);
}
\node[font=\tiny, red!70]         at (0.7,  2.65) {animales};
\node[font=\tiny, blue!70]        at (2.3,  2.65) {vehículos};
\node[font=\tiny, green!60!black] at (3.5,  2.65) {frutas};

% ─── Flecha central ──────────────────────────────────────────
\draw[-{Stealth}, very thick, blue!60] (4.4, 3.7) -- (5.4, 3.7);
\node[font=\small, blue!70, align=center] at (4.9, 4.1) {entrenamiento\\SOM};

% ─── Después: mapa 4×4 organizado ───────────────────────────
\node[font=\small\bfseries, gray] at (8.5, 5.2) {Mapa SOM resultante};
\node[font=\tiny, gray]           at (8.5, 4.85) {(2D, topología preservada)};

% Celdas del mapa coloreadas por categoría emergente
\foreach \row/\col/\fillcol in {
    0/0/red!40,   0/1/red!30,   0/2/blue!20,  0/3/blue!40,
    1/0/red!35,   1/1/red!20,   1/2/blue!30,  1/3/blue!35,
    2/0/orange!20,2/1/orange!15,2/2/green!25, 2/3/green!40,
    3/0/orange!30,3/1/orange!25,3/2/green!35, 3/3/green!45
}{
    \fill[\fillcol] (5.8+\col*0.95, 4.4-\row*0.95)
        rectangle (5.8+\col*0.95+0.9, 4.4-\row*0.95-0.9);
    \draw[gray!40] (5.8+\col*0.95, 4.4-\row*0.95)
        rectangle (5.8+\col*0.95+0.9, 4.4-\row*0.95-0.9);
}

% Etiquetas de regiones
\node[font=\footnotesize\bfseries, red!70]
    at (6.75, 3.95) {animales};
\node[font=\footnotesize\bfseries, blue!70]
    at (8.15, 3.95) {vehículos};
\node[font=\footnotesize\bfseries, orange!80]
    at (6.75, 2.1) {mixto};
\node[font=\footnotesize\bfseries, green!60!black]
    at (8.15, 2.1) {frutas};

% Nota de gradiente suave
\draw[dashed, gray!50] (7.25, 4.4) -- (7.25, 0.55);
\draw[dashed, gray!50] (5.8, 2.55) -- (9.6, 2.55);
\node[font=\tiny, gray] at (8.5, 0.35)
    {Regiones similares son vecinas en el mapa (topología preservada)};

\end{tikzpicture}
\caption{Resultado del entrenamiento de un SOM. Los datos de entrada de alta
dimensión (izquierda) se proyectan en una grilla 2D (derecha) donde entradas
similares quedan agrupadas en regiones contiguas. La topología del espacio
original se preserva: las regiones de categorías cercanas semánticamente
(como \textit{animales} y \textit{mixto}) quedan adyacentes en el mapa~\cite{kohonen1990som}.}
\label{fig:som_resultado}
\end{figure}

% ============================================================
\begin{figure}[H]
\centering
\begin{tikzpicture}[scale=0.9]

% ─── Antes: datos sin organizar ──────────────────────────────
\node[font=\small\bfseries, gray] at (2.0, 5.2) {Datos de entrada};
\node[font=\tiny, gray]           at (2.0, 4.85) {(alta dimensión, sin orden)};

\foreach \x/\y/\col/\lbl in {
    0.4/4.2/red!60/animal,
    1.1/3.5/red!60/animal,
    0.6/3.1/red!60/animal,
    1.8/4.3/blue!60/vehículo,
    2.5/3.8/blue!60/vehículo,
    2.1/3.0/blue!60/vehículo,
    3.2/4.1/green!60!black/fruta,
    3.7/3.4/green!60!black/fruta,
    3.2/2.8/green!60!black/fruta
}{
    \fill[\col] (\x,\y) circle (3.5pt);
}
\node[font=\tiny, red!70]         at (0.7,  2.65) {animales};
\node[font=\tiny, blue!70]        at (2.3,  2.65) {vehículos};
\node[font=\tiny, green!60!black] at (3.5,  2.65) {frutas};

% ─── Flecha central ──────────────────────────────────────────
\draw[-{Stealth}, very thick, blue!60] (4.4, 3.7) -- (5.4, 3.7);
\node[font=\small, blue!70, align=center] at (4.9, 4.1) {entrenamiento\\SOM};

% ─── Después: mapa 4×4 organizado ───────────────────────────
\node[font=\small\bfseries, gray] at (8.5, 5.2) {Mapa SOM resultante};
\node[font=\tiny, gray]           at (8.5, 4.85) {(2D, topología preservada)};

% Celdas del mapa coloreadas por categoría emergente
\foreach \row/\col/\fillcol in {
    0/0/red!40,   0/1/red!30,   0/2/blue!20,  0/3/blue!40,
    1/0/red!35,   1/1/red!20,   1/2/blue!30,  1/3/blue!35,
    2/0/orange!20,2/1/orange!15,2/2/green!25, 2/3/green!40,
    3/0/orange!30,3/1/orange!25,3/2/green!35, 3/3/green!45
}{
    \fill[\fillcol] (5.8+\col*0.95, 4.4-\row*0.95)
        rectangle (5.8+\col*0.95+0.9, 4.4-\row*0.95-0.9);
    \draw[gray!40] (5.8+\col*0.95, 4.4-\row*0.95)
        rectangle (5.8+\col*0.95+0.9, 4.4-\row*0.95-0.9);
}

% Etiquetas de regiones
\node[font=\footnotesize\bfseries, red!70]
    at (6.75, 3.95) {animales};
\node[font=\footnotesize\bfseries, blue!70]
    at (8.15, 3.95) {vehículos};
\node[font=\footnotesize\bfseries, orange!80]
    at (6.75, 2.1) {mixto};
\node[font=\footnotesize\bfseries, green!60!black]
    at (8.15, 2.1) {frutas};

% Nota de gradiente suave
\draw[dashed, gray!50] (7.25, 4.4) -- (7.25, 0.55);
\draw[dashed, gray!50] (5.8, 2.55) -- (9.6, 2.55);
\node[font=\tiny, gray] at (8.5, 0.35)
    {Regiones similares son vecinas en el mapa (topología preservada)};

\end{tikzpicture}
\caption{Resultado del entrenamiento de un SOM. Los datos de entrada de alta
dimensión (izquierda) se proyectan en una grilla 2D (derecha) donde entradas
similares quedan agrupadas en regiones contiguas. La topología del espacio
original se preserva: las regiones de categorías cercanas semánticamente
(como \textit{animales} y \textit{mixto}) quedan adyacentes en el mapa~\cite{kohonen1990som}.}
\label{fig:som_resultado}
\end{figure}


\subsubsection*{Aporte conceptual: representación distribuida}

El aporte fundamental de los SOM es la \textit{representación distribuida y
topológica}: entidades similares activan regiones contiguas del mapa,
estableciendo un correlato geométrico de la similitud semántica. Esta idea
es precursora directa del concepto de \textit{embedding} que describe la
sección siguiente.

%------------------------------------------------------------
\subsection{Representaciones distribuidas y embeddings}
\label{subsec:embeddings}

La idea de representar entidades discretas como vectores continuos,
introducida conceptualmente por los SOM, se formalizó en el procesamiento
del lenguaje bajo el nombre de \textit{embedding}~\cite{mikolov2013word2vec}.
Comprender esta representación es esencial, ya que constituye el primer
bloque de procesamiento de todo Transformer.

\subsubsection*{Representación one-hot}

La alternativa más simple es la codificación \textit{one-hot}: cada token
recibe un vector de dimensión $|\mathcal{V}|$ (tamaño del vocabulario) con
un único componente igual a 1 en la posición correspondiente y 0 en el resto.
Esta representación es dispersa, no captura ninguna relación semántica entre
tokens, y la distancia euclídea entre cualquier par de tokens es idéntica
($\sqrt{2}$), independientemente de su similitud de significado.

\subsubsection*{Proyección mediante una matriz de embeddings}

Un \textit{embedding} proyecta cada token a un vector denso de dimensión
$d_{\text{model}} \ll |\mathcal{V}|$ mediante una matriz de embeddings
$E \in \mathbb{R}^{d_{\text{model}} \times |\mathcal{V}|}$, aprendida durante
el entrenamiento:

\begin{equation}
    \bm{e}_i = E\,\bm{v}_i \quad \in \mathbb{R}^{d_{\text{model}}}
\end{equation}

\noindent donde $\bm{v}_i$ es el vector \textit{one-hot} del token $i$.
Operacionalmente, multiplicar $E$ por un vector \textit{one-hot} equivale a
seleccionar la columna correspondiente de $E$: el embedding no es más que una
\textit{tabla de consulta} en la que cada token tiene asignado un vector
aprendido. En este espacio, tokens semánticamente similares quedan próximos
entre sí, y ciertas relaciones semánticas se expresan como operaciones
vectoriales regulares.

% ============================================================
% fig_proyeccion_embedding.tex — Proyección: one-hot → embedding
% Incluye tabla numérica + figura geométrica
% Uso: % ============================================================
% fig_proyeccion_embedding.tex — Proyección: one-hot → embedding
% Incluye tabla numérica + figura geométrica
% Uso: \input{figures/fig_proyeccion_embedding}
% ============================================================
\begin{figure}[H]
\centering

%──────────────────────────────────────────────────────────────
% PARTE SUPERIOR: tabla numérica de la proyección
%──────────────────────────────────────────────────────────────
\begin{tikzpicture}[
    cell/.style={draw=gray!40, minimum width=0.72cm, minimum height=0.60cm,
                 font=\scriptsize, inner sep=1pt},
    hdr/.style={font=\footnotesize\bfseries},
    hot/.style={fill=blue!50, text=white},
    emb/.style={fill=orange!20}
]

% ── Encabezado columnas ───────────────────────────────────
\node[hdr] at (-1.6, 2.8) {Token};
\node[hdr] at ( 1.4, 2.8) {One-hot $\bm{v}$ $\,(|\mathcal{V}|=4)$};
\node[hdr] at ( 5.6, 2.8) {Embedding $\bm{e} = E\bm{v}$ $\,(d_{\text{model}}=2)$};

% ── Fila: "gato" ─────────────────────────────────────────
\node[font=\small\itshape] at (-1.6, 2.1) {gato};
\foreach \v/\i in {1/0, 0/1, 0/2, 0/3}{
    \pgfmathparse{int(\v==1)}
    \ifnum\pgfmathresult=1
        \node[cell, hot] at (-0.3+\i*0.75, 2.1) {\v};
    \else
        \node[cell]      at (-0.3+\i*0.75, 2.1) {\v};
    \fi
}
\node[font=\Large] at (3.55, 2.1) {$\xrightarrow{\;E\;}$};
\node[cell, emb] at (4.6,  2.1) {0.8};
\node[cell, emb] at (5.35, 2.1) {0.3};

% ── Fila: "perro" ────────────────────────────────────────
\node[font=\small\itshape] at (-1.6, 1.3) {perro};
\foreach \v/\i in {0/0, 1/1, 0/2, 0/3}{
    \pgfmathparse{int(\v==1)}
    \ifnum\pgfmathresult=1
        \node[cell, hot] at (-0.3+\i*0.75, 1.3) {\v};
    \else
        \node[cell]      at (-0.3+\i*0.75, 1.3) {\v};
    \fi
}
\node[font=\Large] at (3.55, 1.3) {$\xrightarrow{\;E\;}$};
\node[cell, emb] at (4.6,  1.3) {0.7};
\node[cell, emb] at (5.35, 1.3) {0.4};

% ── Fila: "auto" ─────────────────────────────────────────
\node[font=\small\itshape] at (-1.6, 0.5) {auto};
\foreach \v/\i in {0/0, 0/1, 1/2, 0/3}{
    \pgfmathparse{int(\v==1)}
    \ifnum\pgfmathresult=1
        \node[cell, hot] at (-0.3+\i*0.75, 0.5) {\v};
    \else
        \node[cell]      at (-0.3+\i*0.75, 0.5) {\v};
    \fi
}
\node[font=\Large] at (3.55, 0.5) {$\xrightarrow{\;E\;}$};
\node[cell, emb] at (4.6,  0.5) {0.1};
\node[cell, emb] at (5.35, 0.5) {0.9};

% ── Fila: "tren" ─────────────────────────────────────────
\node[font=\small\itshape] at (-1.6, -0.3) {tren};
\foreach \v/\i in {0/0, 0/1, 0/2, 1/3}{
    \pgfmathparse{int(\v==1)}
    \ifnum\pgfmathresult=1
        \node[cell, hot] at (-0.3+\i*0.75, -0.3) {\v};
    \else
        \node[cell]      at (-0.3+\i*0.75, -0.3) {\v};
    \fi
}
\node[font=\Large] at (3.55, -0.3) {$\xrightarrow{\;E\;}$};
\node[cell, emb] at (4.6,  -0.3) {0.2};
\node[cell, emb] at (5.35, -0.3) {0.8};

% ── Nota ─────────────────────────────────────────────────
\node[font=\tiny, gray, align=left] at (1.8, -0.85)
    {La representación one-hot tiene distancia idéntica entre todos los pares.
     El embedding captura similitud.};

\end{tikzpicture}

\vspace{0.8em}

%──────────────────────────────────────────────────────────────
% PARTE INFERIOR: figura geométrica del espacio de embeddings
%──────────────────────────────────────────────────────────────
\begin{tikzpicture}[scale=1.05]

\node[font=\small\bfseries, gray] at (3.0, 4.1) {Espacio de embeddings ($d_{\text{model}}=2$)};

% Ejes
\draw[-{Stealth}] (-0.2, 0) -- (4.5, 0) node[right, font=\scriptsize] {dim 1};
\draw[-{Stealth}] (0, -0.2) -- (0, 4.0) node[above, font=\scriptsize] {dim 2};

% Puntos
\fill[red!60]    (3.2, 1.2) circle (4pt) node[right, font=\footnotesize] {\textit{gato}  $(0.8,\,0.3)$};
\fill[red!40]    (2.8, 1.6) circle (4pt) node[left,  font=\footnotesize] {\textit{perro} $(0.7,\,0.4)$};
\fill[blue!60]   (0.4, 3.6) circle (4pt) node[right, font=\footnotesize] {\textit{auto}  $(0.1,\,0.9)$};
\fill[blue!40]   (0.8, 3.2) circle (4pt) node[right, font=\footnotesize] {\textit{tren}  $(0.2,\,0.8)$};

% Elipses de agrupamiento
\draw[red!50,  dashed, thick] (3.0, 1.4) ellipse (0.55cm and 0.45cm);
\draw[blue!50, dashed, thick] (0.6, 3.4) ellipse (0.55cm and 0.45cm);

% Etiquetas de cluster
\node[font=\footnotesize, red!70]  at (3.75, 0.55) {animales};
\node[font=\footnotesize, blue!70] at (0.2,  2.55) {vehículos};

% Distancia entre clusters (flecha doble)
\draw[{Stealth}-{Stealth}, gray, thick] (2.55,1.45) -- (1.0,3.1);
\node[font=\tiny, gray, rotate=-52] at (1.65, 2.45) {distancia semántica};

% Nota
\node[font=\tiny, gray, align=center] at (3.0, -0.55)
    {Tokens semánticamente similares quedan próximos en el espacio.};

\end{tikzpicture}

\caption{Proyección de tokens mediante embeddings. \textbf{Arriba:} cada token
se representa primero como un vector \textit{one-hot} disperso y luego se
multiplica por la matriz de embeddings $E$, obteniendo un vector denso de
dimensión $d_{\text{model}}$. \textbf{Abajo:} en el espacio resultante, los
tokens semánticamente relacionados (\textit{gato}/\textit{perro},
\textit{auto}/\textit{tren}) quedan agrupados, mientras que categorías
distintas se separan geométricamente. Los valores numéricos son ilustrativos.}
\label{fig:proyeccion_embedding}
\end{figure}

% ============================================================
\begin{figure}[H]
\centering

%──────────────────────────────────────────────────────────────
% PARTE SUPERIOR: tabla numérica de la proyección
%──────────────────────────────────────────────────────────────
\begin{tikzpicture}[
    cell/.style={draw=gray!40, minimum width=0.72cm, minimum height=0.60cm,
                 font=\scriptsize, inner sep=1pt},
    hdr/.style={font=\footnotesize\bfseries},
    hot/.style={fill=blue!50, text=white},
    emb/.style={fill=orange!20}
]

% ── Encabezado columnas ───────────────────────────────────
\node[hdr] at (-1.6, 2.8) {Token};
\node[hdr] at ( 1.4, 2.8) {One-hot $\bm{v}$ $\,(|\mathcal{V}|=4)$};
\node[hdr] at ( 5.6, 2.8) {Embedding $\bm{e} = E\bm{v}$ $\,(d_{\text{model}}=2)$};

% ── Fila: "gato" ─────────────────────────────────────────
\node[font=\small\itshape] at (-1.6, 2.1) {gato};
\foreach \v/\i in {1/0, 0/1, 0/2, 0/3}{
    \pgfmathparse{int(\v==1)}
    \ifnum\pgfmathresult=1
        \node[cell, hot] at (-0.3+\i*0.75, 2.1) {\v};
    \else
        \node[cell]      at (-0.3+\i*0.75, 2.1) {\v};
    \fi
}
\node[font=\Large] at (3.55, 2.1) {$\xrightarrow{\;E\;}$};
\node[cell, emb] at (4.6,  2.1) {0.8};
\node[cell, emb] at (5.35, 2.1) {0.3};

% ── Fila: "perro" ────────────────────────────────────────
\node[font=\small\itshape] at (-1.6, 1.3) {perro};
\foreach \v/\i in {0/0, 1/1, 0/2, 0/3}{
    \pgfmathparse{int(\v==1)}
    \ifnum\pgfmathresult=1
        \node[cell, hot] at (-0.3+\i*0.75, 1.3) {\v};
    \else
        \node[cell]      at (-0.3+\i*0.75, 1.3) {\v};
    \fi
}
\node[font=\Large] at (3.55, 1.3) {$\xrightarrow{\;E\;}$};
\node[cell, emb] at (4.6,  1.3) {0.7};
\node[cell, emb] at (5.35, 1.3) {0.4};

% ── Fila: "auto" ─────────────────────────────────────────
\node[font=\small\itshape] at (-1.6, 0.5) {auto};
\foreach \v/\i in {0/0, 0/1, 1/2, 0/3}{
    \pgfmathparse{int(\v==1)}
    \ifnum\pgfmathresult=1
        \node[cell, hot] at (-0.3+\i*0.75, 0.5) {\v};
    \else
        \node[cell]      at (-0.3+\i*0.75, 0.5) {\v};
    \fi
}
\node[font=\Large] at (3.55, 0.5) {$\xrightarrow{\;E\;}$};
\node[cell, emb] at (4.6,  0.5) {0.1};
\node[cell, emb] at (5.35, 0.5) {0.9};

% ── Fila: "tren" ─────────────────────────────────────────
\node[font=\small\itshape] at (-1.6, -0.3) {tren};
\foreach \v/\i in {0/0, 0/1, 0/2, 1/3}{
    \pgfmathparse{int(\v==1)}
    \ifnum\pgfmathresult=1
        \node[cell, hot] at (-0.3+\i*0.75, -0.3) {\v};
    \else
        \node[cell]      at (-0.3+\i*0.75, -0.3) {\v};
    \fi
}
\node[font=\Large] at (3.55, -0.3) {$\xrightarrow{\;E\;}$};
\node[cell, emb] at (4.6,  -0.3) {0.2};
\node[cell, emb] at (5.35, -0.3) {0.8};

% ── Nota ─────────────────────────────────────────────────
\node[font=\tiny, gray, align=left] at (1.8, -0.85)
    {La representación one-hot tiene distancia idéntica entre todos los pares.
     El embedding captura similitud.};

\end{tikzpicture}

\vspace{0.8em}

%──────────────────────────────────────────────────────────────
% PARTE INFERIOR: figura geométrica del espacio de embeddings
%──────────────────────────────────────────────────────────────
\begin{tikzpicture}[scale=1.05]

\node[font=\small\bfseries, gray] at (3.0, 4.1) {Espacio de embeddings ($d_{\text{model}}=2$)};

% Ejes
\draw[-{Stealth}] (-0.2, 0) -- (4.5, 0) node[right, font=\scriptsize] {dim 1};
\draw[-{Stealth}] (0, -0.2) -- (0, 4.0) node[above, font=\scriptsize] {dim 2};

% Puntos
\fill[red!60]    (3.2, 1.2) circle (4pt) node[right, font=\footnotesize] {\textit{gato}  $(0.8,\,0.3)$};
\fill[red!40]    (2.8, 1.6) circle (4pt) node[left,  font=\footnotesize] {\textit{perro} $(0.7,\,0.4)$};
\fill[blue!60]   (0.4, 3.6) circle (4pt) node[right, font=\footnotesize] {\textit{auto}  $(0.1,\,0.9)$};
\fill[blue!40]   (0.8, 3.2) circle (4pt) node[right, font=\footnotesize] {\textit{tren}  $(0.2,\,0.8)$};

% Elipses de agrupamiento
\draw[red!50,  dashed, thick] (3.0, 1.4) ellipse (0.55cm and 0.45cm);
\draw[blue!50, dashed, thick] (0.6, 3.4) ellipse (0.55cm and 0.45cm);

% Etiquetas de cluster
\node[font=\footnotesize, red!70]  at (3.75, 0.55) {animales};
\node[font=\footnotesize, blue!70] at (0.2,  2.55) {vehículos};

% Distancia entre clusters (flecha doble)
\draw[{Stealth}-{Stealth}, gray, thick] (2.55,1.45) -- (1.0,3.1);
\node[font=\tiny, gray, rotate=-52] at (1.65, 2.45) {distancia semántica};

% Nota
\node[font=\tiny, gray, align=center] at (3.0, -0.55)
    {Tokens semánticamente similares quedan próximos en el espacio.};

\end{tikzpicture}

\caption{Proyección de tokens mediante embeddings. \textbf{Arriba:} cada token
se representa primero como un vector \textit{one-hot} disperso y luego se
multiplica por la matriz de embeddings $E$, obteniendo un vector denso de
dimensión $d_{\text{model}}$. \textbf{Abajo:} en el espacio resultante, los
tokens semánticamente relacionados (\textit{gato}/\textit{perro},
\textit{auto}/\textit{tren}) quedan agrupados, mientras que categorías
distintas se separan geométricamente. Los valores numéricos son ilustrativos.}
\label{fig:proyeccion_embedding}
\end{figure}


Los embeddings constituyen el primer bloque de todo Transformer: antes de
aplicar cualquier operación de atención, cada token de la secuencia de entrada
es \textit{proyectado} a este espacio continuo aprendido. La calidad de esta
representación inicial tiene impacto directo en la capacidad del modelo para
capturar relaciones lingüísticas.

%------------------------------------------------------------
\subsection{Redes recurrentes: el puente histórico}
\label{subsec:rnn}

Antes de la aparición de los Transformers, las redes neuronales recurrentes
(RNNs) y sus variantes —LSTM~\cite{hochreiter1997lstm} y GRU— representaban
el estado del arte en el modelado de secuencias. Estas arquitecturas
incorporan un estado oculto $h_t$ que se actualiza en cada paso temporal
acumulando información de la secuencia pasada:

\begin{equation}
    h_t = \varphi(W_h \, h_{t-1} + W_x \, x_t + b)
\end{equation}

Sin embargo, presentan dos problemas estructurales que los Transformers
resuelven directamente:

\begin{enumerate}
    \item \textbf{Gradiente desvanecido.} Al propagar el error hacia atrás
    en el tiempo, el gradiente se multiplica repetidamente por pesos pequeños
    y se atenúa exponencialmente, haciendo que las dependencias muy lejanas
    sean virtualmente invisibles para el optimizador.

    \item \textbf{Procesamiento secuencial obligatorio.} El paso $t$ depende
    del paso $t{-}1$: no hay forma de paralelizar el cómputo sobre la
    secuencia. En secuencias largas, esto hace el entrenamiento prohibitivamente
    lento.
\end{enumerate}

La eliminación de ambas limitaciones es una de las motivaciones centrales
del diseño de los Transformers.

%------------------------------------------------------------
\subsection{Síntesis comparativa}
\label{subsec:comparacion}

La Tabla~\ref{tab:comparacion_clasicos} resume las capacidades de cada
arquitectura frente a los requisitos del procesamiento del lenguaje natural.

\begin{table}[H]
\centering
\caption{Comparación de capacidades de los modelos clásicos frente a los
requisitos del procesamiento del lenguaje natural.}
\label{tab:comparacion_clasicos}
\begin{tabularx}{\textwidth}{lXXXX}
\toprule
\textbf{Modelo} & \textbf{Maneja secuencias} & \textbf{Contexto global} & \textbf{Dependencias largas} & \textbf{Paralelizable} \\
\midrule
Perceptrón           & No       & No       & No       & Sí      \\
MLP                  & Limitado & No       & No       & Sí      \\
Hopfield             & No       & Sí       & No       & No      \\
SOM                  & No       & No       & No       & Parcial \\
RNN / LSTM           & Sí       & Parcial  & Limitado & No      \\
\midrule
\textbf{Transformer} & \textbf{Sí} & \textbf{Sí} & \textbf{Sí} & \textbf{Sí} \\
\bottomrule
\end{tabularx}
\end{table}

La ausencia de un mecanismo capaz de combinar las cuatro propiedades de la
última fila es la limitación estructural que motiva la sección siguiente,
donde se caracteriza el problema del procesamiento del lenguaje natural.

\section{El problema del procesamiento del lenguaje natural}
\label{sec:problema_nlp}

El procesamiento del lenguaje natural (NLP, por sus siglas en inglés) tiene
como objetivo dotar a los sistemas computacionales de la capacidad de
comprender y generar lenguaje humano. A diferencia de otros dominios donde
las entradas son vectores numéricos bien definidos, el lenguaje presenta
propiedades estructurales que lo hacen particularmente difícil de modelar
con las arquitecturas revisadas en la sección anterior.

Esta sección caracteriza las cuatro propiedades del lenguaje natural que
resultan más problemáticas para los modelos clásicos y que motivan
directamente el diseño del mecanismo de atención.

%------------------------------------------------------------
\subsection{El lenguaje como secuencia de símbolos discretos}
\label{subsec:lenguaje_secuencia}

El lenguaje natural se expresa como secuencias de unidades discretas:
palabras, subpalabras o caracteres, denominados genéricamente
\textit{tokens}. A diferencia de una señal continua, cada token pertenece
a un vocabulario finito $\mathcal{V}$ y su representación numérica inmediata
(\textit{one-hot}) no captura ninguna relación con los demás tokens.

Más aún, el lenguaje es fundamentalmente \textit{variable en longitud}: una
oración puede tener 3 tokens o 300. Esto contrasta con los MLPs, que
requieren entradas de dimensión fija, y con las RNNs, que procesan la
secuencia de forma estrictamente secuencial. Cualquier arquitectura que
pretenda modelar lenguaje debe ser capaz de operar sobre secuencias de
longitud arbitraria sin perder información posicional ni contextual.

%------------------------------------------------------------
\subsection{Ambigüedad semántica}
\label{subsec:ambiguedad}

Una de las propiedades más características del lenguaje natural es su
ambigüedad: una misma forma superficial puede tener significados
completamente distintos dependiendo del contexto en que aparece. Esta
ambigüedad ocurre a múltiples niveles:

\begin{itemize}
    \item \textbf{Léxica:} una misma palabra tiene distintos significados
    (\textit{polisemia}). Ejemplo: \textit{banco} puede ser una institución
    financiera o un asiento exterior.

    \item \textbf{Sintáctica:} una misma oración admite distintas
    estructuras gramaticales. Ejemplo: \textit{``Vi al espía con el
    telescopio''} puede significar que yo usé el telescopio, o que el
    espía lo tenía.

    \item \textbf{Pragmática:} el significado depende del contexto
    situacional más allá del texto mismo.
\end{itemize}

% ============================================================
% fig_ambiguedad.tex: Ambigüedad semántica según contexto
% Uso: % ============================================================
% fig_ambiguedad.tex: Ambigüedad semántica según contexto
% Uso: \input{figures/fig_ambiguedad}
% ============================================================
\begin{figure}[H]
\centering
\begin{tikzpicture}[
    word/.style={draw=gray!40, rounded corners, fill=gray!8,
                 minimum width=1.1cm, minimum height=0.65cm, font=\small},
    key/.style={draw=blue!60, rounded corners, fill=blue!12,
                minimum width=1.1cm, minimum height=0.65cm, font=\small\bfseries},
    meaning/.style={draw, rounded corners, fill=white,
                    minimum width=2.6cm, minimum height=0.7cm,
                    font=\footnotesize, align=center},
    arrow/.style={-{Stealth}, thick},
    ctx/.style={font=\footnotesize\itshape, gray}
]

% ── Oración 1 ────────────────────────────────────────────────
\node[ctx] at (3.5, 5.6) {Contexto A:};
\foreach \w/\i in {Fui/0, al/1, banco/2, a/3, depositar/4, dinero/5}{
    \pgfmathparse{int(\i==2)}
    \ifnum\pgfmathresult=1
        \node[key]  at (\i*1.18, 4.95) {\w};
    \else
        \node[word] at (\i*1.18, 4.95) {\w};
    \fi
}
\draw[arrow, blue!60] (2.36, 4.62) -- (2.36, 4.0);
\node[meaning, fill=green!10, draw=green!50!black]
    at (2.36, 3.6) {institución\\financiera};

% ── Oración 2 ────────────────────────────────────────────────
\node[ctx] at (3.5, 2.9) {Contexto B:};
\foreach \w/\i in {Me/0, senté/1, en/2, el/3, banco/4, del/5, parque/6}{
    \pgfmathparse{int(\i==4)}
    \ifnum\pgfmathresult=1
        \node[key]  at (\i*1.1, 2.3) {\w};
    \else
        \node[word] at (\i*1.1, 2.3) {\w};
    \fi
}
\draw[arrow, blue!60] (4.4, 1.97) -- (4.4, 1.35);
\node[meaning, fill=orange!10, draw=orange!60]
    at (4.4, 0.95) {asiento\\exterior};

% ── Anotación central ────────────────────────────────────────
\node[draw=blue!20, fill=blue!5, rounded corners, font=\footnotesize,
      inner sep=6pt, align=center] at (7.8, 2.8)
    {\textbf{``banco''} es la misma\\palabra en ambos casos.\\El significado\\
     solo emerge del\\contexto circundante.};

\draw[dashed, gray!50] (5.9, 3.6) -- (6.8, 3.2);
\draw[dashed, gray!50] (5.8, 0.95) -- (6.8, 2.4);

\end{tikzpicture}
\caption{Ambigüedad semántica léxica. La palabra \textit{banco} tiene
significados completamente distintos según el contexto oracional. Un modelo
que procese cada token de forma independiente, sin considerar las palabras
circundantes, no puede resolver esta ambigüedad.}
\label{fig:ambiguedad}
\end{figure}

% ============================================================
\begin{figure}[H]
\centering
\begin{tikzpicture}[
    word/.style={draw=gray!40, rounded corners, fill=gray!8,
                 minimum width=1.1cm, minimum height=0.65cm, font=\small},
    key/.style={draw=blue!60, rounded corners, fill=blue!12,
                minimum width=1.1cm, minimum height=0.65cm, font=\small\bfseries},
    meaning/.style={draw, rounded corners, fill=white,
                    minimum width=2.6cm, minimum height=0.7cm,
                    font=\footnotesize, align=center},
    arrow/.style={-{Stealth}, thick},
    ctx/.style={font=\footnotesize\itshape, gray}
]

% ── Oración 1 ────────────────────────────────────────────────
\node[ctx] at (3.5, 5.6) {Contexto A:};
\foreach \w/\i in {Fui/0, al/1, banco/2, a/3, depositar/4, dinero/5}{
    \pgfmathparse{int(\i==2)}
    \ifnum\pgfmathresult=1
        \node[key]  at (\i*1.18, 4.95) {\w};
    \else
        \node[word] at (\i*1.18, 4.95) {\w};
    \fi
}
\draw[arrow, blue!60] (2.36, 4.62) -- (2.36, 4.0);
\node[meaning, fill=green!10, draw=green!50!black]
    at (2.36, 3.6) {institución\\financiera};

% ── Oración 2 ────────────────────────────────────────────────
\node[ctx] at (3.5, 2.9) {Contexto B:};
\foreach \w/\i in {Me/0, senté/1, en/2, el/3, banco/4, del/5, parque/6}{
    \pgfmathparse{int(\i==4)}
    \ifnum\pgfmathresult=1
        \node[key]  at (\i*1.1, 2.3) {\w};
    \else
        \node[word] at (\i*1.1, 2.3) {\w};
    \fi
}
\draw[arrow, blue!60] (4.4, 1.97) -- (4.4, 1.35);
\node[meaning, fill=orange!10, draw=orange!60]
    at (4.4, 0.95) {asiento\\exterior};

% ── Anotación central ────────────────────────────────────────
\node[draw=blue!20, fill=blue!5, rounded corners, font=\footnotesize,
      inner sep=6pt, align=center] at (7.8, 2.8)
    {\textbf{``banco''} es la misma\\palabra en ambos casos.\\El significado\\
     solo emerge del\\contexto circundante.};

\draw[dashed, gray!50] (5.9, 3.6) -- (6.8, 3.2);
\draw[dashed, gray!50] (5.8, 0.95) -- (6.8, 2.4);

\end{tikzpicture}
\caption{Ambigüedad semántica léxica. La palabra \textit{banco} tiene
significados completamente distintos según el contexto oracional. Un modelo
que procese cada token de forma independiente, sin considerar las palabras
circundantes, no puede resolver esta ambigüedad.}
\label{fig:ambiguedad}
\end{figure}


La consecuencia directa para el modelado es que no existe una representación
fija correcta para un token: su significado es una función de su contexto.
Un modelo que asigne a cada token un vector de representación estático
(independiente del resto de la secuencia) no puede resolver esta ambigüedad.
Esto descarta, en principio, cualquier arquitectura sin mecanismo de
contextualización dinámica.

%------------------------------------------------------------
\subsection{Dependencias de largo alcance}
\label{subsec:dependencias_largas}

En el lenguaje natural, dos elementos de una oración pueden estar
relacionados gramaticalmente o semánticamente aunque estén separados por
muchos otros tokens. Estas relaciones se denominan \textit{dependencias de
largo alcance} y son ubicuas en el lenguaje real~\cite{bengio2003neural}.

Un ejemplo paradigmático es la concordancia sujeto-verbo en oraciones con
cláusulas subordinadas intercaladas:

\begin{quote}
\textit{``El científico que ganó el premio Nobel de física presentó su investigación.''}
\end{quote}

El verbo \textit{presentó} debe concordar en número y persona con el sujeto
\textit{científico}, aunque ambos estén separados por la cláusula relativa
\textit{``que ganó el premio Nobel de física''} (6 tokens). Un modelo con
ventana de contexto fija menor que esa distancia nunca puede aprender esta
relación correctamente.

% ============================================================
% fig_dependencias.tex — Dependencias largas y correferencia
% Uso: % ============================================================
% fig_dependencias.tex — Dependencias largas y correferencia
% Uso: \input{figures/fig_dependencias}
% ============================================================
\begin{figure}[H]
\centering
\begin{tikzpicture}[
    word/.style={draw=gray!35, rounded corners, fill=gray!8,
                 minimum width=1.05cm, minimum height=0.62cm,
                 font=\scriptsize, inner sep=2pt},
    subj/.style={draw=blue!60, rounded corners, fill=blue!12,
                 minimum width=1.05cm, minimum height=0.62cm,
                 font=\scriptsize\bfseries, inner sep=2pt},
    verb/.style={draw=red!60, rounded corners, fill=red!10,
                 minimum width=1.05cm, minimum height=0.62cm,
                 font=\scriptsize\bfseries, inner sep=2pt},
    pron/.style={draw=orange!70, rounded corners, fill=orange!12,
                 minimum width=1.05cm, minimum height=0.62cm,
                 font=\scriptsize\bfseries, inner sep=2pt},
    lbl/.style={font=\tiny, align=center}
]

% ── Tokens colocados manualmente (12 tokens) ─────────────────
\node[subj] (t0)  at (0*1.08, 3.2) {El};
\node[subj] (t1)  at (1*1.08, 3.2) {cient.};
\node[word] (t2)  at (2*1.08, 3.2) {que};
\node[word] (t3)  at (3*1.08, 3.2) {gan\'o};
\node[word] (t4)  at (4*1.08, 3.2) {el};
\node[word] (t5)  at (5*1.08, 3.2) {premio};
\node[word] (t6)  at (6*1.08, 3.2) {Nobel};
\node[word] (t7)  at (7*1.08, 3.2) {de};
\node[word] (t8)  at (8*1.08, 3.2) {f\'isica};
\node[verb] (t9)  at (9*1.08, 3.2) {present\'o};
\node[pron] (t10) at (10*1.08,3.2) {su};
\node[word] (t11) at (11*1.08,3.2) {invest.};

% ── Arco: sujeto → verbo (dependencia larga) ─────────────────
\draw[-{Stealth}, blue!60, thick]
    (t1.north) .. controls (1.08, 4.6) and (9.72, 4.6) .. (t9.north);
\node[lbl, blue!70] at (5.4, 4.85)
    {dependencia sujeto--verbo \\ (8 tokens de distancia)};

% ── Arco: pronombre → sujeto (correferencia) ─────────────────
\draw[-{Stealth}, orange!80, thick]
    (t10.south) .. controls (10.8, 1.8) and (1.08, 1.8) .. (t1.south);
\node[lbl, orange!70] at (5.4, 1.4)
    {\textit{``su''} refiere a \textit{``cient.''} \\ (correferencia a distancia)};

% ── Leyenda ──────────────────────────────────────────────────
\node[subj, label=right:{\scriptsize sujeto}]    at (0.0,  0.55) {};
\node[verb, label=right:{\scriptsize verbo}]     at (2.5,  0.55) {};
\node[pron, label=right:{\scriptsize pronombre}] at (5.0,  0.55) {};

% ── Nota ─────────────────────────────────────────────────────
\node[draw=gray!30, fill=gray!5, rounded corners, font=\tiny,
      inner sep=5pt, align=left] at (9.5, 0.55)
    {Un MLP con ventana fija no puede\\ conectar tokens tan alejados.\\ Un Transformer s\'i: la atenci\'on es global.};

\end{tikzpicture}
\caption{Dependencia de largo alcance y correferencia. El verbo
\textit{present\'o} concuerda con \textit{cient\'ifico}, su sujeto,
separados por 8 tokens. El pronombre \textit{su} refiere tambi\'en al mismo
sujeto. Ambas relaciones requieren acceder a contexto distante.}
\label{fig:dependencias}
\end{figure}

% ============================================================
\begin{figure}[H]
\centering
\begin{tikzpicture}[
    word/.style={draw=gray!35, rounded corners, fill=gray!8,
                 minimum width=1.05cm, minimum height=0.62cm,
                 font=\scriptsize, inner sep=2pt},
    subj/.style={draw=blue!60, rounded corners, fill=blue!12,
                 minimum width=1.05cm, minimum height=0.62cm,
                 font=\scriptsize\bfseries, inner sep=2pt},
    verb/.style={draw=red!60, rounded corners, fill=red!10,
                 minimum width=1.05cm, minimum height=0.62cm,
                 font=\scriptsize\bfseries, inner sep=2pt},
    pron/.style={draw=orange!70, rounded corners, fill=orange!12,
                 minimum width=1.05cm, minimum height=0.62cm,
                 font=\scriptsize\bfseries, inner sep=2pt},
    lbl/.style={font=\tiny, align=center}
]

% ── Tokens colocados manualmente (12 tokens) ─────────────────
\node[subj] (t0)  at (0*1.08, 3.2) {El};
\node[subj] (t1)  at (1*1.08, 3.2) {cient.};
\node[word] (t2)  at (2*1.08, 3.2) {que};
\node[word] (t3)  at (3*1.08, 3.2) {gan\'o};
\node[word] (t4)  at (4*1.08, 3.2) {el};
\node[word] (t5)  at (5*1.08, 3.2) {premio};
\node[word] (t6)  at (6*1.08, 3.2) {Nobel};
\node[word] (t7)  at (7*1.08, 3.2) {de};
\node[word] (t8)  at (8*1.08, 3.2) {f\'isica};
\node[verb] (t9)  at (9*1.08, 3.2) {present\'o};
\node[pron] (t10) at (10*1.08,3.2) {su};
\node[word] (t11) at (11*1.08,3.2) {invest.};

% ── Arco: sujeto → verbo (dependencia larga) ─────────────────
\draw[-{Stealth}, blue!60, thick]
    (t1.north) .. controls (1.08, 4.6) and (9.72, 4.6) .. (t9.north);
\node[lbl, blue!70] at (5.4, 4.85)
    {dependencia sujeto--verbo \\ (8 tokens de distancia)};

% ── Arco: pronombre → sujeto (correferencia) ─────────────────
\draw[-{Stealth}, orange!80, thick]
    (t10.south) .. controls (10.8, 1.8) and (1.08, 1.8) .. (t1.south);
\node[lbl, orange!70] at (5.4, 1.4)
    {\textit{``su''} refiere a \textit{``cient.''} \\ (correferencia a distancia)};

% ── Leyenda ──────────────────────────────────────────────────
\node[subj, label=right:{\scriptsize sujeto}]    at (0.0,  0.55) {};
\node[verb, label=right:{\scriptsize verbo}]     at (2.5,  0.55) {};
\node[pron, label=right:{\scriptsize pronombre}] at (5.0,  0.55) {};

% ── Nota ─────────────────────────────────────────────────────
\node[draw=gray!30, fill=gray!5, rounded corners, font=\tiny,
      inner sep=5pt, align=left] at (9.5, 0.55)
    {Un MLP con ventana fija no puede\\ conectar tokens tan alejados.\\ Un Transformer s\'i: la atenci\'on es global.};

\end{tikzpicture}
\caption{Dependencia de largo alcance y correferencia. El verbo
\textit{present\'o} concuerda con \textit{cient\'ifico}, su sujeto,
separados por 8 tokens. El pronombre \textit{su} refiere tambi\'en al mismo
sujeto. Ambas relaciones requieren acceder a contexto distante.}
\label{fig:dependencias}
\end{figure}


Formalmente, si definimos la \textit{longitud de dependencia} entre dos
tokens $i$ y $j$ como $|i - j|$, los modelos de ventana fija de ancho $k$
son ciegos a toda dependencia con $|i - j| > k$. En textos reales, las
dependencias relevantes pueden extenderse a decenas de tokens~\cite{vaswani2017attention}.

%------------------------------------------------------------
\subsection{Orden y estructura sintáctica}
\label{subsec:orden}

El orden de los tokens en una secuencia es portador de significado: no es
simplemente un atributo accidental sino una dimensión informativa
fundamental. Considérense los pares:

\begin{center}
\begin{tabular}{ll}
\textit{``El perro muerde al hombre''} & $\neq$ \textit{``El hombre muerde al perro''} \\[2pt]
\textit{``No es verdad que miente''}   & $\neq$ \textit{``Es verdad que no miente''}   \\
\end{tabular}
\end{center}

Ambos pares contienen exactamente los mismos tokens; el significado cambia
radicalmente con el orden. Esto implica que cualquier modelo que procese los
tokens como un conjunto desordenado (es decir, que sea \textit{invariante a
permutaciones}) es estructuralmente incapaz de modelar lenguaje de forma
correcta~\cite{goodfellow2016deeplearning}.

Los MLPs son invariantes a permutaciones salvo que la posición esté
codificada explícitamente en la entrada. Las RNNs incorporan el orden
implícitamente a través del estado oculto secuencial, pero al costo de
procesar los tokens uno a uno. Los Transformers deben resolver este problema
de una tercera forma: codificando el orden \textit{explícitamente} sin
sacrificar el paralelismo. Esta solución, la codificación posicional, se
describe en la Sección~\ref{sec:arquitectura_transformer}.

%------------------------------------------------------------
\subsection{Correferencialidad}
\label{subsec:correferencia}

La \textit{correferencialidad} ocurre cuando distintas expresiones en un
texto refieren a la misma entidad. El caso más común son los pronombres:

\begin{quote}
\textit{``La investigadora publicó sus resultados. \textbf{Ella} había
trabajado durante diez años en el problema.''}
\end{quote}

Para interpretar correctamente \textit{ella}, el modelo debe establecer que
este pronombre corefiere con \textit{investigadora}, mencionada en la oración
anterior. Esta relación puede extenderse a través de múltiples oraciones y
requiere rastrear entidades a lo largo de un contexto potencialmente largo.

La correferencia es conceptualmente análoga al problema de Hopfield: así
como en una red de Hopfield el estado de cada unidad está determinado por el
estado global del sistema, resolver la correferencia requiere que cada token
sea interpretado en función de todos los demás. La diferencia clave es la
escala y la dirección: mientras Hopfield opera sobre patrones binarios fijos,
el lenguaje requiere relaciones dinámicas y asimétricas entre tokens.

%------------------------------------------------------------
\subsection{Análisis integrado: por qué el contexto es indispensable}
\label{subsec:contexto_indispensable}

Las cuatro propiedades anteriores convergen en una conclusión unificada:
\textbf{el significado de un token no puede determinarse en aislamiento}. Es
siempre una función del contexto en que aparece.

La siguiente figura analiza paso a paso cómo el contexto es imprescindible
para resolver incluso una ambigüedad simple:

% ============================================================
% fig_ejemplo_contexto.tex — Análisis paso a paso de una oración
% Uso: % ============================================================
% fig_ejemplo_contexto.tex — Análisis paso a paso de una oración
% Uso: \input{figures/fig_ejemplo_contexto}
% ============================================================
\begin{figure}[H]
\centering
\begin{tikzpicture}[
    word/.style={draw=gray!40, rounded corners, fill=gray!8,
                 minimum width=1.0cm, minimum height=0.6cm, font=\small},
    focus/.style={draw=blue!60, rounded corners, fill=blue!12,
                  minimum width=1.0cm, minimum height=0.6cm,
                  font=\small\bfseries},
    ctx/.style={draw=green!50!black, rounded corners, fill=green!8,
                minimum width=1.0cm, minimum height=0.6cm, font=\small},
    box/.style={draw=gray!40, rounded corners, fill=white,
                inner sep=6pt, font=\footnotesize, align=left},
    arrow/.style={-{Stealth}, thick, gray},
    step/.style={font=\small\bfseries, blue!70}
]

% ── Oración completa (referencia) ───────────────────────────
\node[font=\small\bfseries, gray] at (3.5, 6.8)
    {Oración: \textit{``Vi a la estrella con el telescopio''}};

% ── PASO 1: token aislado ────────────────────────────────────
\node[step] at (-0.5, 5.9) {Paso 1};
\node[box]  at (4.2, 5.85) {Sin contexto: \textbf{``estrella''} puede ser\\
    \textcolor{blue!70}{$\bullet$ astro celeste} \quad
    \textcolor{orange!80}{$\bullet$ persona famosa}};

\node[focus] at (1.8, 5.85) {estrella};
\node[font=\tiny, red!60] at (1.8, 5.45) {ambigua};

% ── PASO 2: ventana pequeña (±2) ─────────────────────────────
\node[step] at (-0.5, 4.6) {Paso 2};
\node[box]  at (4.5, 4.55) {Ventana $\pm 2$: \textit{``la estrella con el''}\\
    Sigue ambigua: sin información resolutoria.};

\foreach \w/\i/\st in {la/0/ctx, estrella/1/focus, con/2/ctx, el/3/ctx}{
    \pgfmathparse{int(\i==1)}
    \ifnum\pgfmathresult=1
        \node[focus] at (\i*1.1, 4.55) {\w};
    \else
        \node[ctx]   at (\i*1.1, 4.55) {\w};
    \fi
}
\node[font=\tiny, red!60] at (1.1, 4.15) {aún ambigua};

% ── PASO 3: oración completa ────────────────────────────────
\node[step] at (-0.5, 3.2) {Paso 3};
\node[box, fill=green!5] at (5.0, 3.15)
    {Oración completa: \textit{``con el telescopio''}\\
    $\Rightarrow$ instrumento óptico $\Rightarrow$
    \textcolor{green!50!black}{\textbf{astro celeste}}};

\foreach \w/\i in {Vi/0, a/1, la/2, estrella/3, con/4, el/5, telescopio/6}{
    \pgfmathparse{int(\i==3)}
    \ifnum\pgfmathresult=1
        \node[focus] at (\i*1.05, 3.15) {\w};
    \else
        \pgfmathparse{int(\i==6 || \i==5 || \i==4)}
        \ifnum\pgfmathresult=1
            \node[ctx, draw=green!50!black, fill=green!10] at (\i*1.05, 3.15) {\w};
        \else
            \node[word] at (\i*1.05, 3.15) {\w};
        \fi
    \fi
}

% Arco resolutorio
\draw[-{Stealth}, green!60!black, thick]
    (6.3, 2.82) .. controls (6.3, 2.2) and (3.15, 2.2) .. (3.15, 2.82);
\node[font=\tiny, green!50!black] at (4.7, 1.95)
    {\textit{``telescopio''} resuelve la ambigüedad de \textit{``estrella''}};

% ── Conclusión ───────────────────────────────────────────────
\node[draw=blue!25, fill=blue!5, rounded corners, font=\footnotesize,
      inner sep=7pt, align=center] at (3.5, 1.1)
    {El significado correcto requiere acceso a \textbf{toda} la oración,\\
     no solo a una ventana local. Un modelo que no tenga acceso\\
     global al contexto no puede resolver esta ambigüedad.};

\end{tikzpicture}
\caption{Análisis paso a paso de la resolución de ambigüedad en la oración
\textit{``Vi a la estrella con el telescopio''}. En el Paso 1, el token
\textit{estrella} es ambiguo en aislamiento. En el Paso 2, una ventana local
no aporta información resolutoria. Solo en el Paso 3, con acceso a toda la
oración, la palabra \textit{telescopio} permite inferir que \textit{estrella}
se refiere a un astro y no a una persona famosa.}
\label{fig:ejemplo_contexto}
\end{figure}

% ============================================================
\begin{figure}[H]
\centering
\begin{tikzpicture}[
    word/.style={draw=gray!40, rounded corners, fill=gray!8,
                 minimum width=1.0cm, minimum height=0.6cm, font=\small},
    focus/.style={draw=blue!60, rounded corners, fill=blue!12,
                  minimum width=1.0cm, minimum height=0.6cm,
                  font=\small\bfseries},
    ctx/.style={draw=green!50!black, rounded corners, fill=green!8,
                minimum width=1.0cm, minimum height=0.6cm, font=\small},
    box/.style={draw=gray!40, rounded corners, fill=white,
                inner sep=6pt, font=\footnotesize, align=left},
    arrow/.style={-{Stealth}, thick, gray},
    step/.style={font=\small\bfseries, blue!70}
]

% ── Oración completa (referencia) ───────────────────────────
\node[font=\small\bfseries, gray] at (3.5, 6.8)
    {Oración: \textit{``Vi a la estrella con el telescopio''}};

% ── PASO 1: token aislado ────────────────────────────────────
\node[step] at (-0.5, 5.9) {Paso 1};
\node[box]  at (4.2, 5.85) {Sin contexto: \textbf{``estrella''} puede ser\\
    \textcolor{blue!70}{$\bullet$ astro celeste} \quad
    \textcolor{orange!80}{$\bullet$ persona famosa}};

\node[focus] at (1.8, 5.85) {estrella};
\node[font=\tiny, red!60] at (1.8, 5.45) {ambigua};

% ── PASO 2: ventana pequeña (±2) ─────────────────────────────
\node[step] at (-0.5, 4.6) {Paso 2};
\node[box]  at (4.5, 4.55) {Ventana $\pm 2$: \textit{``la estrella con el''}\\
    Sigue ambigua: sin información resolutoria.};

\foreach \w/\i/\st in {la/0/ctx, estrella/1/focus, con/2/ctx, el/3/ctx}{
    \pgfmathparse{int(\i==1)}
    \ifnum\pgfmathresult=1
        \node[focus] at (\i*1.1, 4.55) {\w};
    \else
        \node[ctx]   at (\i*1.1, 4.55) {\w};
    \fi
}
\node[font=\tiny, red!60] at (1.1, 4.15) {aún ambigua};

% ── PASO 3: oración completa ────────────────────────────────
\node[step] at (-0.5, 3.2) {Paso 3};
\node[box, fill=green!5] at (5.0, 3.15)
    {Oración completa: \textit{``con el telescopio''}\\
    $\Rightarrow$ instrumento óptico $\Rightarrow$
    \textcolor{green!50!black}{\textbf{astro celeste}}};

\foreach \w/\i in {Vi/0, a/1, la/2, estrella/3, con/4, el/5, telescopio/6}{
    \pgfmathparse{int(\i==3)}
    \ifnum\pgfmathresult=1
        \node[focus] at (\i*1.05, 3.15) {\w};
    \else
        \pgfmathparse{int(\i==6 || \i==5 || \i==4)}
        \ifnum\pgfmathresult=1
            \node[ctx, draw=green!50!black, fill=green!10] at (\i*1.05, 3.15) {\w};
        \else
            \node[word] at (\i*1.05, 3.15) {\w};
        \fi
    \fi
}

% Arco resolutorio
\draw[-{Stealth}, green!60!black, thick]
    (6.3, 2.82) .. controls (6.3, 2.2) and (3.15, 2.2) .. (3.15, 2.82);
\node[font=\tiny, green!50!black] at (4.7, 1.95)
    {\textit{``telescopio''} resuelve la ambigüedad de \textit{``estrella''}};

% ── Conclusión ───────────────────────────────────────────────
\node[draw=blue!25, fill=blue!5, rounded corners, font=\footnotesize,
      inner sep=7pt, align=center] at (3.5, 1.1)
    {El significado correcto requiere acceso a \textbf{toda} la oración,\\
     no solo a una ventana local. Un modelo que no tenga acceso\\
     global al contexto no puede resolver esta ambigüedad.};

\end{tikzpicture}
\caption{Análisis paso a paso de la resolución de ambigüedad en la oración
\textit{``Vi a la estrella con el telescopio''}. En el Paso 1, el token
\textit{estrella} es ambiguo en aislamiento. En el Paso 2, una ventana local
no aporta información resolutoria. Solo en el Paso 3, con acceso a toda la
oración, la palabra \textit{telescopio} permite inferir que \textit{estrella}
se refiere a un astro y no a una persona famosa.}
\label{fig:ejemplo_contexto}
\end{figure}


Este ejemplo ilustra un principio general: un modelo que procese cada token
de forma independiente, o que solo tenga acceso a una ventana local, no puede
garantizar la interpretación correcta. Se requiere un mecanismo que permita
que \textit{cualquier} token acceda a la información de \textit{cualquier
otro} token de la secuencia, sin restricción de distancia.

%------------------------------------------------------------
\subsection{Requerimientos para un modelo de lenguaje adecuado}
\label{subsec:requerimientos}

Las propiedades del lenguaje descritas en esta sección imponen requerimientos
concretos a cualquier arquitectura que pretenda modelarlo correctamente.
La Tabla~\ref{tab:requerimientos_nlp} los resume y los contrasta con las
capacidades de los modelos clásicos ya analizados.

\begin{table}[H]
\centering
\caption{Requerimientos del modelado del lenguaje natural y cobertura por
arquitectura.}
\label{tab:requerimientos_nlp}
\begin{tabularx}{\textwidth}{lXcc}
\toprule
\textbf{Requerimiento} & \textbf{Motivación} & \textbf{Modelos clásicos} & \textbf{Transformer} \\
\midrule
Longitud variable        & Secuencias arbitrarias        & Limitado / No & Sí \\
Representación contextual & Ambigüedad semántica          & No            & Sí \\
Dependencias globales    & Largo alcance, correferencia  & No            & Sí \\
Sensibilidad al orden    & Estructura sintáctica         & Parcial       & Sí \\
Paralelización           & Eficiencia en entrenamiento   & No (RNN)      & Sí \\
\bottomrule
\end{tabularx}
\end{table}

Ninguna de las arquitecturas clásicas cubre simultáneamente todos estos
requerimientos. La siguiente sección introduce el mecanismo que permite
satisfacerlos: la \textit{atención}.

\section{El mecanismo de atención}
\label{sec:mecanismo_atencion}

Las secciones anteriores identificaron la limitación central de los modelos
clásicos: ninguno dispone de un mecanismo que permita que cada elemento de
una secuencia acceda, de forma dinámica y dependiente del contenido, a la
información de cualquier otro elemento. Esta sección presenta la solución
que los Transformers proponen: el \textit{mecanismo de atención}.

%------------------------------------------------------------
\subsection{Intuición: ¿qué significa ``prestar atención''?}
\label{subsec:atencion_intuicion}

La idea central es sencilla. Cuando un ser humano lee la oración
\textit{``El gato comió el ratón''} y quiere entender el verbo
\textit{comió}, no lo procesa en aislamiento: busca activamente quién
realiza la acción (sujeto: \textit{gato}) y sobre qué recae (objeto:
\textit{ratón}). Esta capacidad de enfocarse selectivamente en las partes
relevantes del contexto es exactamente lo que el mecanismo de atención
implementa de forma computacional.

Formalmente, para cada token de la secuencia, la atención calcula un
\textit{peso} $\alpha_{ij}$ que indica cuánta relevancia tiene el token $j$
para representar al token $i$. La representación final del token $i$ es
entonces una \textit{suma ponderada} de las representaciones de todos los
demás tokens:

\begin{equation}
    \bm{z}_i = \sum_{j=1}^{n} \alpha_{ij}\, \bm{v}_j
    \qquad \text{con} \quad \alpha_{ij} \geq 0, \quad \sum_j \alpha_{ij} = 1
\end{equation}

% ============================================================
% fig_atencion_intuicion.tex — Intuición del mecanismo de atención
% Uso: % ============================================================
% fig_atencion_intuicion.tex — Intuición del mecanismo de atención
% Uso: \input{figures/fig_atencion_intuicion}
% ============================================================
\begin{figure}[H]
\centering
\begin{tikzpicture}[
    word/.style={draw=gray!40, rounded corners, fill=gray!8,
                 minimum width=1.1cm, minimum height=0.62cm, font=\small},
    focus/.style={draw=blue!60, rounded corners, fill=blue!12,
                  minimum width=1.1cm, minimum height=0.62cm,
                  font=\small\bfseries},
    weight/.style={font=\footnotesize\bfseries},
    bar/.style={fill=blue!50, draw=blue!70}
]

\node[font=\small\bfseries, gray] at (3.3, 5.6)
    {Token consultado: \textit{``comi\'o''}};

\node[word]  at (0*1.25, 5.0) {El};
\node[word]  at (1*1.25, 5.0) {gato};
\node[focus] at (2*1.25, 5.0) {comi\'o};
\node[word]  at (3*1.25, 5.0) {el};
\node[word]  at (4*1.25, 5.0) {rat\'on};

\node[font=\small\bfseries, gray] at (3.3, 4.2)
    {Pesos de atenci\'on $\alpha_i$ asignados por \textit{``comi\'o''}:};

\node[font=\small\itshape] at (-0.3, 3.55) {El};
\fill[bar, fill opacity=0.35] (0.2, 3.25) rectangle (0.475, 3.55);
\node[weight, blue!70] at (0.9, 3.40) {0.05};

\node[font=\small\itshape] at (-0.3, 2.93) {gato};
\fill[bar, fill opacity=0.685] (0.2, 2.63) rectangle (3.225, 2.93);
\node[weight, blue!70] at (3.65, 2.78) {0.55};

\node[font=\small\itshape] at (-0.3, 2.31) {comi\'o};
\fill[bar, fill opacity=0.40] (0.2, 2.01) rectangle (0.75, 2.31);
\node[weight, blue!70] at (1.15, 2.16) {0.10};

\node[font=\small\itshape] at (-0.3, 1.69) {el};
\fill[bar, fill opacity=0.35] (0.2, 1.39) rectangle (0.475, 1.69);
\node[weight, blue!70] at (0.9, 1.54) {0.05};

\node[font=\small\itshape] at (-0.3, 1.07) {rat\'on};
\fill[bar, fill opacity=0.525] (0.2, 0.77) rectangle (1.575, 1.07);
\node[weight, blue!70] at (2.0, 0.92) {0.25};

\node[draw=blue!20, fill=blue!5, rounded corners, font=\footnotesize,
      inner sep=7pt, align=left] at (7.8, 2.5)
    {\textit{``comi\'o''} presta m\'as atenci\'on\\
     a \textit{``gato''} (sujeto, 0.55) y\\
     \textit{``rat\'on''} (objeto, 0.25).\\[4pt]
     Los pesos son \textbf{din\'amicos}:\\
     dependen del contenido,\\
     no de la posici\'on.};

\node[font=\footnotesize, gray] at (3.3, 0.2)
    {Representaci\'on contextualizada:
     $\bm{z} = \sum_i \alpha_i \, \bm{v}_i$};

\end{tikzpicture}
\caption{Intuici\'on del mecanismo de atenci\'on. Para el token
\textit{comi\'o}, la atenci\'on calcula un peso $\alpha_i$ para cada
token de la secuencia. La representaci\'on final es una suma ponderada
de los valores, donde los pesos reflejan la relevancia sem\'antica.}
\label{fig:atencion_intuicion}
\end{figure}

% ============================================================
\begin{figure}[H]
\centering
\begin{tikzpicture}[
    word/.style={draw=gray!40, rounded corners, fill=gray!8,
                 minimum width=1.1cm, minimum height=0.62cm, font=\small},
    focus/.style={draw=blue!60, rounded corners, fill=blue!12,
                  minimum width=1.1cm, minimum height=0.62cm,
                  font=\small\bfseries},
    weight/.style={font=\footnotesize\bfseries},
    bar/.style={fill=blue!50, draw=blue!70}
]

\node[font=\small\bfseries, gray] at (3.3, 5.6)
    {Token consultado: \textit{``comi\'o''}};

\node[word]  at (0*1.25, 5.0) {El};
\node[word]  at (1*1.25, 5.0) {gato};
\node[focus] at (2*1.25, 5.0) {comi\'o};
\node[word]  at (3*1.25, 5.0) {el};
\node[word]  at (4*1.25, 5.0) {rat\'on};

\node[font=\small\bfseries, gray] at (3.3, 4.2)
    {Pesos de atenci\'on $\alpha_i$ asignados por \textit{``comi\'o''}:};

\node[font=\small\itshape] at (-0.3, 3.55) {El};
\fill[bar, fill opacity=0.35] (0.2, 3.25) rectangle (0.475, 3.55);
\node[weight, blue!70] at (0.9, 3.40) {0.05};

\node[font=\small\itshape] at (-0.3, 2.93) {gato};
\fill[bar, fill opacity=0.685] (0.2, 2.63) rectangle (3.225, 2.93);
\node[weight, blue!70] at (3.65, 2.78) {0.55};

\node[font=\small\itshape] at (-0.3, 2.31) {comi\'o};
\fill[bar, fill opacity=0.40] (0.2, 2.01) rectangle (0.75, 2.31);
\node[weight, blue!70] at (1.15, 2.16) {0.10};

\node[font=\small\itshape] at (-0.3, 1.69) {el};
\fill[bar, fill opacity=0.35] (0.2, 1.39) rectangle (0.475, 1.69);
\node[weight, blue!70] at (0.9, 1.54) {0.05};

\node[font=\small\itshape] at (-0.3, 1.07) {rat\'on};
\fill[bar, fill opacity=0.525] (0.2, 0.77) rectangle (1.575, 1.07);
\node[weight, blue!70] at (2.0, 0.92) {0.25};

\node[draw=blue!20, fill=blue!5, rounded corners, font=\footnotesize,
      inner sep=7pt, align=left] at (7.8, 2.5)
    {\textit{``comi\'o''} presta m\'as atenci\'on\\
     a \textit{``gato''} (sujeto, 0.55) y\\
     \textit{``rat\'on''} (objeto, 0.25).\\[4pt]
     Los pesos son \textbf{din\'amicos}:\\
     dependen del contenido,\\
     no de la posici\'on.};

\node[font=\footnotesize, gray] at (3.3, 0.2)
    {Representaci\'on contextualizada:
     $\bm{z} = \sum_i \alpha_i \, \bm{v}_i$};

\end{tikzpicture}
\caption{Intuici\'on del mecanismo de atenci\'on. Para el token
\textit{comi\'o}, la atenci\'on calcula un peso $\alpha_i$ para cada
token de la secuencia. La representaci\'on final es una suma ponderada
de los valores, donde los pesos reflejan la relevancia sem\'antica.}
\label{fig:atencion_intuicion}
\end{figure}


Tres propiedades de este mecanismo lo distinguen radicalmente de los modelos
clásicos:

\begin{enumerate}
    \item \textbf{Global:} todos los tokens pueden relacionarse con todos,
    sin restricción de distancia.
    \item \textbf{Dinámico:} los pesos $\alpha_{ij}$ se calculan en función
    del contenido de la secuencia, no son parámetros fijos.
    \item \textbf{Paralelo:} el cómputo de todos los $\alpha_{ij}$ puede
    realizarse simultáneamente, sin procesamiento secuencial.
\end{enumerate}

%------------------------------------------------------------
\subsection{De la intuición a las matrices: Q, K y V}
\label{subsec:qkv}

La implementación práctica del mecanismo requiere operacionalizar la noción
de relevancia entre tokens. Vaswani et al.~\cite{vaswani2017attention}
proponen una analogía con los sistemas de recuperación de información: un
token emite una \textit{consulta} sobre qué necesita, el resto de los tokens
expone sus \textit{claves} (qué ofrecen), y la similitud entre consulta y
claves determina qué \textit{valores} se recuperan.

Concretamente, a partir de la representación de entrada $\bm{x}_i$ de cada
token se generan tres vectores mediante proyecciones lineales aprendidas:

\begin{align}
    \bm{q}_i &= W_Q\, \bm{x}_i \quad \in \mathbb{R}^{d_k}
    \quad \text{(\textit{query}: qué busca este token)} \\
    \bm{k}_j &= W_K\, \bm{x}_j \quad \in \mathbb{R}^{d_k}
    \quad \text{(\textit{key}: qué ofrece cada token)}   \\
    \bm{v}_j &= W_V\, \bm{x}_j \quad \in \mathbb{R}^{d_v}
    \quad \text{(\textit{value}: qué información aporta)}
\end{align}

\noindent donde $W_Q, W_K \in \mathbb{R}^{d_k \times d_{\text{model}}}$ y
$W_V \in \mathbb{R}^{d_v \times d_{\text{model}}}$ son matrices de peso
aprendidas. La relevancia entre los tokens $i$ y $j$ se mide mediante el
producto escalar $\bm{q}_i \cdot \bm{k}_j^\top$.

% ============================================================
% fig_qkv.tex: Mecanismo Q, K, V
% Uso: % ============================================================
% fig_qkv.tex: Mecanismo Q, K, V
% Uso: \input{figures/fig_qkv}
% ============================================================
\begin{figure}[H]
\centering
\begin{tikzpicture}[
    vec/.style={draw, rounded corners, minimum width=1.0cm,
                minimum height=0.55cm, font=\footnotesize, inner sep=3pt},
    q/.style={vec, fill=blue!15,  draw=blue!50},
    k/.style={vec, fill=red!12,   draw=red!45},
    v/.style={vec, fill=green!12, draw=green!50!black},
    op/.style={draw=gray!50, fill=gray!8, rounded corners,
               minimum width=1.1cm, minimum height=0.55cm,
               font=\footnotesize, align=center}
]

% ── Entrada ──────────────────────────────────────────────────
\node[vec, fill=gray!12, draw=gray!50, minimum width=2.0cm]
    (x) at (0, 0) {$\bm{x}_i \in \mathbb{R}^{d_{\text{model}}}$};

% ── Proyecciones ─────────────────────────────────────────────
\node[q, minimum width=1.5cm] (q) at (-2.8, -1.8)
    {$\bm{q}_i = W_Q \bm{x}_i$};
\node[k, minimum width=1.5cm] (k) at ( 0.0, -1.8)
    {$\bm{k}_j = W_K \bm{x}_j$};
\node[v, minimum width=1.5cm] (v) at ( 2.8, -1.8)
    {$\bm{v}_j = W_V \bm{x}_j$};

\draw[-{Stealth}, thick, blue!60]  (x.south) -- ++(-2.8,0) -- (q.north);
\draw[-{Stealth}, thick, red!50]   (x.south) -- (k.north);
\draw[-{Stealth}, thick, green!60!black] (x.south) -- ++(2.8,0) -- (v.north);

\node[blue!60,  font=\tiny] at (-1.8,-0.8) {$W_Q$};
\node[red!50,   font=\tiny] at ( 0.3,-0.9) {$W_K$};
\node[green!60!black, font=\tiny] at (1.8,-0.8) {$W_V$};

% ── Producto Q·K^T ───────────────────────────────────────────
\node[op] (dot) at (-1.4, -3.3) {$\bm{q}_i \cdot \bm{k}_j^{\top}$};
\draw[-{Stealth}, thick, gray] (q.south) -- ++(0,-0.7) -| (dot.north);
\draw[-{Stealth}, thick, gray] (k.south) -- ++(0,-0.7) -| (dot.north);

% ── Escalado ─────────────────────────────────────────────────
\node[op] (scale) at (-1.4, -4.5) {$/\sqrt{d_k}$};
\draw[-{Stealth}, thick] (dot) -- (scale);

% ── Softmax ──────────────────────────────────────────────────
\node[op, fill=orange!10, draw=orange!50] (softmax) at (-1.4, -5.7) {softmax};
\draw[-{Stealth}, thick] (scale) -- (softmax);
\node[font=\footnotesize\bfseries, orange!70] at (0.2,-5.7) {$\alpha_{ij}$};

% ── Suma ponderada ───────────────────────────────────────────
\node[op, fill=green!8, draw=green!50!black] (out) at (1.4, -7.0)
    {$\sum_j \alpha_{ij}\,\bm{v}_j$};
\draw[-{Stealth}, thick, orange!60]    (softmax.south) -- ++(0,-0.6) -| (out.north);
\draw[-{Stealth}, thick, green!60!black] (v.south)     -- ++(0,-4.5) -- (out.north);

% ── Salida ───────────────────────────────────────────────────
\node[vec, fill=purple!10, draw=purple!40, minimum width=1.8cm]
    (z) at (1.4, -8.2) {$\bm{z}_i \in \mathbb{R}^{d_v}$};
\draw[-{Stealth}, thick] (out) -- (z);
\node[font=\footnotesize, gray, align=center] at (3.4,-8.2)
    {representaci\'on\\contextualizada};

% ── Fórmula ──────────────────────────────────────────────────
\node[draw=blue!20, fill=blue!5, rounded corners, font=\small,
      inner sep=7pt, align=center] at (5.5, -4.5)
    {$\text{Attention}(Q,K,V)$\\[4pt]
     $=\,\text{softmax}\!\left(\dfrac{QK^{\top}}{\sqrt{d_k}}\right)V$};

\end{tikzpicture}
\caption{Mecanismo de atenci\'on de producto escalar escalado. A partir de
$\bm{x}_i$ se generan tres proyecciones: consulta $\bm{q}_i$, clave
$\bm{k}_j$ y valor $\bm{v}_j$. El producto $\bm{q}_i \cdot \bm{k}_j^\top$
mide la afinidad, se escala por $\sqrt{d_k}$ y se normaliza con softmax
para obtener los pesos $\alpha_{ij}$. La salida $\bm{z}_i$ es la suma
ponderada de los valores~\cite{vaswani2017attention}.}
\label{fig:qkv}
\end{figure}

% ============================================================
\begin{figure}[H]
\centering
\begin{tikzpicture}[
    vec/.style={draw, rounded corners, minimum width=1.0cm,
                minimum height=0.55cm, font=\footnotesize, inner sep=3pt},
    q/.style={vec, fill=blue!15,  draw=blue!50},
    k/.style={vec, fill=red!12,   draw=red!45},
    v/.style={vec, fill=green!12, draw=green!50!black},
    op/.style={draw=gray!50, fill=gray!8, rounded corners,
               minimum width=1.1cm, minimum height=0.55cm,
               font=\footnotesize, align=center}
]

% ── Entrada ──────────────────────────────────────────────────
\node[vec, fill=gray!12, draw=gray!50, minimum width=2.0cm]
    (x) at (0, 0) {$\bm{x}_i \in \mathbb{R}^{d_{\text{model}}}$};

% ── Proyecciones ─────────────────────────────────────────────
\node[q, minimum width=1.5cm] (q) at (-2.8, -1.8)
    {$\bm{q}_i = W_Q \bm{x}_i$};
\node[k, minimum width=1.5cm] (k) at ( 0.0, -1.8)
    {$\bm{k}_j = W_K \bm{x}_j$};
\node[v, minimum width=1.5cm] (v) at ( 2.8, -1.8)
    {$\bm{v}_j = W_V \bm{x}_j$};

\draw[-{Stealth}, thick, blue!60]  (x.south) -- ++(-2.8,0) -- (q.north);
\draw[-{Stealth}, thick, red!50]   (x.south) -- (k.north);
\draw[-{Stealth}, thick, green!60!black] (x.south) -- ++(2.8,0) -- (v.north);

\node[blue!60,  font=\tiny] at (-1.8,-0.8) {$W_Q$};
\node[red!50,   font=\tiny] at ( 0.3,-0.9) {$W_K$};
\node[green!60!black, font=\tiny] at (1.8,-0.8) {$W_V$};

% ── Producto Q·K^T ───────────────────────────────────────────
\node[op] (dot) at (-1.4, -3.3) {$\bm{q}_i \cdot \bm{k}_j^{\top}$};
\draw[-{Stealth}, thick, gray] (q.south) -- ++(0,-0.7) -| (dot.north);
\draw[-{Stealth}, thick, gray] (k.south) -- ++(0,-0.7) -| (dot.north);

% ── Escalado ─────────────────────────────────────────────────
\node[op] (scale) at (-1.4, -4.5) {$/\sqrt{d_k}$};
\draw[-{Stealth}, thick] (dot) -- (scale);

% ── Softmax ──────────────────────────────────────────────────
\node[op, fill=orange!10, draw=orange!50] (softmax) at (-1.4, -5.7) {softmax};
\draw[-{Stealth}, thick] (scale) -- (softmax);
\node[font=\footnotesize\bfseries, orange!70] at (0.2,-5.7) {$\alpha_{ij}$};

% ── Suma ponderada ───────────────────────────────────────────
\node[op, fill=green!8, draw=green!50!black] (out) at (1.4, -7.0)
    {$\sum_j \alpha_{ij}\,\bm{v}_j$};
\draw[-{Stealth}, thick, orange!60]    (softmax.south) -- ++(0,-0.6) -| (out.north);
\draw[-{Stealth}, thick, green!60!black] (v.south)     -- ++(0,-4.5) -- (out.north);

% ── Salida ───────────────────────────────────────────────────
\node[vec, fill=purple!10, draw=purple!40, minimum width=1.8cm]
    (z) at (1.4, -8.2) {$\bm{z}_i \in \mathbb{R}^{d_v}$};
\draw[-{Stealth}, thick] (out) -- (z);
\node[font=\footnotesize, gray, align=center] at (3.4,-8.2)
    {representaci\'on\\contextualizada};

% ── Fórmula ──────────────────────────────────────────────────
\node[draw=blue!20, fill=blue!5, rounded corners, font=\small,
      inner sep=7pt, align=center] at (5.5, -4.5)
    {$\text{Attention}(Q,K,V)$\\[4pt]
     $=\,\text{softmax}\!\left(\dfrac{QK^{\top}}{\sqrt{d_k}}\right)V$};

\end{tikzpicture}
\caption{Mecanismo de atenci\'on de producto escalar escalado. A partir de
$\bm{x}_i$ se generan tres proyecciones: consulta $\bm{q}_i$, clave
$\bm{k}_j$ y valor $\bm{v}_j$. El producto $\bm{q}_i \cdot \bm{k}_j^\top$
mide la afinidad, se escala por $\sqrt{d_k}$ y se normaliza con softmax
para obtener los pesos $\alpha_{ij}$. La salida $\bm{z}_i$ es la suma
ponderada de los valores~\cite{vaswani2017attention}.}
\label{fig:qkv}
\end{figure}


%------------------------------------------------------------
\subsection{Atención de producto escalar escalado}
\label{subsec:scaled_dot_product}

Operando sobre toda la secuencia de forma matricial, las consultas, claves
y valores se apilan en matrices $Q \in \mathbb{R}^{n \times d_k}$,
$K \in \mathbb{R}^{n \times d_k}$ y $V \in \mathbb{R}^{n \times d_v}$.
La atención se computa entonces como:

\begin{equation}
    \text{Attention}(Q, K, V) =
    \text{softmax}\!\left(\frac{QK^{\top}}{\sqrt{d_k}}\right) V
    \label{eq:attention}
\end{equation}

Cada componente cumple un rol preciso:

\begin{itemize}
    \item $QK^{\top}$: matriz de \textit{scores} de afinidad entre todos los
    pares de tokens. El elemento $(i,j)$ mide cuánto se parece la consulta
    del token $i$ a la clave del token $j$.

    \item $1/\sqrt{d_k}$: factor de escala que estabiliza el entrenamiento.
    Para dimensiones $d_k$ grandes, los productos escalares crecen en
    magnitud y empujan a la función softmax hacia regiones de gradiente
    muy pequeño; dividir por $\sqrt{d_k}$ corrige este efecto~\cite{vaswani2017attention}.

    \item $\text{softmax}(\cdot)$: normaliza los scores por fila, produciendo
    pesos $\alpha_{ij} \in (0,1)$ con $\sum_j \alpha_{ij} = 1$, interpretables
    como una distribución de probabilidad sobre la secuencia.

    \item $V$: los pesos se aplican sobre los valores, produciendo para cada
    token una representación contextualizada que integra información de toda
    la secuencia ponderada por relevancia.
\end{itemize}

%------------------------------------------------------------
\subsection{Ejemplo numérico paso a paso}
\label{subsec:atencion_ejemplo}

Para concretar el mecanismo, se presenta un ejemplo con la secuencia
\textit{gato–comió–ratón} ($n = 3$ tokens, $d_k = d_v = 2$). Los valores
numéricos son ilustrativos y están elegidos para producir un resultado
interpretable.

% ============================================================
% fig_atencion_ejemplo.tex: Ejemplo numérico Q, K, V (3 tokens)
% Uso: % ============================================================
% fig_atencion_ejemplo.tex: Ejemplo numérico Q, K, V (3 tokens)
% Uso: \input{figures/fig_atencion_ejemplo}
% ============================================================
\begin{figure}[H]
\centering
\begin{tikzpicture}[
    cell/.style={draw=gray!40, minimum width=1.1cm, minimum height=0.60cm,
                 font=\footnotesize, inner sep=2pt},
    hdr/.style={font=\footnotesize\bfseries},
    qc/.style={fill=blue!12},
    kc/.style={fill=red!10},
    vc/.style={fill=green!10},
    sc/.style={fill=orange!12},
    ac/.style={fill=purple!10},
    zc/.style={fill=teal!10},
    stepnode/.style={font=\small\bfseries, blue!70}
]

\node[font=\small\bfseries, gray] at (5.0, 8.5)
    {Secuencia: \textit{gato} $(t_1)$ \quad \textit{comi\'o} $(t_2)$ \quad
     \textit{rat\'on} $(t_3)$ \quad $|$ \quad $d_k = d_v = 2$};

% ── PASO 1: Matrices Q, K, V ─────────────────────────────────
\node[stepnode] at (-0.5, 7.7) {Paso 1};
\foreach \lbl/\x in {Q/1.5, K/4.0, V/6.5}{
    \node[hdr] at (\x, 7.7) {$\lbl$};
}
\foreach \lbl/\y in {$t_1$/7.1, $t_2$/6.45, $t_3$/5.80}{
    \node[font=\footnotesize] at (0.3, \y) {\lbl};
}
% Q
\node[cell,qc] at (1.1,7.10){1.0}; \node[cell,qc] at (2.0,7.10){0.0};
\node[cell,qc] at (1.1,6.45){0.0}; \node[cell,qc] at (2.0,6.45){1.0};
\node[cell,qc] at (1.1,5.80){1.0}; \node[cell,qc] at (2.0,5.80){1.0};
% K
\node[cell,kc] at (3.6,7.10){1.0}; \node[cell,kc] at (4.5,7.10){0.0};
\node[cell,kc] at (3.6,6.45){1.0}; \node[cell,kc] at (4.5,6.45){1.0};
\node[cell,kc] at (3.6,5.80){0.0}; \node[cell,kc] at (4.5,5.80){1.0};
% V
\node[cell,vc] at (6.1,7.10){1.0}; \node[cell,vc] at (7.0,7.10){0.0};
\node[cell,vc] at (6.1,6.45){0.5}; \node[cell,vc] at (7.0,6.45){0.5};
\node[cell,vc] at (6.1,5.80){0.0}; \node[cell,vc] at (7.0,5.80){1.0};

% ── PASO 2: Scores ───────────────────────────────────────────
\node[stepnode] at (-0.5, 4.85) {Paso 2};
\node[hdr, align=center] at (1.8, 4.85) {$S = QK^{\top}/\sqrt{2}$};
\foreach \lbl/\x in {$t_1$/1.1, $t_2$/2.2, $t_3$/3.3}{
    \node[font=\footnotesize] at (\x, 4.35) {\lbl};
}
\foreach \lbl/\y in {$t_1$/3.85, $t_2$/3.20, $t_3$/2.55}{
    \node[font=\footnotesize] at (0.3, \y) {\lbl};
}
\node[cell,sc] at (1.1,3.85){0.71}; \node[cell,sc] at (2.2,3.85){0.71}; \node[cell,sc] at (3.3,3.85){0.00};
\node[cell,sc] at (1.1,3.20){0.00}; \node[cell,sc] at (2.2,3.20){0.71}; \node[cell,sc] at (3.3,3.20){0.71};
\node[cell,sc] at (1.1,2.55){0.71}; \node[cell,sc] at (2.2,2.55){1.41}; \node[cell,sc] at (3.3,2.55){0.71};

% ── Flecha softmax ───────────────────────────────────────────
\draw[-{Stealth}, thick, gray!60] (3.8, 3.2) -- (4.6, 3.2);
\node[font=\tiny, gray, align=center] at (4.2, 3.5) {softmax\\por fila};

% ── PASO 3: Pesos α ──────────────────────────────────────────
\node[stepnode] at (5.2, 4.85) {Paso 3};
\node[hdr, align=center] at (6.8, 4.85) {$A = \text{softmax}(S)$};
\foreach \lbl/\x in {$t_1$/5.6, $t_2$/6.7, $t_3$/7.8}{
    \node[font=\footnotesize] at (\x, 4.35) {\lbl};
}
\foreach \lbl/\y in {$t_1$/3.85, $t_2$/3.20, $t_3$/2.55}{
    \node[font=\footnotesize] at (4.9, \y) {\lbl};
}
\node[cell,ac] at (5.6,3.85){0.40}; \node[cell,ac] at (6.7,3.85){0.40}; \node[cell,ac] at (7.8,3.85){0.20};
\node[cell,ac] at (5.6,3.20){0.20}; \node[cell,ac] at (6.7,3.20){0.40}; \node[cell,ac] at (7.8,3.20){0.40};
\node[cell,ac] at (5.6,2.55){0.25}; \node[cell,ac] at (6.7,2.55){0.50}; \node[cell,ac] at (7.8,2.55){0.25};

% ── PASO 4: Z = A·V ──────────────────────────────────────────
\node[stepnode] at (3.0, 1.85) {Paso 4};
\node[hdr, align=center] at (5.2, 1.85) {$Z = A \cdot V$};
\foreach \lbl/\x in {dim1/4.6, dim2/5.7}{
    \node[font=\footnotesize] at (\x, 1.35) {\lbl};
}
\foreach \lbl/\y in {$z_1$/0.90, $z_2$/0.25, $z_3$/-0.40}{
    \node[font=\footnotesize] at (3.8, \y) {\lbl};
}
\node[cell,zc] at (4.6, 0.90){0.60}; \node[cell,zc] at (5.7, 0.90){0.40};
\node[cell,zc] at (4.6, 0.25){0.40}; \node[cell,zc] at (5.7, 0.25){0.60};
\node[cell,zc] at (4.6,-0.40){0.50}; \node[cell,zc] at (5.7,-0.40){0.50};

\draw[-{Stealth}, thick, gray!60]
    (7.3, 2.55) -- ++(0,-0.4) -- ++(-1.1,0) -- ++(0,-0.4);
\node[font=\tiny, gray] at (8.3, 2.1) {$Z = A \cdot V$};

\node[draw=gray!30, fill=gray!5, rounded corners, font=\tiny,
      inner sep=5pt, align=left] at (8.8, 0.3)
    {Fila $i$ de $Z$: representaci\'on de $t_i$\\
     contextualizada por toda la secuencia.};

\end{tikzpicture}
\caption{Ejemplo num\'erico de atenci\'on escalada con 3 tokens ($d_k=d_v=2$).
\textbf{Paso 1:} matrices $Q$, $K$, $V$.
\textbf{Paso 2:} scores $S=QK^\top/\sqrt{d_k}$.
\textbf{Paso 3:} pesos $\alpha_{ij}$ tras softmax por fila.
\textbf{Paso 4:} salida $Z=AV$. Valores ilustrativos.}
\label{fig:atencion_ejemplo}
\end{figure}

% ============================================================
\begin{figure}[H]
\centering
\begin{tikzpicture}[
    cell/.style={draw=gray!40, minimum width=1.1cm, minimum height=0.60cm,
                 font=\footnotesize, inner sep=2pt},
    hdr/.style={font=\footnotesize\bfseries},
    qc/.style={fill=blue!12},
    kc/.style={fill=red!10},
    vc/.style={fill=green!10},
    sc/.style={fill=orange!12},
    ac/.style={fill=purple!10},
    zc/.style={fill=teal!10},
    stepnode/.style={font=\small\bfseries, blue!70}
]

\node[font=\small\bfseries, gray] at (5.0, 8.5)
    {Secuencia: \textit{gato} $(t_1)$ \quad \textit{comi\'o} $(t_2)$ \quad
     \textit{rat\'on} $(t_3)$ \quad $|$ \quad $d_k = d_v = 2$};

% ── PASO 1: Matrices Q, K, V ─────────────────────────────────
\node[stepnode] at (-0.5, 7.7) {Paso 1};
\foreach \lbl/\x in {Q/1.5, K/4.0, V/6.5}{
    \node[hdr] at (\x, 7.7) {$\lbl$};
}
\foreach \lbl/\y in {$t_1$/7.1, $t_2$/6.45, $t_3$/5.80}{
    \node[font=\footnotesize] at (0.3, \y) {\lbl};
}
% Q
\node[cell,qc] at (1.1,7.10){1.0}; \node[cell,qc] at (2.0,7.10){0.0};
\node[cell,qc] at (1.1,6.45){0.0}; \node[cell,qc] at (2.0,6.45){1.0};
\node[cell,qc] at (1.1,5.80){1.0}; \node[cell,qc] at (2.0,5.80){1.0};
% K
\node[cell,kc] at (3.6,7.10){1.0}; \node[cell,kc] at (4.5,7.10){0.0};
\node[cell,kc] at (3.6,6.45){1.0}; \node[cell,kc] at (4.5,6.45){1.0};
\node[cell,kc] at (3.6,5.80){0.0}; \node[cell,kc] at (4.5,5.80){1.0};
% V
\node[cell,vc] at (6.1,7.10){1.0}; \node[cell,vc] at (7.0,7.10){0.0};
\node[cell,vc] at (6.1,6.45){0.5}; \node[cell,vc] at (7.0,6.45){0.5};
\node[cell,vc] at (6.1,5.80){0.0}; \node[cell,vc] at (7.0,5.80){1.0};

% ── PASO 2: Scores ───────────────────────────────────────────
\node[stepnode] at (-0.5, 4.85) {Paso 2};
\node[hdr, align=center] at (1.8, 4.85) {$S = QK^{\top}/\sqrt{2}$};
\foreach \lbl/\x in {$t_1$/1.1, $t_2$/2.2, $t_3$/3.3}{
    \node[font=\footnotesize] at (\x, 4.35) {\lbl};
}
\foreach \lbl/\y in {$t_1$/3.85, $t_2$/3.20, $t_3$/2.55}{
    \node[font=\footnotesize] at (0.3, \y) {\lbl};
}
\node[cell,sc] at (1.1,3.85){0.71}; \node[cell,sc] at (2.2,3.85){0.71}; \node[cell,sc] at (3.3,3.85){0.00};
\node[cell,sc] at (1.1,3.20){0.00}; \node[cell,sc] at (2.2,3.20){0.71}; \node[cell,sc] at (3.3,3.20){0.71};
\node[cell,sc] at (1.1,2.55){0.71}; \node[cell,sc] at (2.2,2.55){1.41}; \node[cell,sc] at (3.3,2.55){0.71};

% ── Flecha softmax ───────────────────────────────────────────
\draw[-{Stealth}, thick, gray!60] (3.8, 3.2) -- (4.6, 3.2);
\node[font=\tiny, gray, align=center] at (4.2, 3.5) {softmax\\por fila};

% ── PASO 3: Pesos α ──────────────────────────────────────────
\node[stepnode] at (5.2, 4.85) {Paso 3};
\node[hdr, align=center] at (6.8, 4.85) {$A = \text{softmax}(S)$};
\foreach \lbl/\x in {$t_1$/5.6, $t_2$/6.7, $t_3$/7.8}{
    \node[font=\footnotesize] at (\x, 4.35) {\lbl};
}
\foreach \lbl/\y in {$t_1$/3.85, $t_2$/3.20, $t_3$/2.55}{
    \node[font=\footnotesize] at (4.9, \y) {\lbl};
}
\node[cell,ac] at (5.6,3.85){0.40}; \node[cell,ac] at (6.7,3.85){0.40}; \node[cell,ac] at (7.8,3.85){0.20};
\node[cell,ac] at (5.6,3.20){0.20}; \node[cell,ac] at (6.7,3.20){0.40}; \node[cell,ac] at (7.8,3.20){0.40};
\node[cell,ac] at (5.6,2.55){0.25}; \node[cell,ac] at (6.7,2.55){0.50}; \node[cell,ac] at (7.8,2.55){0.25};

% ── PASO 4: Z = A·V ──────────────────────────────────────────
\node[stepnode] at (3.0, 1.85) {Paso 4};
\node[hdr, align=center] at (5.2, 1.85) {$Z = A \cdot V$};
\foreach \lbl/\x in {dim1/4.6, dim2/5.7}{
    \node[font=\footnotesize] at (\x, 1.35) {\lbl};
}
\foreach \lbl/\y in {$z_1$/0.90, $z_2$/0.25, $z_3$/-0.40}{
    \node[font=\footnotesize] at (3.8, \y) {\lbl};
}
\node[cell,zc] at (4.6, 0.90){0.60}; \node[cell,zc] at (5.7, 0.90){0.40};
\node[cell,zc] at (4.6, 0.25){0.40}; \node[cell,zc] at (5.7, 0.25){0.60};
\node[cell,zc] at (4.6,-0.40){0.50}; \node[cell,zc] at (5.7,-0.40){0.50};

\draw[-{Stealth}, thick, gray!60]
    (7.3, 2.55) -- ++(0,-0.4) -- ++(-1.1,0) -- ++(0,-0.4);
\node[font=\tiny, gray] at (8.3, 2.1) {$Z = A \cdot V$};

\node[draw=gray!30, fill=gray!5, rounded corners, font=\tiny,
      inner sep=5pt, align=left] at (8.8, 0.3)
    {Fila $i$ de $Z$: representaci\'on de $t_i$\\
     contextualizada por toda la secuencia.};

\end{tikzpicture}
\caption{Ejemplo num\'erico de atenci\'on escalada con 3 tokens ($d_k=d_v=2$).
\textbf{Paso 1:} matrices $Q$, $K$, $V$.
\textbf{Paso 2:} scores $S=QK^\top/\sqrt{d_k}$.
\textbf{Paso 3:} pesos $\alpha_{ij}$ tras softmax por fila.
\textbf{Paso 4:} salida $Z=AV$. Valores ilustrativos.}
\label{fig:atencion_ejemplo}
\end{figure}


La fila $i$ de la salida $Z$ contiene la representación contextualizada
del token $i$. Por ejemplo, la fila de \textit{comió} ($t_2$) de la matriz de
pesos es $(0.20,\, 0.40,\, 0.40)$: el verbo reparte su atención por igual
entre sí mismo y \textit{ratón}, y le presta menos a \textit{gato}. Su
representación $z_2 = (0.40,\, 0.60)$ es entonces una mezcla de los valores
de los tres tokens en esas proporciones. Este resultado emerge del producto
escalar entre vectores, sin reglas explícitas. Los valores de $Q$, $K$ y $V$
son ilustrativos: en un modelo real los aprende el entrenamiento.

%------------------------------------------------------------
\subsection{Atención multi-cabeza}
\label{subsec:multihead}

Un único mecanismo de atención puede capturar un tipo de relación a la vez.
Sin embargo, una oración contiene simultáneamente relaciones sintácticas,
semánticas, de correferencia, posicionales, entre otras. La
\textit{atención multi-cabeza} (\textit{Multi-Head Attention}) resuelve
esto ejecutando $h$ mecanismos de atención en paralelo, cada uno con sus
propias matrices de proyección:

\begin{equation}
    \text{head}_l = \text{Attention}(Q W_Q^l,\; K W_K^l,\; V W_V^l)
    \qquad l = 1, \ldots, h
\end{equation}

Las salidas de cada cabeza se concatenan y se proyectan con una matriz
$W^O \in \mathbb{R}^{h d_v \times d_{\text{model}}}$:

\begin{equation}
    \text{MultiHead}(Q, K, V) =
    \text{Concat}(\text{head}_1, \ldots, \text{head}_h)\, W^O
\end{equation}

Para mantener el costo computacional constante, cada cabeza opera en un
subespacio de dimensión reducida: $d_k = d_v = d_{\text{model}} / h$.

% ============================================================
% fig_multihead.tex: Multi-Head Attention
% Uso: % ============================================================
% fig_multihead.tex: Multi-Head Attention
% Uso: \input{figures/fig_multihead}
% ============================================================
\begin{figure}[H]
\centering
\begin{tikzpicture}[
    head/.style={draw, rounded corners, minimum width=1.5cm,
                 minimum height=0.7cm, font=\footnotesize\bfseries,
                 align=center},
    h1/.style={head, fill=blue!12,   draw=blue!50},
    h2/.style={head, fill=red!10,    draw=red!45},
    h3/.style={head, fill=green!10,  draw=green!50!black},
    h4/.style={head, fill=orange!10, draw=orange!50},
    box/.style={draw=gray!50, fill=gray!6, rounded corners,
                minimum width=5.5cm, minimum height=0.7cm,
                font=\small, align=center}
]

% ── Entrada ──────────────────────────────────────────────────
\node[box] (inp) at (3.5, 6.5)
    {Entrada $X \in \mathbb{R}^{n \times d_{\text{model}}}$};

% ── 4 cabezas ────────────────────────────────────────────────
\node[h1] (h1) at (0.7,  5.0) {Head 1\\$W_Q^1,W_K^1,W_V^1$};
\node[h2] (h2) at (2.5,  5.0) {Head 2\\$W_Q^2,W_K^2,W_V^2$};
\node[h3] (h3) at (4.3,  5.0) {Head 3\\$W_Q^3,W_K^3,W_V^3$};
\node[h4] (h4) at (6.1,  5.0) {Head 4\\$W_Q^4,W_K^4,W_V^4$};

\draw[-{Stealth}, thick, blue!50]        (inp.south) -| (h1.north);
\draw[-{Stealth}, thick, red!45]         (inp.south) -| (h2.north);
\draw[-{Stealth}, thick, green!50!black] (inp.south) -| (h3.north);
\draw[-{Stealth}, thick, orange!50]      (inp.south) -| (h4.north);

\node[font=\footnotesize, gray] at (3.5, 5.85)
    {proyecciones independientes ($d_k = d_{\text{model}}/h$)};

% ── Salidas ──────────────────────────────────────────────────
\node[h1, minimum height=0.55cm] (z1) at (0.7,  3.7) {$Z_1$};
\node[h2, minimum height=0.55cm] (z2) at (2.5,  3.7) {$Z_2$};
\node[h3, minimum height=0.55cm] (z3) at (4.3,  3.7) {$Z_3$};
\node[h4, minimum height=0.55cm] (z4) at (6.1,  3.7) {$Z_4$};

\draw[-{Stealth}, thick] (h1)--(z1);
\draw[-{Stealth}, thick] (h2)--(z2);
\draw[-{Stealth}, thick] (h3)--(z3);
\draw[-{Stealth}, thick] (h4)--(z4);

% ── Concatenación ────────────────────────────────────────────
\node[box, fill=yellow!8, draw=yellow!60!black] (concat) at (3.5, 2.6)
    {Concat$(Z_1, Z_2, Z_3, Z_4)$};
\draw[-{Stealth}, thick] (z1.south) |- (concat.west);
\draw[-{Stealth}, thick] (z2.south) -- (concat.north);
\draw[-{Stealth}, thick] (z3.south) -- (concat.north);
\draw[-{Stealth}, thick] (z4.south) |- (concat.east);

% ── Proyección final ─────────────────────────────────────────
\node[box, fill=purple!8, draw=purple!40] (proj) at (3.5, 1.5)
    {Proyecci\'on lineal $W^O$};
\draw[-{Stealth}, thick] (concat)--(proj);

% ── Salida ───────────────────────────────────────────────────
\node[box] (out) at (3.5, 0.5)
    {Salida $\in \mathbb{R}^{n \times d_{\text{model}}}$};
\draw[-{Stealth}, thick] (proj)--(out);

% ── Anotación ────────────────────────────────────────────────
\node[draw=blue!15, fill=blue!4, rounded corners, font=\footnotesize,
      inner sep=6pt, align=left] at (9.6, 3.8)
    {Cada cabeza aprende\\un tipo de relaci\'on\\distinto en paralelo:\\[4pt]
     \textcolor{blue!70}{Head 1:} sint\'actica\\
     \textcolor{red!60}{Head 2:} sem\'antica\\
     \textcolor{green!60!black}{Head 3:} correferencia\\
     \textcolor{orange!70}{Head 4:} posici\'on};

\node[draw=gray!30, fill=gray!5, rounded corners, font=\small,
      inner sep=6pt, align=center] at (9.6, 1.2)
    {$\text{MHA}(Q,K,V) =$\\[4pt]
     $\text{Concat}(\text{head}_1,\ldots,\text{head}_h)W^O$};

\end{tikzpicture}
\caption{Atenci\'on multi-cabeza. La entrada se proyecta $h$ veces con
matrices independientes; cada cabeza opera en un subespacio de dimensi\'on
$d_k = d_{\text{model}}/h$. Las salidas se concatenan y se proyectan con
$W^O$~\cite{vaswani2017attention}.}
\label{fig:multihead}
\end{figure}

% ============================================================
\begin{figure}[H]
\centering
\begin{tikzpicture}[
    head/.style={draw, rounded corners, minimum width=1.5cm,
                 minimum height=0.7cm, font=\footnotesize\bfseries,
                 align=center},
    h1/.style={head, fill=blue!12,   draw=blue!50},
    h2/.style={head, fill=red!10,    draw=red!45},
    h3/.style={head, fill=green!10,  draw=green!50!black},
    h4/.style={head, fill=orange!10, draw=orange!50},
    box/.style={draw=gray!50, fill=gray!6, rounded corners,
                minimum width=5.5cm, minimum height=0.7cm,
                font=\small, align=center}
]

% ── Entrada ──────────────────────────────────────────────────
\node[box] (inp) at (3.5, 6.5)
    {Entrada $X \in \mathbb{R}^{n \times d_{\text{model}}}$};

% ── 4 cabezas ────────────────────────────────────────────────
\node[h1] (h1) at (0.7,  5.0) {Head 1\\$W_Q^1,W_K^1,W_V^1$};
\node[h2] (h2) at (2.5,  5.0) {Head 2\\$W_Q^2,W_K^2,W_V^2$};
\node[h3] (h3) at (4.3,  5.0) {Head 3\\$W_Q^3,W_K^3,W_V^3$};
\node[h4] (h4) at (6.1,  5.0) {Head 4\\$W_Q^4,W_K^4,W_V^4$};

\draw[-{Stealth}, thick, blue!50]        (inp.south) -| (h1.north);
\draw[-{Stealth}, thick, red!45]         (inp.south) -| (h2.north);
\draw[-{Stealth}, thick, green!50!black] (inp.south) -| (h3.north);
\draw[-{Stealth}, thick, orange!50]      (inp.south) -| (h4.north);

\node[font=\footnotesize, gray] at (3.5, 5.85)
    {proyecciones independientes ($d_k = d_{\text{model}}/h$)};

% ── Salidas ──────────────────────────────────────────────────
\node[h1, minimum height=0.55cm] (z1) at (0.7,  3.7) {$Z_1$};
\node[h2, minimum height=0.55cm] (z2) at (2.5,  3.7) {$Z_2$};
\node[h3, minimum height=0.55cm] (z3) at (4.3,  3.7) {$Z_3$};
\node[h4, minimum height=0.55cm] (z4) at (6.1,  3.7) {$Z_4$};

\draw[-{Stealth}, thick] (h1)--(z1);
\draw[-{Stealth}, thick] (h2)--(z2);
\draw[-{Stealth}, thick] (h3)--(z3);
\draw[-{Stealth}, thick] (h4)--(z4);

% ── Concatenación ────────────────────────────────────────────
\node[box, fill=yellow!8, draw=yellow!60!black] (concat) at (3.5, 2.6)
    {Concat$(Z_1, Z_2, Z_3, Z_4)$};
\draw[-{Stealth}, thick] (z1.south) |- (concat.west);
\draw[-{Stealth}, thick] (z2.south) -- (concat.north);
\draw[-{Stealth}, thick] (z3.south) -- (concat.north);
\draw[-{Stealth}, thick] (z4.south) |- (concat.east);

% ── Proyección final ─────────────────────────────────────────
\node[box, fill=purple!8, draw=purple!40] (proj) at (3.5, 1.5)
    {Proyecci\'on lineal $W^O$};
\draw[-{Stealth}, thick] (concat)--(proj);

% ── Salida ───────────────────────────────────────────────────
\node[box] (out) at (3.5, 0.5)
    {Salida $\in \mathbb{R}^{n \times d_{\text{model}}}$};
\draw[-{Stealth}, thick] (proj)--(out);

% ── Anotación ────────────────────────────────────────────────
\node[draw=blue!15, fill=blue!4, rounded corners, font=\footnotesize,
      inner sep=6pt, align=left] at (9.6, 3.8)
    {Cada cabeza aprende\\un tipo de relaci\'on\\distinto en paralelo:\\[4pt]
     \textcolor{blue!70}{Head 1:} sint\'actica\\
     \textcolor{red!60}{Head 2:} sem\'antica\\
     \textcolor{green!60!black}{Head 3:} correferencia\\
     \textcolor{orange!70}{Head 4:} posici\'on};

\node[draw=gray!30, fill=gray!5, rounded corners, font=\small,
      inner sep=6pt, align=center] at (9.6, 1.2)
    {$\text{MHA}(Q,K,V) =$\\[4pt]
     $\text{Concat}(\text{head}_1,\ldots,\text{head}_h)W^O$};

\end{tikzpicture}
\caption{Atenci\'on multi-cabeza. La entrada se proyecta $h$ veces con
matrices independientes; cada cabeza opera en un subespacio de dimensi\'on
$d_k = d_{\text{model}}/h$. Las salidas se concatenan y se proyectan con
$W^O$~\cite{vaswani2017attention}.}
\label{fig:multihead}
\end{figure}


%------------------------------------------------------------
\subsection{Auto-atención y máscara causal}
\label{subsec:autoatencion}

Cuando $Q$, $K$ y $V$ se derivan de la misma secuencia de entrada, el
mecanismo se denomina \textit{auto-atención} (\textit{self-attention}):
cada token se relaciona con todos los demás tokens de la misma secuencia.
Es el mecanismo central del encoder de un Transformer.

Para tareas de generación autoregresiva, donde el modelo predice el
siguiente token sin acceder a tokens futuros, se aplica una
\textit{máscara causal}: una matriz que fuerza a $-\infty$ los scores
de los pares $(i, j)$ con $j > i$ antes del softmax:

\begin{equation}
    \text{Attention}_{\text{causal}}(Q,K,V) =
    \text{softmax}\!\left(\frac{QK^{\top} + M}{\sqrt{d_k}}\right) V
\end{equation}

\noindent donde $M_{ij} = -\infty$ si $j > i$ y $0$ en caso contrario.
Esto garantiza que el token $i$ solo vea los tokens anteriores, manteniendo
la propiedad autoregresiva sin sacrificar el paralelismo durante el
entrenamiento.

%------------------------------------------------------------
\subsection{Conexión con los modelos clásicos}
\label{subsec:atencion_conexion}

El mecanismo de atención generaliza varios conceptos ya estudiados:

\begin{itemize}
    \item \textbf{Hopfield}: la suma ponderada $\sum_j \alpha_{ij} \bm{v}_j$
    es formalmente análoga a la recuperación de memoria asociativa. Los
    pesos $\alpha_{ij}$ reemplazan la dinámica de minimización de energía
    por un cómputo diferenciable y paralelizable. Ramsauer et
    al.~\cite{ramsauer2020hopfield} demuestran esta equivalencia formalmente.

    \item \textbf{SOM}: los embeddings sobre los que opera la atención son
    espacios de representación distribuida donde la proximidad geométrica
    refleja similitud semántica, análogos a los mapas de Kohonen.

    \item \textbf{MLP}: las proyecciones $W_Q$, $W_K$, $W_V$ son
    transformaciones lineales aprendidas (el mismo bloque fundamental
    de los MLPs), pero utilizadas aquí para construir representaciones
    relacionales en lugar de clasificar directamente.
\end{itemize}

La atención no descarta lo construido por los modelos clásicos: lo integra
y generaliza para resolver las limitaciones identificadas en la
Sección~\ref{sec:problema_nlp}. La siguiente sección describe cómo este
mecanismo se ensambla dentro de la arquitectura completa del Transformer.

\section{La arquitectura Transformer}
\label{sec:arquitectura_transformer}

Las secciones anteriores construyeron los conceptos fundamentales: las
limitaciones de los modelos clásicos, el problema del lenguaje, y el
mecanismo de atención como solución. Esta sección ensambla esos conceptos
en la arquitectura completa del Transformer, describiendo cada bloque en
el orden en que los datos los atraviesan.

%------------------------------------------------------------
\subsection{Visión general del pipeline}
\label{subsec:pipeline}

Antes de entrar en los detalles de cada componente, conviene tener una
imagen clara del flujo completo de información. Dada una secuencia de
tokens de entrada, el Transformer realiza los siguientes pasos en orden:

\begin{enumerate}
    \item Cada token se convierte en un vector denso mediante una capa de
    \textbf{embeddings}.
    \item Se suma un vector de \textbf{codificación posicional} para
    inyectar información de orden.
    \item La secuencia pasa por $N$ \textbf{bloques encoder} apilados,
    cada uno con atención multi-cabeza y una red feed-forward.
    \item (En modelos seq2seq) Una secuencia de salida parcial se procesa
    en paralelo por $N$ \textbf{bloques decoder}, que combinan
    auto-atención enmascarada con atención cruzada sobre el encoder.
    \item Una proyección lineal y un softmax producen una distribución
    de probabilidad sobre el vocabulario para el siguiente token.
\end{enumerate}

%------------------------------------------------------------
\subsection{Codificación posicional}
\label{subsec:positional_encoding}

El mecanismo de atención, tal como se describió en la
Sección~\ref{sec:mecanismo_atencion}, es invariante a permutaciones: si se
reordena la secuencia, los pesos de atención cambian pero el modelo no
tiene forma de saber que el orden cambió, porque trata todos los tokens
de forma simétrica. Esto contradice una de las propiedades más importantes
del lenguaje: el orden importa.

La solución propuesta por Vaswani et al.~\cite{vaswani2017attention} es
sumar a cada embedding un vector de \textit{codificación posicional} $PE_t$
que depende únicamente de la posición $t$ del token en la secuencia:

\begin{equation}
    \bm{x}'_t = \bm{e}_t + PE_t
\end{equation}

La codificación utiliza funciones sinusoidales de distintas frecuencias
para cada dimensión del vector:

\begin{align}
    PE_{(t,\, 2i)}   &= \sin\!\left(\frac{t}{10000^{2i/d_{\text{model}}}}\right) \\[4pt]
    PE_{(t,\, 2i+1)} &= \cos\!\left(\frac{t}{10000^{2i/d_{\text{model}}}}\right)
\end{align}

\noindent donde $t$ es la posición del token y $i$ es el índice de la
dimensión. Las dimensiones bajas oscilan con frecuencia alta (capturan
relaciones locales) y las altas con frecuencia baja (capturan relaciones
globales), de forma análoga a cómo un sistema numérico posicional codifica
valores de distinta magnitud en cada dígito.

Esta codificación tiene tres propiedades convenientes: es determinista
(no requiere aprendizaje adicional), generaliza a secuencias más largas
que las vistas en entrenamiento, y permite al modelo inferir distancias
relativas entre posiciones a través de diferencias entre vectores.

% ============================================================
% fig_transformer_completo.tex: Arquitectura Transformer completa
% Uso: % ============================================================
% fig_transformer_completo.tex — Arquitectura Transformer completa
% Uso: \input{figures/fig_transformer_completo}
% ============================================================
\begin{figure}[H]
\centering
\begin{tikzpicture}[
    blk/.style={draw, rounded corners, minimum width=3.0cm,
                minimum height=0.62cm, font=\footnotesize, align=center},
    enc/.style={blk, fill=blue!10, draw=blue!40},
    dec/.style={blk, fill=green!10, draw=green!50!black},
    norm/.style={blk, fill=orange!10, draw=orange!50,
                 minimum height=0.5cm},
    emb/.style={blk, fill=purple!10, draw=purple!40},
    out/.style={blk, fill=red!8, draw=red!40},
    add/.style={draw, circle, fill=yellow!20, draw=yellow!60!black,
                minimum size=0.42cm, font=\scriptsize\bfseries},
    arrow/.style={-{Stealth}, thick},
    skip/.style={-{Stealth}, thick, gray!50, dashed},
    lbl/.style={font=\tiny, gray}
]

%═══════════════════════════════════════════════════════════
% ENCODER (izquierda)
%═══════════════════════════════════════════════════════════
\node[font=\small\bfseries, blue!60] at (-3.2, 10.8) {ENCODER};

% Entrada
\node[lbl] (src) at (-3.2, -0.4) {secuencia de entrada};
\node[emb] (src_emb) at (-3.2, 0.55) {Input Embedding};
\node[emb] (src_pe)  at (-3.2, 1.45) {Positional Encoding};

% Bloque encoder (representado como caja grande)
\draw[blue!30, rounded corners, thick, fill=blue!3]
    (-4.9, 2.1) rectangle (-1.5, 7.9);
\node[lbl] at (-4.55, 7.65) {$\times N$};

\node[enc]  (mha_e)  at (-3.2, 2.75) {Multi-Head\\Self-Attention};
\node[add]  (add1e)  at (-3.2, 3.75) {$+$};
\node[norm] (ln1e)   at (-3.2, 4.4)  {Add \& Norm};
\node[enc]  (ffn_e)  at (-3.2, 5.3)  {Feed-Forward};
\node[add]  (add2e)  at (-3.2, 6.3)  {$+$};
\node[norm] (ln2e)   at (-3.2, 6.95) {Add \& Norm};

% Salida encoder
\node[lbl] (enc_out) at (-3.2, 8.3) {$K, V$ al decoder};

% Flechas encoder
\draw[arrow] (src)     -- (src_emb);
\draw[arrow] (src_emb) -- (src_pe);
\draw[arrow] (src_pe)  -- (mha_e);
\draw[arrow] (mha_e)   -- (add1e);
\draw[arrow] (add1e)   -- (ln1e);
\draw[arrow] (ln1e)    -- (ffn_e);
\draw[arrow] (ffn_e)   -- (add2e);
\draw[arrow] (add2e)   -- (ln2e);
\draw[arrow] (ln2e)    -- (enc_out);

% Residuales encoder
\draw[skip] (src_pe.west) -- ++(-0.7,0) -- ++(0,3.75) -- (add1e.west);
\draw[skip] (ln1e.west)   -- ++(-0.5,0) -- ++(0,2.0)  -- (add2e.west);

%═══════════════════════════════════════════════════════════
% DECODER (derecha)
%═══════════════════════════════════════════════════════════
\node[font=\small\bfseries, green!60!black] at (3.2, 10.8) {DECODER};

% Entrada
\node[lbl] (tgt) at (3.2, -0.4) {secuencia de salida (desplazada)};
\node[emb] (tgt_emb) at (3.2, 0.55) {Output Embedding};
\node[emb] (tgt_pe)  at (3.2, 1.45) {Positional Encoding};

% Bloque decoder
\draw[green!30, rounded corners, thick, fill=green!3]
    (1.5, 2.1) rectangle (4.9, 9.1);
\node[lbl] at (4.55, 8.85) {$\times N$};

\node[dec]  (mha_d1) at (3.2, 2.75) {Masked Multi-Head\\Self-Attention};
\node[add]  (add1d)  at (3.2, 3.75) {$+$};
\node[norm] (ln1d)   at (3.2, 4.4)  {Add \& Norm};
\node[dec]  (mha_d2) at (3.2, 5.3)  {Multi-Head\\Cross-Attention};
\node[add]  (add2d)  at (3.2, 6.3)  {$+$};
\node[norm] (ln2d)   at (3.2, 6.95) {Add \& Norm};
\node[dec]  (ffn_d)  at (3.2, 7.85) {Feed-Forward};
\node[add]  (add3d)  at (3.2, 8.75) {$+$};
% (ln3 dentro del bloque, salida)

% Flechas decoder
\draw[arrow] (tgt)     -- (tgt_emb);
\draw[arrow] (tgt_emb) -- (tgt_pe);
\draw[arrow] (tgt_pe)  -- (mha_d1);
\draw[arrow] (mha_d1)  -- (add1d);
\draw[arrow] (add1d)   -- (ln1d);
\draw[arrow] (ln1d)    -- (mha_d2);
\draw[arrow] (mha_d2)  -- (add2d);
\draw[arrow] (add2d)   -- (ln2d);
\draw[arrow] (ln2d)    -- (ffn_d);
\draw[arrow] (ffn_d)   -- (add3d);

% Residuales decoder
\draw[skip] (tgt_pe.east) -- ++(0.7,0) -- ++(0,3.75) -- (add1d.east);
\draw[skip] (ln1d.east)   -- ++(0.5,0) -- ++(0,2.0)  -- (add2d.east);
\draw[skip] (ln2d.east)   -- ++(0.5,0) -- ++(0,2.0)  -- (add3d.east);

% ── Conexión encoder → decoder (cross-attention) ─────────────
\draw[arrow, red!60, very thick]
    (enc_out.east) -- ++(0.5,0) -- ++(0,-3.0) -- (mha_d2.west);
\node[font=\tiny, red!60, align=center] at (0, 6.8)
    {$K$, $V$ del\\encoder};

% ── Salida final ──────────────────────────────────────────────
\node[out]  (linear)  at (3.2, 9.8)  {Linear};
\node[out]  (softmax) at (3.2, 10.5) {Softmax};
\node[lbl]           at (3.2, 11.1)  {probabilidades de salida};

\draw[arrow] (add3d)  -- (linear);
\draw[arrow] (linear)  -- (softmax);

% ── Nota máscara ─────────────────────────────────────────────
\node[draw=gray!30, fill=gray!5, rounded corners, font=\tiny,
      inner sep=4pt, align=left] at (6.8, 2.75)
    {\textbf{Masked:} m\'ascara causal\\impide ver tokens futuros.\\Garantiza generaci\'on\\autoregresiva.};
\draw[-{Stealth}, gray!40, thin] (4.9, 2.75) -- (5.5, 2.75);

\end{tikzpicture}
\caption{Arquitectura completa del Transformer original~\cite{vaswani2017attention}.
El encoder (izquierda) procesa la secuencia de entrada en paralelo y produce
representaciones contextualizadas. El decoder (derecha) genera la secuencia
de salida token a token, utilizando atenci\'on enmascarada sobre su propia
salida parcial y atenci\'on cruzada (\textit{cross-attention}) sobre las
representaciones del encoder. Las flechas rojas indican el flujo de $K$ y $V$
del encoder al decoder.}
\label{fig:transformer_completo}
\end{figure}

% ============================================================
\begin{figure}[H]
\centering
\begin{tikzpicture}[
    blk/.style={draw, rounded corners, minimum width=3.0cm,
                minimum height=0.62cm, font=\footnotesize, align=center},
    enc/.style={blk, fill=blue!10, draw=blue!40},
    dec/.style={blk, fill=green!10, draw=green!50!black},
    norm/.style={blk, fill=orange!10, draw=orange!50,
                 minimum height=0.5cm},
    emb/.style={blk, fill=purple!10, draw=purple!40},
    outblk/.style={blk, fill=red!8, draw=red!40},
    add/.style={draw, circle, fill=yellow!20, draw=yellow!60!black,
                minimum size=0.42cm, font=\scriptsize\bfseries},
    arrow/.style={-{Stealth}, thick},
    skip/.style={-{Stealth}, thick, gray!50, dashed},
    lbl/.style={font=\tiny, gray}
]

%═══════════════════════════════════════════════════════════
% ENCODER (izquierda)
%═══════════════════════════════════════════════════════════
\node[font=\small\bfseries, blue!60] at (-3.2, 10.8) {ENCODER};

% Entrada
\node[lbl] (src) at (-3.2, -0.4) {secuencia de entrada};
\node[emb] (src_emb) at (-3.2, 0.55) {Input Embedding};
\node[emb] (src_pe)  at (-3.2, 1.45) {Positional Encoding};

% Bloque encoder (representado como caja grande)
\draw[blue!30, rounded corners, thick, fill=blue!3]
    (-4.9, 2.1) rectangle (-1.5, 7.9);
\node[lbl] at (-4.55, 7.65) {$\times N$};

\node[enc]  (mha_e)  at (-3.2, 2.75) {Multi-Head\\Self-Attention};
\node[add]  (add1e)  at (-3.2, 3.75) {$+$};
\node[norm] (ln1e)   at (-3.2, 4.4)  {Add \& Norm};
\node[enc]  (ffn_e)  at (-3.2, 5.3)  {Feed-Forward};
\node[add]  (add2e)  at (-3.2, 6.3)  {$+$};
\node[norm] (ln2e)   at (-3.2, 6.95) {Add \& Norm};

% Salida encoder
\node[lbl] (enc_out) at (-3.2, 8.3) {$K, V$ al decoder};

% Flechas encoder
\draw[arrow] (src)     -- (src_emb);
\draw[arrow] (src_emb) -- (src_pe);
\draw[arrow] (src_pe)  -- (mha_e);
\draw[arrow] (mha_e)   -- (add1e);
\draw[arrow] (add1e)   -- (ln1e);
\draw[arrow] (ln1e)    -- (ffn_e);
\draw[arrow] (ffn_e)   -- (add2e);
\draw[arrow] (add2e)   -- (ln2e);
\draw[arrow] (ln2e)    -- (enc_out);

% Residuales encoder
\draw[skip] (src_pe.west) -- ++(-0.7,0) -- ++(0,3.75) -- (add1e.west);
\draw[skip] (ln1e.west)   -- ++(-0.5,0) -- ++(0,2.0)  -- (add2e.west);

%═══════════════════════════════════════════════════════════
% DECODER (derecha)
%═══════════════════════════════════════════════════════════
\node[font=\small\bfseries, green!60!black] at (3.2, 10.8) {DECODER};

% Entrada
\node[lbl] (tgt) at (3.2, -0.4) {secuencia de salida (desplazada)};
\node[emb] (tgt_emb) at (3.2, 0.55) {Output Embedding};
\node[emb] (tgt_pe)  at (3.2, 1.45) {Positional Encoding};

% Bloque decoder
\draw[green!30, rounded corners, thick, fill=green!3]
    (1.5, 2.1) rectangle (4.9, 9.1);
\node[lbl] at (4.55, 8.85) {$\times N$};

\node[dec]  (mha_d1) at (3.2, 2.75) {Masked Multi-Head\\Self-Attention};
\node[add]  (add1d)  at (3.2, 3.75) {$+$};
\node[norm] (ln1d)   at (3.2, 4.4)  {Add \& Norm};
\node[dec]  (mha_d2) at (3.2, 5.3)  {Multi-Head\\Cross-Attention};
\node[add]  (add2d)  at (3.2, 6.3)  {$+$};
\node[norm] (ln2d)   at (3.2, 6.95) {Add \& Norm};
\node[dec]  (ffn_d)  at (3.2, 7.85) {Feed-Forward};
\node[add]  (add3d)  at (3.2, 8.75) {$+$};
% (ln3 dentro del bloque, salida)

% Flechas decoder
\draw[arrow] (tgt)     -- (tgt_emb);
\draw[arrow] (tgt_emb) -- (tgt_pe);
\draw[arrow] (tgt_pe)  -- (mha_d1);
\draw[arrow] (mha_d1)  -- (add1d);
\draw[arrow] (add1d)   -- (ln1d);
\draw[arrow] (ln1d)    -- (mha_d2);
\draw[arrow] (mha_d2)  -- (add2d);
\draw[arrow] (add2d)   -- (ln2d);
\draw[arrow] (ln2d)    -- (ffn_d);
\draw[arrow] (ffn_d)   -- (add3d);

% Residuales decoder
\draw[skip] (tgt_pe.east) -- ++(0.7,0) -- ++(0,3.75) -- (add1d.east);
\draw[skip] (ln1d.east)   -- ++(0.5,0) -- ++(0,2.0)  -- (add2d.east);
\draw[skip] (ln2d.east)   -- ++(0.5,0) -- ++(0,2.0)  -- (add3d.east);

% ── Conexión encoder → decoder (cross-attention) ─────────────
\draw[arrow, red!60, very thick]
    (enc_out.east) -- ++(0.5,0) -- ++(0,-3.0) -- (mha_d2.west);
\node[font=\tiny, red!60, align=center] at (0, 6.8)
    {$K$, $V$ del\\encoder};

% ── Salida final ──────────────────────────────────────────────
\node[outblk]  (linear)  at (3.2, 9.8)  {Linear};
\node[outblk]  (softmax) at (3.2, 10.5) {Softmax};
\node[lbl]           at (3.2, 11.1)  {probabilidades de salida};

\draw[arrow] (add3d)  -- (linear);
\draw[arrow] (linear)  -- (softmax);

% ── Nota máscara ─────────────────────────────────────────────
\node[draw=gray!30, fill=gray!5, rounded corners, font=\tiny,
      inner sep=4pt, align=left] at (6.8, 2.75)
    {\textbf{Masked:} m\'ascara causal\\impide ver tokens futuros.\\Garantiza generaci\'on\\autoregresiva.};
\draw[-{Stealth}, gray!40, thin] (4.9, 2.75) -- (5.5, 2.75);

\end{tikzpicture}
\caption{Arquitectura completa del Transformer original~\cite{vaswani2017attention}.
El encoder (izquierda) procesa la secuencia de entrada en paralelo y produce
representaciones contextualizadas. El decoder (derecha) genera la secuencia
de salida token a token, utilizando atenci\'on enmascarada sobre su propia
salida parcial y atenci\'on cruzada (\textit{cross-attention}) sobre las
representaciones del encoder. Las flechas rojas indican el flujo de $K$ y $V$
del encoder al decoder.}
\label{fig:positional_encoding}
\end{figure}


%------------------------------------------------------------
\subsection{El bloque encoder}
\label{subsec:bloque_encoder}

El bloque encoder es la unidad fundamental de procesamiento. Recibe una
secuencia de vectores y devuelve una secuencia del mismo tamaño, donde
cada vector es una representación contextualizada de su token
correspondiente. Está compuesto por dos sub-capas:

\begin{enumerate}
    \item \textbf{Multi-Head Self-Attention:} cada token consulta a todos
    los demás para construir su representación contextualizada, como se
    describió en la Sección~\ref{sec:mecanismo_atencion}.

    \item \textbf{Feed-Forward Network (FFN):} una red de dos capas
    lineales con activación no lineal, aplicada de forma idéntica e
    independiente a cada posición:
    \begin{equation}
        \text{FFN}(\bm{x}) = \max(0,\, \bm{x} W_1 + \bm{b}_1)\, W_2 + \bm{b}_2
    \end{equation}
    La FFN actúa como una memoria de patrones por posición: mientras que
    la atención mezcla información entre tokens, la FFN transforma cada
    representación de forma local.
\end{enumerate}

Ambas sub-capas están envueltas en dos mecanismos de estabilización:

\begin{itemize}
    \item \textbf{Conexión residual (\textit{skip connection}):} la entrada
    de cada sub-capa se suma a su salida antes de la normalización:
    $\text{salida} = \text{sub-capa}(\bm{x}) + \bm{x}$. Esto permite que
    los gradientes fluyan directamente hacia las capas anteriores durante
    el entrenamiento, resolviendo el problema del gradiente desvanecido
    que afectaba a las RNNs.

    \item \textbf{Layer Normalization:} normaliza las activaciones de cada
    vector de forma independiente, estabilizando el entrenamiento y
    acelerando la convergencia~\cite{goodfellow2016deeplearning}.
\end{itemize}

% ============================================================
% fig_encoder_block.tex: Bloque Encoder detallado
% Uso: % ============================================================
% fig_encoder_block.tex: Bloque Encoder detallado
% Uso: \input{figures/fig_encoder_block}
% ============================================================
\begin{figure}[H]
\centering
\begin{tikzpicture}[
    block/.style={draw, rounded corners, fill=blue!8, draw=blue!40,
                  minimum width=4.2cm, minimum height=0.7cm,
                  font=\small, align=center},
    norm/.style={draw, rounded corners, fill=orange!10, draw=orange!50,
                 minimum width=4.2cm, minimum height=0.55cm,
                 font=\small, align=center},
    add/.style={draw, circle, fill=yellow!15, draw=yellow!60!black,
                minimum size=0.5cm, font=\small\bfseries},
    arrow/.style={-{Stealth}, thick},
    skip/.style={-{Stealth}, thick, gray!60, dashed}
]

% ── Entrada ──────────────────────────────────────────────────
\node[font=\footnotesize, gray] (inp) at (0, 0) {$\bm{X}$ (embeddings + PE)};

% ── Sub-bloque 1: Multi-Head Self-Attention ──────────────────
\node[block] (mha)  at (0, 1.4) {Multi-Head\\Self-Attention};
\node[add]   (add1) at (0, 2.6) {$+$};
\node[norm]  (ln1)  at (0, 3.4) {Layer Norm};

% ── Sub-bloque 2: Feed-Forward ───────────────────────────────
\node[block] (ffn)  at (0, 4.6) {Feed-Forward\\Network};
\node[add]   (add2) at (0, 5.8) {$+$};
\node[norm]  (ln2)  at (0, 6.6) {Layer Norm};

% ── Salida ───────────────────────────────────────────────────
\node[font=\footnotesize, gray] (out) at (0, 7.6) {$\bm{Z}$ (representaciones contextualizadas)};

% ── Flechas principales ──────────────────────────────────────
\draw[arrow] (inp)  -- (mha);
\draw[arrow] (mha)  -- (add1);
\draw[arrow] (add1) -- (ln1);
\draw[arrow] (ln1)  -- (ffn);
\draw[arrow] (ffn)  -- (add2);
\draw[arrow] (add2) -- (ln2);
\draw[arrow] (ln2)  -- (out);

% ── Conexiones residuales ─────────────────────────────────────
% Residual 1: entrada → add1
\draw[skip] (inp.east) -- ++(1.6,0) -- ++(0,2.6) -- (add1.east);
\node[font=\tiny, gray] at (2.0, 1.3) {residual};

% Residual 2: después de ln1 → add2
\draw[skip] (ln1.east) -- ++(1.4,0) -- ++(0,2.3) -- (add2.east);

% ── Anotaciones laterales ─────────────────────────────────────
\node[draw=blue!20, fill=blue!4, rounded corners, font=\tiny,
      inner sep=4pt, align=left] at (-3.2, 1.4)
    {Cada token consulta\\a todos los dem\'as.\\Produce representaci\'on\\contextualizada.};
\draw[-{Stealth}, gray!50, thin] (-1.8, 1.4) -- (-2.1, 1.4);

\node[draw=green!20, fill=green!4, rounded corners, font=\tiny,
      inner sep=4pt, align=left] at (-3.2, 4.6)
    {Transf. no lineal\\aplicada a cada token\\independientemente.\\ReLU o GELU.};
\draw[-{Stealth}, gray!50, thin] (-1.8, 4.6) -- (-2.1, 4.6);

\node[draw=orange!20, fill=orange!4, rounded corners, font=\tiny,
      inner sep=4pt, align=left] at (3.2, 3.4)
    {Normaliza activaciones.\\Estabiliza entrenamiento.};
\draw[-{Stealth}, gray!50, thin] (1.8, 3.4) -- (2.1, 3.4);

% ── Corchete del bloque ───────────────────────────────────────
\draw[gray!50, thick]
    (-2.6, 1.0) -- (-2.8, 1.0) -- (-2.8, 7.0) -- (-2.6, 7.0);
\node[font=\footnotesize, gray, rotate=90] at (-3.1, 4.0)
    {$\times N$ bloques};

\end{tikzpicture}
\caption{Estructura de un bloque encoder. Cada bloque contiene dos
sub-capas: atenci\'on multi-cabeza y una red feed-forward. Ambas est\'an
rodeadas por una conexi\'on residual (\textit{skip connection}) y una
normalizaci\'on de capa (\textit{Layer Norm}). El bloque se apila $N$ veces
para aumentar la capacidad representacional del modelo~\cite{vaswani2017attention}.}
\label{fig:encoder_block}
\end{figure}

% ============================================================
\begin{figure}[H]
\centering
\begin{tikzpicture}[
    block/.style={draw, rounded corners, fill=blue!8, draw=blue!40,
                  minimum width=4.2cm, minimum height=0.7cm,
                  font=\small, align=center},
    norm/.style={draw, rounded corners, fill=orange!10, draw=orange!50,
                 minimum width=4.2cm, minimum height=0.55cm,
                 font=\small, align=center},
    add/.style={draw, circle, fill=yellow!15, draw=yellow!60!black,
                minimum size=0.5cm, font=\small\bfseries},
    arrow/.style={-{Stealth}, thick},
    skip/.style={-{Stealth}, thick, gray!60, dashed}
]

% ── Entrada ──────────────────────────────────────────────────
\node[font=\footnotesize, gray] (inp) at (0, 0) {$\bm{X}$ (embeddings + PE)};

% ── Sub-bloque 1: Multi-Head Self-Attention ──────────────────
\node[block] (mha)  at (0, 1.4) {Multi-Head\\Self-Attention};
\node[add]   (add1) at (0, 2.6) {$+$};
\node[norm]  (ln1)  at (0, 3.4) {Layer Norm};

% ── Sub-bloque 2: Feed-Forward ───────────────────────────────
\node[block] (ffn)  at (0, 4.6) {Feed-Forward\\Network};
\node[add]   (add2) at (0, 5.8) {$+$};
\node[norm]  (ln2)  at (0, 6.6) {Layer Norm};

% ── Salida ───────────────────────────────────────────────────
\node[font=\footnotesize, gray] (out) at (0, 7.6) {$\bm{Z}$ (representaciones contextualizadas)};

% ── Flechas principales ──────────────────────────────────────
\draw[arrow] (inp)  -- (mha);
\draw[arrow] (mha)  -- (add1);
\draw[arrow] (add1) -- (ln1);
\draw[arrow] (ln1)  -- (ffn);
\draw[arrow] (ffn)  -- (add2);
\draw[arrow] (add2) -- (ln2);
\draw[arrow] (ln2)  -- (out);

% ── Conexiones residuales ─────────────────────────────────────
% Residual 1: entrada → add1
\draw[skip] (inp.east) -- ++(1.6,0) -- ++(0,2.6) -- (add1.east);
\node[font=\tiny, gray] at (2.0, 1.3) {residual};

% Residual 2: después de ln1 → add2
\draw[skip] (ln1.east) -- ++(1.4,0) -- ++(0,2.3) -- (add2.east);

% ── Anotaciones laterales ─────────────────────────────────────
\node[draw=blue!20, fill=blue!4, rounded corners, font=\tiny,
      inner sep=4pt, align=left] at (-3.2, 1.4)
    {Cada token consulta\\a todos los dem\'as.\\Produce representaci\'on\\contextualizada.};
\draw[-{Stealth}, gray!50, thin] (-1.8, 1.4) -- (-2.1, 1.4);

\node[draw=green!20, fill=green!4, rounded corners, font=\tiny,
      inner sep=4pt, align=left] at (-3.2, 4.6)
    {Transf. no lineal\\aplicada a cada token\\independientemente.\\ReLU o GELU.};
\draw[-{Stealth}, gray!50, thin] (-1.8, 4.6) -- (-2.1, 4.6);

\node[draw=orange!20, fill=orange!4, rounded corners, font=\tiny,
      inner sep=4pt, align=left] at (3.2, 3.4)
    {Normaliza activaciones.\\Estabiliza entrenamiento.};
\draw[-{Stealth}, gray!50, thin] (1.8, 3.4) -- (2.1, 3.4);

% ── Corchete del bloque ───────────────────────────────────────
\draw[gray!50, thick]
    (-2.6, 1.0) -- (-2.8, 1.0) -- (-2.8, 7.0) -- (-2.6, 7.0);
\node[font=\footnotesize, gray, rotate=90] at (-3.1, 4.0)
    {$\times N$ bloques};

\end{tikzpicture}
\caption{Estructura de un bloque encoder. Cada bloque contiene dos
sub-capas: atenci\'on multi-cabeza y una red feed-forward. Ambas est\'an
rodeadas por una conexi\'on residual (\textit{skip connection}) y una
normalizaci\'on de capa (\textit{Layer Norm}). El bloque se apila $N$ veces
para aumentar la capacidad representacional del modelo~\cite{vaswani2017attention}.}
\label{fig:encoder_block}
\end{figure}


El encoder completo apila $N$ de estos bloques. Cada bloque puede
especializarse en distintos tipos de relaciones: los primeros bloques
tienden a capturar patrones locales y sintácticos, mientras que los
bloques más profundos capturan relaciones semánticas de mayor
abstracción~\cite{vaswani2017attention}.

%------------------------------------------------------------
\subsection{El bloque decoder}
\label{subsec:bloque_decoder}

El decoder extiende el encoder con un mecanismo adicional para incorporar
información de la secuencia de entrada. Se utiliza en tareas donde la
salida es una secuencia distinta de la entrada, como la traducción
automática o la generación de resúmenes. Consta de tres sub-capas:

\begin{enumerate}
    \item \textbf{Masked Multi-Head Self-Attention:} idéntica a la del
    encoder, pero con una máscara causal que impide que el token en
    posición $t$ acceda a tokens en posiciones $t+1, t+2, \ldots$
    Esto garantiza la propiedad autoregresiva: el modelo genera cada
    token condicionado únicamente en los anteriores.

    \item \textbf{Multi-Head Cross-Attention:} en esta sub-capa, las
    \textit{queries} provienen del decoder (la secuencia generada hasta
    el momento) mientras que las \textit{keys} y \textit{values} provienen
    de la salida del encoder (la representación de la entrada). Esto
    permite al decoder ``consultar'' la secuencia de entrada para decidir
    qué información es relevante en cada paso de generación.

    \item \textbf{Feed-Forward Network:} idéntica a la del encoder.
\end{enumerate}

Las tres sub-capas también están envueltas en conexiones residuales y
Layer Normalization.

%------------------------------------------------------------
\subsection{Arquitectura completa}
\label{subsec:arquitectura_completa}

La Figura~\ref{fig:transformer_completo} muestra la arquitectura completa
del Transformer original, con encoder y decoder.

% ============================================================
% fig_transformer_completo.tex — Arquitectura Transformer completa
% Uso: % ============================================================
% fig_transformer_completo.tex — Arquitectura Transformer completa
% Uso: \input{figures/fig_transformer_completo}
% ============================================================
\begin{figure}[H]
\centering
\begin{tikzpicture}[
    blk/.style={draw, rounded corners, minimum width=3.0cm,
                minimum height=0.62cm, font=\footnotesize, align=center},
    enc/.style={blk, fill=blue!10, draw=blue!40},
    dec/.style={blk, fill=green!10, draw=green!50!black},
    norm/.style={blk, fill=orange!10, draw=orange!50,
                 minimum height=0.5cm},
    emb/.style={blk, fill=purple!10, draw=purple!40},
    out/.style={blk, fill=red!8, draw=red!40},
    add/.style={draw, circle, fill=yellow!20, draw=yellow!60!black,
                minimum size=0.42cm, font=\scriptsize\bfseries},
    arrow/.style={-{Stealth}, thick},
    skip/.style={-{Stealth}, thick, gray!50, dashed},
    lbl/.style={font=\tiny, gray}
]

%═══════════════════════════════════════════════════════════
% ENCODER (izquierda)
%═══════════════════════════════════════════════════════════
\node[font=\small\bfseries, blue!60] at (-3.2, 10.8) {ENCODER};

% Entrada
\node[lbl] (src) at (-3.2, -0.4) {secuencia de entrada};
\node[emb] (src_emb) at (-3.2, 0.55) {Input Embedding};
\node[emb] (src_pe)  at (-3.2, 1.45) {Positional Encoding};

% Bloque encoder (representado como caja grande)
\draw[blue!30, rounded corners, thick, fill=blue!3]
    (-4.9, 2.1) rectangle (-1.5, 7.9);
\node[lbl] at (-4.55, 7.65) {$\times N$};

\node[enc]  (mha_e)  at (-3.2, 2.75) {Multi-Head\\Self-Attention};
\node[add]  (add1e)  at (-3.2, 3.75) {$+$};
\node[norm] (ln1e)   at (-3.2, 4.4)  {Add \& Norm};
\node[enc]  (ffn_e)  at (-3.2, 5.3)  {Feed-Forward};
\node[add]  (add2e)  at (-3.2, 6.3)  {$+$};
\node[norm] (ln2e)   at (-3.2, 6.95) {Add \& Norm};

% Salida encoder
\node[lbl] (enc_out) at (-3.2, 8.3) {$K, V$ al decoder};

% Flechas encoder
\draw[arrow] (src)     -- (src_emb);
\draw[arrow] (src_emb) -- (src_pe);
\draw[arrow] (src_pe)  -- (mha_e);
\draw[arrow] (mha_e)   -- (add1e);
\draw[arrow] (add1e)   -- (ln1e);
\draw[arrow] (ln1e)    -- (ffn_e);
\draw[arrow] (ffn_e)   -- (add2e);
\draw[arrow] (add2e)   -- (ln2e);
\draw[arrow] (ln2e)    -- (enc_out);

% Residuales encoder
\draw[skip] (src_pe.west) -- ++(-0.7,0) -- ++(0,3.75) -- (add1e.west);
\draw[skip] (ln1e.west)   -- ++(-0.5,0) -- ++(0,2.0)  -- (add2e.west);

%═══════════════════════════════════════════════════════════
% DECODER (derecha)
%═══════════════════════════════════════════════════════════
\node[font=\small\bfseries, green!60!black] at (3.2, 10.8) {DECODER};

% Entrada
\node[lbl] (tgt) at (3.2, -0.4) {secuencia de salida (desplazada)};
\node[emb] (tgt_emb) at (3.2, 0.55) {Output Embedding};
\node[emb] (tgt_pe)  at (3.2, 1.45) {Positional Encoding};

% Bloque decoder
\draw[green!30, rounded corners, thick, fill=green!3]
    (1.5, 2.1) rectangle (4.9, 9.1);
\node[lbl] at (4.55, 8.85) {$\times N$};

\node[dec]  (mha_d1) at (3.2, 2.75) {Masked Multi-Head\\Self-Attention};
\node[add]  (add1d)  at (3.2, 3.75) {$+$};
\node[norm] (ln1d)   at (3.2, 4.4)  {Add \& Norm};
\node[dec]  (mha_d2) at (3.2, 5.3)  {Multi-Head\\Cross-Attention};
\node[add]  (add2d)  at (3.2, 6.3)  {$+$};
\node[norm] (ln2d)   at (3.2, 6.95) {Add \& Norm};
\node[dec]  (ffn_d)  at (3.2, 7.85) {Feed-Forward};
\node[add]  (add3d)  at (3.2, 8.75) {$+$};
% (ln3 dentro del bloque, salida)

% Flechas decoder
\draw[arrow] (tgt)     -- (tgt_emb);
\draw[arrow] (tgt_emb) -- (tgt_pe);
\draw[arrow] (tgt_pe)  -- (mha_d1);
\draw[arrow] (mha_d1)  -- (add1d);
\draw[arrow] (add1d)   -- (ln1d);
\draw[arrow] (ln1d)    -- (mha_d2);
\draw[arrow] (mha_d2)  -- (add2d);
\draw[arrow] (add2d)   -- (ln2d);
\draw[arrow] (ln2d)    -- (ffn_d);
\draw[arrow] (ffn_d)   -- (add3d);

% Residuales decoder
\draw[skip] (tgt_pe.east) -- ++(0.7,0) -- ++(0,3.75) -- (add1d.east);
\draw[skip] (ln1d.east)   -- ++(0.5,0) -- ++(0,2.0)  -- (add2d.east);
\draw[skip] (ln2d.east)   -- ++(0.5,0) -- ++(0,2.0)  -- (add3d.east);

% ── Conexión encoder → decoder (cross-attention) ─────────────
\draw[arrow, red!60, very thick]
    (enc_out.east) -- ++(0.5,0) -- ++(0,-3.0) -- (mha_d2.west);
\node[font=\tiny, red!60, align=center] at (0, 6.8)
    {$K$, $V$ del\\encoder};

% ── Salida final ──────────────────────────────────────────────
\node[out]  (linear)  at (3.2, 9.8)  {Linear};
\node[out]  (softmax) at (3.2, 10.5) {Softmax};
\node[lbl]           at (3.2, 11.1)  {probabilidades de salida};

\draw[arrow] (add3d)  -- (linear);
\draw[arrow] (linear)  -- (softmax);

% ── Nota máscara ─────────────────────────────────────────────
\node[draw=gray!30, fill=gray!5, rounded corners, font=\tiny,
      inner sep=4pt, align=left] at (6.8, 2.75)
    {\textbf{Masked:} m\'ascara causal\\impide ver tokens futuros.\\Garantiza generaci\'on\\autoregresiva.};
\draw[-{Stealth}, gray!40, thin] (4.9, 2.75) -- (5.5, 2.75);

\end{tikzpicture}
\caption{Arquitectura completa del Transformer original~\cite{vaswani2017attention}.
El encoder (izquierda) procesa la secuencia de entrada en paralelo y produce
representaciones contextualizadas. El decoder (derecha) genera la secuencia
de salida token a token, utilizando atenci\'on enmascarada sobre su propia
salida parcial y atenci\'on cruzada (\textit{cross-attention}) sobre las
representaciones del encoder. Las flechas rojas indican el flujo de $K$ y $V$
del encoder al decoder.}
\label{fig:transformer_completo}
\end{figure}

% ============================================================
\begin{figure}[H]
\centering
\begin{tikzpicture}[
    blk/.style={draw, rounded corners, minimum width=3.0cm,
                minimum height=0.62cm, font=\footnotesize, align=center},
    enc/.style={blk, fill=blue!10, draw=blue!40},
    dec/.style={blk, fill=green!10, draw=green!50!black},
    norm/.style={blk, fill=orange!10, draw=orange!50,
                 minimum height=0.5cm},
    emb/.style={blk, fill=purple!10, draw=purple!40},
    out/.style={blk, fill=red!8, draw=red!40},
    add/.style={draw, circle, fill=yellow!20, draw=yellow!60!black,
                minimum size=0.42cm, font=\scriptsize\bfseries},
    arrow/.style={-{Stealth}, thick},
    skip/.style={-{Stealth}, thick, gray!50, dashed},
    lbl/.style={font=\tiny, gray}
]

%═══════════════════════════════════════════════════════════
% ENCODER (izquierda)
%═══════════════════════════════════════════════════════════
\node[font=\small\bfseries, blue!60] at (-3.2, 10.8) {ENCODER};

% Entrada
\node[lbl] (src) at (-3.2, -0.4) {secuencia de entrada};
\node[emb] (src_emb) at (-3.2, 0.55) {Input Embedding};
\node[emb] (src_pe)  at (-3.2, 1.45) {Positional Encoding};

% Bloque encoder (representado como caja grande)
\draw[blue!30, rounded corners, thick, fill=blue!3]
    (-4.9, 2.1) rectangle (-1.5, 7.9);
\node[lbl] at (-4.55, 7.65) {$\times N$};

\node[enc]  (mha_e)  at (-3.2, 2.75) {Multi-Head\\Self-Attention};
\node[add]  (add1e)  at (-3.2, 3.75) {$+$};
\node[norm] (ln1e)   at (-3.2, 4.4)  {Add \& Norm};
\node[enc]  (ffn_e)  at (-3.2, 5.3)  {Feed-Forward};
\node[add]  (add2e)  at (-3.2, 6.3)  {$+$};
\node[norm] (ln2e)   at (-3.2, 6.95) {Add \& Norm};

% Salida encoder
\node[lbl] (enc_out) at (-3.2, 8.3) {$K, V$ al decoder};

% Flechas encoder
\draw[arrow] (src)     -- (src_emb);
\draw[arrow] (src_emb) -- (src_pe);
\draw[arrow] (src_pe)  -- (mha_e);
\draw[arrow] (mha_e)   -- (add1e);
\draw[arrow] (add1e)   -- (ln1e);
\draw[arrow] (ln1e)    -- (ffn_e);
\draw[arrow] (ffn_e)   -- (add2e);
\draw[arrow] (add2e)   -- (ln2e);
\draw[arrow] (ln2e)    -- (enc_out);

% Residuales encoder
\draw[skip] (src_pe.west) -- ++(-0.7,0) -- ++(0,3.75) -- (add1e.west);
\draw[skip] (ln1e.west)   -- ++(-0.5,0) -- ++(0,2.0)  -- (add2e.west);

%═══════════════════════════════════════════════════════════
% DECODER (derecha)
%═══════════════════════════════════════════════════════════
\node[font=\small\bfseries, green!60!black] at (3.2, 10.8) {DECODER};

% Entrada
\node[lbl] (tgt) at (3.2, -0.4) {secuencia de salida (desplazada)};
\node[emb] (tgt_emb) at (3.2, 0.55) {Output Embedding};
\node[emb] (tgt_pe)  at (3.2, 1.45) {Positional Encoding};

% Bloque decoder
\draw[green!30, rounded corners, thick, fill=green!3]
    (1.5, 2.1) rectangle (4.9, 9.1);
\node[lbl] at (4.55, 8.85) {$\times N$};

\node[dec]  (mha_d1) at (3.2, 2.75) {Masked Multi-Head\\Self-Attention};
\node[add]  (add1d)  at (3.2, 3.75) {$+$};
\node[norm] (ln1d)   at (3.2, 4.4)  {Add \& Norm};
\node[dec]  (mha_d2) at (3.2, 5.3)  {Multi-Head\\Cross-Attention};
\node[add]  (add2d)  at (3.2, 6.3)  {$+$};
\node[norm] (ln2d)   at (3.2, 6.95) {Add \& Norm};
\node[dec]  (ffn_d)  at (3.2, 7.85) {Feed-Forward};
\node[add]  (add3d)  at (3.2, 8.75) {$+$};
% (ln3 dentro del bloque, salida)

% Flechas decoder
\draw[arrow] (tgt)     -- (tgt_emb);
\draw[arrow] (tgt_emb) -- (tgt_pe);
\draw[arrow] (tgt_pe)  -- (mha_d1);
\draw[arrow] (mha_d1)  -- (add1d);
\draw[arrow] (add1d)   -- (ln1d);
\draw[arrow] (ln1d)    -- (mha_d2);
\draw[arrow] (mha_d2)  -- (add2d);
\draw[arrow] (add2d)   -- (ln2d);
\draw[arrow] (ln2d)    -- (ffn_d);
\draw[arrow] (ffn_d)   -- (add3d);

% Residuales decoder
\draw[skip] (tgt_pe.east) -- ++(0.7,0) -- ++(0,3.75) -- (add1d.east);
\draw[skip] (ln1d.east)   -- ++(0.5,0) -- ++(0,2.0)  -- (add2d.east);
\draw[skip] (ln2d.east)   -- ++(0.5,0) -- ++(0,2.0)  -- (add3d.east);

% ── Conexión encoder → decoder (cross-attention) ─────────────
\draw[arrow, red!60, very thick]
    (enc_out.east) -- ++(0.5,0) -- ++(0,-3.0) -- (mha_d2.west);
\node[font=\tiny, red!60, align=center] at (0, 6.8)
    {$K$, $V$ del\\encoder};

% ── Salida final ──────────────────────────────────────────────
\node[out]  (linear)  at (3.2, 9.8)  {Linear};
\node[out]  (softmax) at (3.2, 10.5) {Softmax};
\node[lbl]           at (3.2, 11.1)  {probabilidades de salida};

\draw[arrow] (add3d)  -- (linear);
\draw[arrow] (linear)  -- (softmax);

% ── Nota máscara ─────────────────────────────────────────────
\node[draw=gray!30, fill=gray!5, rounded corners, font=\tiny,
      inner sep=4pt, align=left] at (6.8, 2.75)
    {\textbf{Masked:} m\'ascara causal\\impide ver tokens futuros.\\Garantiza generaci\'on\\autoregresiva.};
\draw[-{Stealth}, gray!40, thin] (4.9, 2.75) -- (5.5, 2.75);

\end{tikzpicture}
\caption{Arquitectura completa del Transformer original~\cite{vaswani2017attention}.
El encoder (izquierda) procesa la secuencia de entrada en paralelo y produce
representaciones contextualizadas. El decoder (derecha) genera la secuencia
de salida token a token, utilizando atenci\'on enmascarada sobre su propia
salida parcial y atenci\'on cruzada (\textit{cross-attention}) sobre las
representaciones del encoder. Las flechas rojas indican el flujo de $K$ y $V$
del encoder al decoder.}
\label{fig:transformer_completo}
\end{figure}


El flujo de información puede resumirse así: la secuencia de entrada pasa
por el encoder y produce un conjunto de representaciones ricas en contexto.
El decoder recibe estas representaciones a través del cross-attention y,
combinándolas con su propio estado interno, genera la secuencia de salida
token a token. En cada paso, el decoder produce una distribución de
probabilidad sobre el vocabulario y selecciona el token más probable (o
muestrea de la distribución), que luego se alimenta como entrada al
siguiente paso.

%------------------------------------------------------------
\subsection{Variantes principales}
\label{subsec:variantes}

La arquitectura original de Vaswani et al.\ incluye encoder y decoder.
Sin embargo, muchas aplicaciones prácticas utilizan únicamente uno de
los dos componentes, dependiendo de la tarea:

% ============================================================
% fig_variantes_transformer.tex — Variantes encoder-only vs enc-dec
% Uso: % ============================================================
% fig_variantes_transformer.tex — Variantes encoder-only vs enc-dec
% Uso: \input{figures/fig_variantes_transformer}
% ============================================================
\begin{figure}[H]
\centering
\begin{tikzpicture}[
    hdr/.style={font=\small\bfseries, align=center},
    cell/.style={font=\footnotesize, align=center},
    enconly/.style={draw=blue!40, fill=blue!8, rounded corners,
                    minimum width=2.2cm, minimum height=0.8cm,
                    font=\footnotesize\bfseries, align=center},
    encdec/.style={draw=green!50!black, fill=green!8, rounded corners,
                   minimum width=2.2cm, minimum height=0.8cm,
                   font=\footnotesize\bfseries, align=center}
]

% ── Mini diagramas ───────────────────────────────────────────

% Encoder-only
\node[enconly] at (1.5, 6.5) {Encoder\\$\times N$};
\node[font=\tiny, gray] at (1.5, 5.9) {entrada};
\node[font=\tiny, gray] at (1.5, 7.1) {repr. contextual};
\draw[-{Stealth}, thick, blue!50] (1.5, 6.1) -- (1.5, 6.15);
\draw[-{Stealth}, thick, blue!50] (1.5, 6.85) -- (1.5, 6.9);
\node[font=\small\bfseries, blue!60] at (1.5, 7.5) {Encoder-only};
\node[font=\tiny, gray] at (1.5, 7.9) {(ej: BERT)};

% Encoder-Decoder
\node[enconly]  at (6.5, 6.5) {Encoder\\$\times N$};
\node[encdec]   at (8.8, 6.5) {Decoder\\$\times N$};
\draw[-{Stealth}, thick, red!50] (7.7, 6.5) -- (8.05, 6.5);
\node[font=\tiny, red!60] at (7.87, 6.75) {$K,V$};
\node[font=\tiny, gray]   at (6.5, 5.9) {entrada};
\node[font=\tiny, gray]   at (8.8, 7.1) {salida};
\draw[-{Stealth}, thick, blue!50]  (6.5, 6.1) -- (6.5, 6.15);
\draw[-{Stealth}, thick, green!50!black] (8.8, 6.85) -- (8.8, 6.9);
\node[font=\small\bfseries, green!60!black] at (7.65, 7.5) {Encoder-Decoder};
\node[font=\tiny, gray] at (7.65, 7.9) {(ej: T5, BART)};

% Decoder-only
\node[encdec] at (12.5, 6.5) {Decoder\\(causal) $\times N$};
\node[font=\tiny, gray] at (12.5, 5.9) {contexto previo};
\node[font=\tiny, gray] at (12.5, 7.1) {siguiente token};
\draw[-{Stealth}, thick, green!50!black] (12.5, 6.1) -- (12.5, 6.15);
\draw[-{Stealth}, thick, green!50!black] (12.5, 6.85) -- (12.5, 6.9);
\node[font=\small\bfseries, green!60!black] at (12.5, 7.5) {Decoder-only};
\node[font=\tiny, gray] at (12.5, 7.9) {(ej: GPT, LLaMA)};

% ── Tabla comparativa ─────────────────────────────────────────
\node[hdr] at (0,    5.0) {Caracter\'istica};
\node[hdr, blue!60] at (3.5,  5.0) {Encoder-only};
\node[hdr, green!60!black] at (7.5,  5.0) {Encoder-Decoder};
\node[hdr, green!60!black] at (11.5, 5.0) {Decoder-only};

\draw[gray!40, thick] (-0.5, 4.7) -- (13.5, 4.7);

% Filas
\foreach \y/\lbl/\a/\b/\c in {
    4.1/{Direcci\'on atenci\'on}/{bidireccional}/{bidireccional (enc)}/{unidireccional},
    3.4/{Tarea principal}/{clasificaci\'on, NER}/{traducci\'on, res\'umenes}/{generaci\'on de texto},
    2.7/{Ve tokens futuros}/{S\'i}/{S\'i (encoder)}/{No (m\'ascara causal)},
    2.0/{Tama\~no relativo}/{mediano}/{grande}/{grande}
}{
    \node[cell, align=left] at (0, \y) {\lbl};
    \node[cell] at (3.5,  \y) {\a};
    \node[cell] at (7.5,  \y) {\b};
    \node[cell] at (11.5, \y) {\c};
    \draw[gray!20] (-0.5, \y-0.35) -- (13.5, \y-0.35);
}

\draw[gray!40] (-0.5, 5.3) -- (-0.5, 1.65);
\draw[gray!40] (2.0,  5.3) -- (2.0,  1.65);
\draw[gray!40] (5.5,  5.3) -- (5.5,  1.65);
\draw[gray!40] (9.5,  5.3) -- (9.5,  1.65);
\draw[gray!40] (13.5, 5.3) -- (13.5, 1.65);

\end{tikzpicture}
\caption{Las tres variantes principales de la arquitectura Transformer.
\textbf{Encoder-only} (BERT): procesa la entrada completa en ambas
direcciones; ideal para tareas de comprensi\'on. \textbf{Encoder-decoder}
(T5, BART): mapea una secuencia a otra; ideal para traducci\'on y res\'umenes.
\textbf{Decoder-only} (GPT, LLaMA): genera texto token a token con m\'ascara
causal; es la base de los modelos de lenguaje actuales.}
\label{fig:variantes_transformer}
\end{figure}

% ============================================================
\begin{figure}[H]
\centering
\begin{tikzpicture}[
    hdr/.style={font=\small\bfseries, align=center},
    cell/.style={font=\footnotesize, align=center},
    enconly/.style={draw=blue!40, fill=blue!8, rounded corners,
                    minimum width=2.2cm, minimum height=0.8cm,
                    font=\footnotesize\bfseries, align=center},
    encdec/.style={draw=green!50!black, fill=green!8, rounded corners,
                   minimum width=2.2cm, minimum height=0.8cm,
                   font=\footnotesize\bfseries, align=center}
]

% ── Mini diagramas ───────────────────────────────────────────

% Encoder-only
\node[enconly] at (1.5, 6.5) {Encoder\\$\times N$};
\node[font=\tiny, gray] at (1.5, 5.9) {entrada};
\node[font=\tiny, gray] at (1.5, 7.1) {repr. contextual};
\draw[-{Stealth}, thick, blue!50] (1.5, 6.1) -- (1.5, 6.15);
\draw[-{Stealth}, thick, blue!50] (1.5, 6.85) -- (1.5, 6.9);
\node[font=\small\bfseries, blue!60] at (1.5, 7.5) {Encoder-only};
\node[font=\tiny, gray] at (1.5, 7.9) {(ej: BERT)};

% Encoder-Decoder
\node[enconly]  at (6.5, 6.5) {Encoder\\$\times N$};
\node[encdec]   at (8.8, 6.5) {Decoder\\$\times N$};
\draw[-{Stealth}, thick, red!50] (7.7, 6.5) -- (8.05, 6.5);
\node[font=\tiny, red!60] at (7.87, 6.75) {$K,V$};
\node[font=\tiny, gray]   at (6.5, 5.9) {entrada};
\node[font=\tiny, gray]   at (8.8, 7.1) {salida};
\draw[-{Stealth}, thick, blue!50]  (6.5, 6.1) -- (6.5, 6.15);
\draw[-{Stealth}, thick, green!50!black] (8.8, 6.85) -- (8.8, 6.9);
\node[font=\small\bfseries, green!60!black] at (7.65, 7.5) {Encoder-Decoder};
\node[font=\tiny, gray] at (7.65, 7.9) {(ej: T5, BART)};

% Decoder-only
\node[encdec] at (12.5, 6.5) {Decoder\\(causal) $\times N$};
\node[font=\tiny, gray] at (12.5, 5.9) {contexto previo};
\node[font=\tiny, gray] at (12.5, 7.1) {siguiente token};
\draw[-{Stealth}, thick, green!50!black] (12.5, 6.1) -- (12.5, 6.15);
\draw[-{Stealth}, thick, green!50!black] (12.5, 6.85) -- (12.5, 6.9);
\node[font=\small\bfseries, green!60!black] at (12.5, 7.5) {Decoder-only};
\node[font=\tiny, gray] at (12.5, 7.9) {(ej: GPT, LLaMA)};

% ── Tabla comparativa ─────────────────────────────────────────
\node[hdr] at (0,    5.0) {Caracter\'istica};
\node[hdr, blue!60] at (3.5,  5.0) {Encoder-only};
\node[hdr, green!60!black] at (7.5,  5.0) {Encoder-Decoder};
\node[hdr, green!60!black] at (11.5, 5.0) {Decoder-only};

\draw[gray!40, thick] (-0.5, 4.7) -- (13.5, 4.7);

% Filas
\foreach \y/\lbl/\a/\b/\c in {
    4.1/{Direcci\'on atenci\'on}/{bidireccional}/{bidireccional (enc)}/{unidireccional},
    3.4/{Tarea principal}/{clasificaci\'on, NER}/{traducci\'on, res\'umenes}/{generaci\'on de texto},
    2.7/{Ve tokens futuros}/{S\'i}/{S\'i (encoder)}/{No (m\'ascara causal)},
    2.0/{Tama\~no relativo}/{mediano}/{grande}/{grande}
}{
    \node[cell, align=left] at (0, \y) {\lbl};
    \node[cell] at (3.5,  \y) {\a};
    \node[cell] at (7.5,  \y) {\b};
    \node[cell] at (11.5, \y) {\c};
    \draw[gray!20] (-0.5, \y-0.35) -- (13.5, \y-0.35);
}

\draw[gray!40] (-0.5, 5.3) -- (-0.5, 1.65);
\draw[gray!40] (2.0,  5.3) -- (2.0,  1.65);
\draw[gray!40] (5.5,  5.3) -- (5.5,  1.65);
\draw[gray!40] (9.5,  5.3) -- (9.5,  1.65);
\draw[gray!40] (13.5, 5.3) -- (13.5, 1.65);

\end{tikzpicture}
\caption{Las tres variantes principales de la arquitectura Transformer.
\textbf{Encoder-only} (BERT): procesa la entrada completa en ambas
direcciones; ideal para tareas de comprensi\'on. \textbf{Encoder-decoder}
(T5, BART): mapea una secuencia a otra; ideal para traducci\'on y res\'umenes.
\textbf{Decoder-only} (GPT, LLaMA): genera texto token a token con m\'ascara
causal; es la base de los modelos de lenguaje actuales.}
\label{fig:variantes_transformer}
\end{figure}


Las tres variantes comparten los mismos bloques fundamentales —embeddings,
codificación posicional, atención multi-cabeza, FFN, conexiones residuales
y Layer Norm— pero difieren en cómo los ensamblan y en qué dirección
procesan la información. Esta modularidad es una de las razones del éxito
del Transformer como arquitectura universal.

%------------------------------------------------------------
\subsection{Hiperparámetros principales}
\label{subsec:hiperparametros}

El comportamiento del Transformer está controlado por un conjunto de
hiperparámetros que determinan su capacidad y costo computacional. La
Tabla~\ref{tab:hiperparametros} resume los más importantes.

\begin{table}[H]
\centering
\caption{Hiperparámetros principales del Transformer y su efecto.}
\label{tab:hiperparametros}
\begin{tabularx}{\textwidth}{llX}
\toprule
\textbf{Parámetro} & \textbf{Notación} & \textbf{Efecto} \\
\midrule
Dimensión del modelo   & $d_{\text{model}}$ & Tamaño de todos los vectores internos. Afecta capacidad y memoria. \\
Número de cabezas      & $h$               & Cuántos tipos de relaciones puede capturar simultáneamente. \\
Número de bloques      & $N$               & Profundidad del modelo. Más bloques = más abstracción. \\
Dimensión FFN          & $d_{\text{ff}}$   & Generalmente $4 \times d_{\text{model}}$. Controla la ``memoria'' por posición. \\
Longitud máx.\ de secuencia & $n_{\text{max}}$ & Límite de contexto. El costo de atención escala como $O(n^2)$. \\
\bottomrule
\end{tabularx}
\end{table}

Para el prototipo mínimo implementado en este trabajo (descrito en la
Sección~\ref{sec:transformer_minimo}), se utilizan valores reducidos de
estos hiperparámetros que permiten el entrenamiento en hardware doméstico
sin sacrificar la comprensión de los mecanismos.

\section{Especificación del Transformer mínimo}
\label{sec:transformer_minimo}

Las secciones anteriores establecieron los fundamentos teóricos del
Transformer: el mecanismo de atención, la arquitectura completa y sus
variantes. Esta sección define formalmente el modelo concreto que se
implementa (ver el repositorio), especificando la tarea,
las entradas y salidas, la arquitectura y los hiperparámetros. El objetivo
es que cada decisión de diseño sea justificable en términos de los
conceptos teóricos ya presentados.

%------------------------------------------------------------
\subsection{Definición del problema}
\label{subsec:def_problema}

El modelo implementado resuelve la tarea de \textit{modelado de lenguaje
a nivel carácter}: dada una secuencia de caracteres $c_1, c_2, \ldots,
c_t$, predecir la distribución de probabilidad sobre el siguiente carácter
$c_{t+1}$.

\begin{equation}
    P(c_{t+1} \mid c_1, c_2, \ldots, c_t)
\end{equation}

Esta tarea es representativa de los modelos de lenguaje modernos a gran
escala, reducida a una escala manejable. Operar a nivel carácter elimina
la necesidad de tokenizadores complejos: el vocabulario $\mathcal{V}$
está compuesto simplemente por todos los caracteres únicos del corpus de
entrenamiento (letras, espacios, signos de puntuación), con
$|\mathcal{V}|$ típicamente entre 60 y 100 símbolos.

La justificación de esta elección frente a alternativas más simples
(como clasificación) es que el modelado de lenguaje autoregresivo
requiere y ejercita todos los componentes del Transformer: embeddings,
codificación posicional, atención con máscara causal y proyección de
salida. Ningún componente es redundante.

%------------------------------------------------------------
\subsection{Representación de la entrada}
\label{subsec:entrada}

Cada carácter $c_t \in \mathcal{V}$ se representa como un índice entero
$v_t \in \{0, 1, \ldots, |\mathcal{V}|-1\}$ mediante un mapeo fijo
aprendido durante el preprocesamiento. Este índice se convierte en un
vector denso mediante la capa de embeddings:

\begin{equation}
    \bm{e}_t = E[v_t] \quad \in \mathbb{R}^{d_{\text{model}}}
\end{equation}

\noindent donde $E \in \mathbb{R}^{|\mathcal{V}| \times d_{\text{model}}}$
es la matriz de embeddings aprendida. A continuación se suma la
codificación posicional sinusoidal descrita en la
Sección~\ref{subsec:positional_encoding}:

\begin{equation}
    \bm{x}'_t = \bm{e}_t + PE_t
\end{equation}

La secuencia resultante $\bm{x}'_1, \ldots, \bm{x}'_T$ constituye la
entrada al encoder, donde $T$ es la longitud de contexto máxima.

%------------------------------------------------------------
\subsection{Arquitectura del modelo}
\label{subsec:arq_modelo}

El modelo sigue la variante \textit{decoder-only} con máscara causal,
que es la base arquitectónica de los modelos generativos modernos
(familia GPT). Se elige esta variante en lugar del encoder-decoder
original porque la tarea es de generación autoregresiva sobre una sola
secuencia, y un encoder-decoder requeriría dos secuencias distintas
(entrada y objetivo).

Formalmente, el modelo está compuesto por:

\begin{enumerate}
    \item \textbf{Capa de embedding:}
    $E \in \mathbb{R}^{|\mathcal{V}| \times d_{\text{model}}}$

    \item \textbf{Codificación posicional sinusoidal:}
    fija, no aprendida, sumada a los embeddings.

    \item \textbf{$N$ bloques encoder con máscara causal,}
    cada uno con:
    \begin{itemize}
        \item Multi-Head Self-Attention con máscara triangular inferior
        \item Feed-Forward Network con activación ReLU
        \item Conexiones residuales y Layer Normalization
    \end{itemize}

    \item \textbf{Proyección lineal de salida:}
    $W_o \in \mathbb{R}^{d_{\text{model}} \times |\mathcal{V}|}$,
    que mapea la representación final al espacio del vocabulario.

    \item \textbf{Softmax:} produce la distribución de probabilidad
    $\hat{P} \in \Delta^{|\mathcal{V}|-1}$ sobre el siguiente carácter.
\end{enumerate}

% ============================================================
% fig_pipeline_minimo.tex: Pipeline del Transformer mínimo
% Uso: % ============================================================
% fig_pipeline_minimo.tex: Pipeline del Transformer mínimo
% Uso: \input{figures/fig_pipeline_minimo}
% ============================================================
\begin{figure}[H]
\centering
\begin{tikzpicture}[
    blk/.style={draw, rounded corners, minimum width=2.4cm,
                minimum height=0.7cm, font=\small, align=center},
    inp/.style={blk, fill=gray!10,    draw=gray!50},
    emb/.style={blk, fill=purple!10,  draw=purple!40},
    posblk/.style={blk, fill=teal!10,    draw=teal!50},
    enc/.style={blk, fill=blue!10,    draw=blue!40},
    outblk/.style={blk, fill=orange!10,  draw=orange!50},
    sft/.style={blk, fill=green!10,   draw=green!50!black},
    arrow/.style={-{Stealth}, thick},
    lbl/.style={font=\tiny, gray, align=center}
]

% ── Tokens de entrada ────────────────────────────────────────
\node[inp] (chars) at (0, 0) {``hola''\\{\tiny $[h,o,l,a]$}};
\node[lbl] at (0, -0.75) {caracteres};

% ── Índices enteros ──────────────────────────────────────────
\node[inp] (idx) at (2.8, 0) {índices\\{\tiny $[7,14,11,0]$}};
\node[lbl] at (2.8, -0.75) {vocabulario $|\mathcal{V}|$};

% ── Embedding ─────────────────────────────────────────────────
\node[emb] (emb) at (5.6, 0) {Embedding\\{\tiny $\mathbb{R}^{n \times d_{\text{model}}}$}};
\node[lbl] at (5.6, -0.75) {$E \in \mathbb{R}^{|\mathcal{V}| \times d}$};

% ── Positional Encoding ───────────────────────────────────────
\node[posblk] (pe) at (8.4, 0) {Positional\\Encoding};
\node[lbl] at (8.4, -0.75) {$\bm{x}'_t = \bm{e}_t + PE_t$};

% ── Encoder (N bloques) ───────────────────────────────────────
\node[enc] (enc) at (11.2, 0) {Encoder\\{\tiny $\times N$ bloques}};
\node[lbl] at (11.2, -0.75) {MHA + FFN + Norm};

% ── Proyección lineal ─────────────────────────────────────────
\node[outblk] (lin) at (11.2, -2.0) {Linear\\{\tiny $\mathbb{R}^{d} \to \mathbb{R}^{|\mathcal{V}|}$}};

% ── Softmax ──────────────────────────────────────────────────
\node[sft] (sft) at (8.4, -2.0) {Softmax};
\node[lbl] at (8.4, -2.75) {$P(\text{siguiente car.})$};

% ── Token predicho ────────────────────────────────────────────
\node[inp] (pred) at (5.6, -2.0) {``?''\\{\tiny token predicho}};
\node[lbl] at (5.6, -2.75) {argmax o muestreo};

% ── Flechas forward ───────────────────────────────────────────
\draw[arrow] (chars) -- (idx);
\draw[arrow] (idx)   -- (emb);
\draw[arrow] (emb)   -- (pe);
\draw[arrow] (pe)    -- (enc);
\draw[arrow] (enc)   -- (lin);
\draw[arrow] (lin)   -- (sft);
\draw[arrow] (sft)   -- (pred);

% ── Anotación: máscara causal ─────────────────────────────────
\node[draw=blue!20, fill=blue!4, rounded corners, font=\tiny,
      inner sep=4pt, align=left] at (11.2, 1.35)
    {M\'ascara causal:\\el token $t$ solo\\ve $1,\ldots,t{-}1$};
\draw[-{Stealth}, gray!40, thin] (11.2, 0.65) -- (11.2, 0.95);

% ── Objetivo de entrenamiento ─────────────────────────────────
\node[draw=orange!30, fill=orange!4, rounded corners, font=\tiny,
      inner sep=5pt, align=center] at (3.0, -2.0)
    {Entrenamiento:\\$\mathcal{L} = -\log P(c_{t+1} \mid c_1,\ldots,c_t)$\\
     Entrop\'ia cruzada sobre el\\siguiente car\'acter correcto};

% ── Flechas de flujo de la pérdida ────────────────────────────
\draw[-{Stealth}, orange!50, dashed, thick]
    (pred.west) -- +(-0.3, 0) -- +(-0.3, -0.7)
    node[below, font=\tiny, orange!70] {$\hat{y}$ vs $y$};

\end{tikzpicture}
\caption{Pipeline completo del Transformer m\'inimo implementado en este
trabajo. La secuencia de caracteres de entrada se convierte en \'indices
enteros, se proyecta al espacio de embeddings, se le suma la codificaci\'on
posicional, y se procesa a trav\'es de $N$ bloques encoder con m\'ascara
causal. La proyecci\'on lineal final y el softmax producen una distribuci\'on
de probabilidad sobre el vocabulario de caracteres. Durante el entrenamiento,
la p\'erdida de entrop\'ia cruzada compara la predicci\'on con el siguiente
car\'acter real de la secuencia.}
\label{fig:pipeline_minimo}
\end{figure}

% ============================================================
\begin{figure}[H]
\centering
\begin{tikzpicture}[
    blk/.style={draw, rounded corners, minimum width=2.4cm,
                minimum height=0.7cm, font=\small, align=center},
    inp/.style={blk, fill=gray!10,    draw=gray!50},
    emb/.style={blk, fill=purple!10,  draw=purple!40},
    posblk/.style={blk, fill=teal!10,    draw=teal!50},
    enc/.style={blk, fill=blue!10,    draw=blue!40},
    outblk/.style={blk, fill=orange!10,  draw=orange!50},
    sft/.style={blk, fill=green!10,   draw=green!50!black},
    arrow/.style={-{Stealth}, thick},
    lbl/.style={font=\tiny, gray, align=center}
]

% ── Tokens de entrada ────────────────────────────────────────
\node[inp] (chars) at (0, 0) {``hola''\\{\tiny $[h,o,l,a]$}};
\node[lbl] at (0, -0.75) {caracteres};

% ── Índices enteros ──────────────────────────────────────────
\node[inp] (idx) at (2.8, 0) {índices\\{\tiny $[7,14,11,0]$}};
\node[lbl] at (2.8, -0.75) {vocabulario $|\mathcal{V}|$};

% ── Embedding ─────────────────────────────────────────────────
\node[emb] (emb) at (5.6, 0) {Embedding\\{\tiny $\mathbb{R}^{n \times d_{\text{model}}}$}};
\node[lbl] at (5.6, -0.75) {$E \in \mathbb{R}^{|\mathcal{V}| \times d}$};

% ── Positional Encoding ───────────────────────────────────────
\node[posblk] (pe) at (8.4, 0) {Positional\\Encoding};
\node[lbl] at (8.4, -0.75) {$\bm{x}'_t = \bm{e}_t + PE_t$};

% ── Encoder (N bloques) ───────────────────────────────────────
\node[enc] (enc) at (11.2, 0) {Encoder\\{\tiny $\times N$ bloques}};
\node[lbl] at (11.2, -0.75) {MHA + FFN + Norm};

% ── Proyección lineal ─────────────────────────────────────────
\node[outblk] (lin) at (11.2, -2.0) {Linear\\{\tiny $\mathbb{R}^{d} \to \mathbb{R}^{|\mathcal{V}|}$}};

% ── Softmax ──────────────────────────────────────────────────
\node[sft] (sft) at (8.4, -2.0) {Softmax};
\node[lbl] at (8.4, -2.75) {$P(\text{siguiente car.})$};

% ── Token predicho ────────────────────────────────────────────
\node[inp] (pred) at (5.6, -2.0) {``?''\\{\tiny token predicho}};
\node[lbl] at (5.6, -2.75) {argmax o muestreo};

% ── Flechas forward ───────────────────────────────────────────
\draw[arrow] (chars) -- (idx);
\draw[arrow] (idx)   -- (emb);
\draw[arrow] (emb)   -- (pe);
\draw[arrow] (pe)    -- (enc);
\draw[arrow] (enc)   -- (lin);
\draw[arrow] (lin)   -- (sft);
\draw[arrow] (sft)   -- (pred);

% ── Anotación: máscara causal ─────────────────────────────────
\node[draw=blue!20, fill=blue!4, rounded corners, font=\tiny,
      inner sep=4pt, align=left] at (11.2, 1.35)
    {M\'ascara causal:\\el token $t$ solo\\ve $1,\ldots,t{-}1$};
\draw[-{Stealth}, gray!40, thin] (11.2, 0.65) -- (11.2, 0.95);

% ── Objetivo de entrenamiento ─────────────────────────────────
\node[draw=orange!30, fill=orange!4, rounded corners, font=\tiny,
      inner sep=5pt, align=center] at (3.0, -2.0)
    {Entrenamiento:\\$\mathcal{L} = -\log P(c_{t+1} \mid c_1,\ldots,c_t)$\\
     Entrop\'ia cruzada sobre el\\siguiente car\'acter correcto};

% ── Flechas de flujo de la pérdida ────────────────────────────
\draw[-{Stealth}, orange!50, dashed, thick]
    (pred.west) -- +(-0.3, 0) -- +(-0.3, -0.7)
    node[below, font=\tiny, orange!70] {$\hat{y}$ vs $y$};

\end{tikzpicture}
\caption{Pipeline completo del Transformer m\'inimo implementado en este
trabajo. La secuencia de caracteres de entrada se convierte en \'indices
enteros, se proyecta al espacio de embeddings, se le suma la codificaci\'on
posicional, y se procesa a trav\'es de $N$ bloques encoder con m\'ascara
causal. La proyecci\'on lineal final y el softmax producen una distribuci\'on
de probabilidad sobre el vocabulario de caracteres. Durante el entrenamiento,
la p\'erdida de entrop\'ia cruzada compara la predicci\'on con el siguiente
car\'acter real de la secuencia.}
\label{fig:pipeline_minimo}
\end{figure}


%------------------------------------------------------------
\subsection{Función de pérdida y entrenamiento}
\label{subsec:entrenamiento}

El modelo se entrena minimizando la entropía cruzada entre la
distribución predicha $\hat{P}(c_{t+1})$ y el carácter real $c_{t+1}$:

\begin{equation}
    \mathcal{L} = -\frac{1}{T} \sum_{t=1}^{T}
    \log \hat{P}(c_{t+1} \mid c_1, \ldots, c_t)
\end{equation}

Esta pérdida se minimiza mediante descenso del gradiente con el
optimizador Adam~\cite{goodfellow2016deeplearning}, que extiende el
descenso del gradiente estocástico con estimaciones adaptativas de
momentos de primer y segundo orden, acelerando la convergencia en la
práctica.

Durante el entrenamiento, la máscara causal garantiza que el modelo no
acceda a información futura: el cómputo de $\hat{P}(c_{t+1})$ solo
utiliza $c_1, \ldots, c_t$. Esto hace posible calcular la pérdida para
todas las posiciones de la secuencia en paralelo, sin procesar los tokens
uno a uno como haría una RNN.

%------------------------------------------------------------
\subsection{Especificación formal de hiperparámetros}
\label{subsec:hiperparametros_minimo}

La Tabla~\ref{tab:hiperparametros_minimo} presenta la especificación completa
del modelo mínimo. Los valores están elegidos para que el entrenamiento
sea viable en hardware doméstico (CPU o GPU de consumo) en minutos, sin
sacrificar la capacidad de demostrar las propiedades fundamentales del
Transformer.

\begin{table}[H]
\centering
\caption{Especificación formal del Transformer mínimo implementado.}
\label{tab:hiperparametros_minimo}
\begin{tabularx}{\textwidth}{llXl}
\toprule
\textbf{Hiperparámetro} & \textbf{Notación} & \textbf{Descripción} & \textbf{Valor} \\
\midrule
Dimensión del modelo   & $d_{\text{model}}$ & Tamaño de embeddings y representaciones internas & 128 \\
Número de cabezas      & $h$               & Cabezas de atención paralelas & 4 \\
Dimensión por cabeza   & $d_k = d_v$       & $d_{\text{model}} / h$ & 32 \\
Número de bloques      & $N$               & Bloques encoder apilados & 2 \\
Dimensión FFN          & $d_{\text{ff}}$   & Capa interna de la red feed-forward & 512 \\
Longitud de contexto   & $T$               & Tokens procesados simultáneamente & 64 \\
Tasa de abandono       & $p_{\text{drop}}$ & Dropout para regularización & 0.1 \\
Tamaño de batch        & $B$               & Secuencias por paso de actualización & 32 \\
Tasa de aprendizaje    & $\eta$            & Paso del optimizador Adam & $3 \times 10^{-4}$ \\
\bottomrule
\end{tabularx}
\end{table}

Estos valores siguen el principio de que cada hiperparámetro debe ser
justificable. En particular:

\begin{itemize}
    \item $d_{\text{model}} = 128$ es suficiente para que el modelo
    aprenda representaciones de caracteres con estructura semántica,
    sin requerir memoria excesiva.

    \item $h = 4$ cabezas permiten capturar cuatro tipos de relaciones
    distintas en paralelo (por ejemplo, patrones fonéticos, morfológicos,
    sintácticos locales y dependencias de posición).

    \item $N = 2$ bloques producen un modelo con dos niveles de
    abstracción, lo que es suficiente para evidenciar el comportamiento
    jerárquico del Transformer.

    \item $T = 64$ garantiza un contexto suficiente para aprender
    patrones de longitud de palabra y estructura de oraciones simples.

    \item $d_{\text{ff}} = 512 = 4 \times d_{\text{model}}$ sigue la
    proporción estándar propuesta por Vaswani et al.~\cite{vaswani2017attention}.
\end{itemize}

%------------------------------------------------------------
\subsection{Comparación con modelos de producción}
\label{subsec:comparacion_escala}

Para contextualizar el modelo mínimo dentro del panorama actual, la
Tabla~\ref{tab:escala} lo compara con modelos de lenguaje representativos.

\begin{table}[H]
\centering
\caption{Comparación de escala entre el prototipo mínimo y modelos representativos.}
\label{tab:escala}
\begin{tabularx}{\textwidth}{lXXXX}
\toprule
\textbf{Modelo} & \textbf{$d_{\text{model}}$} & \textbf{$N$} & \textbf{$h$} & \textbf{Parámetros aprox.} \\
\midrule
Este trabajo (mínimo) & 128  & 2  & 4   & $\sim$0.41M \\
GPT-2 small           & 768  & 12 & 12  & $\sim$117M \\
GPT-2 large           & 1280 & 36 & 20  & $\sim$774M \\
GPT-3                 & 12288 & 96 & 96 & $\sim$175B \\
\bottomrule
\end{tabularx}
\end{table}

La diferencia de escala es enorme, pero los componentes son los mismos.
El prototipo pertenece a la misma familia decoder-only que GPT-2: atención
multi-cabeza con máscara causal, red feed-forward, conexiones residuales y
normalización de capa. Difiere en detalles de implementación (GPT-2 usa
posiciones aprendidas, activación GELU y normalización previa a cada
sub-capa), pero no en la idea. Esta es la propiedad más importante del
Transformer como arquitectura: su diseño escala de forma predecible y sin
cambios estructurales desde los $\sim$0.41M parámetros de este prototipo
(calculados en \texttt{verify/parametros.py} para un vocabulario de 65
caracteres) hasta los 175B de GPT-3~\cite{vaswani2017attention}.

En este sentido, el prototipo no es una simplificación del Transformer
real, sino el mismo modelo operando a escala didáctica.

\input{sections/07_conclusiones}

\newpage
\printbibliography[heading=bibintoc, title={Referencias}]

\newpage
\appendix
%\input{sections/A_implementacion}

\end{document}
